\documentclass[12pt,lettersize,journal,onecolumn]{IEEEtran}
\usepackage{setspace}
\doublespacing
\usepackage{soul}
\usepackage{xcolor}
\usepackage[table]{xcolor}
\sethlcolor{yellow}
\soulregister\cite{7}
\soulregister\ref{7}
\soulregister\eqref{7}
\newcommand{\sep}{\,$>$\,}
\DeclareRobustCommand{\myhl}[1]{{\sethlcolor{yellow}\hl{#1}}}
\usepackage{amsmath,amssymb,amsfonts}
\usepackage{mathrsfs} % for \mathscr
\newcommand{\scrT}{\mathscr{T}}
\newcommand{\setTau}{\scrT}
% \usepackage{algorithmic}
% \usepackage[ruled,vlined,linesnumbered]{algorithm2e}
\usepackage{algorithm}
\usepackage{algpseudocode}
% \SetKwInput{KwIn}{Input}           % \KwIn{...}
% \SetKwInput{KwOut}{Output}
% \SetKwInput{KwHyper}{Hyperparameters} % \KwHyper{...}
% \usepackage[margin=1in]{geometry}
\usepackage[caption=false,font=normalsize,labelfont=rm,textfont=rm]{subfig} % \usepackage[caption=false,font=normalsize,labelfont=sf,textfont=sf]{subfig}
\usepackage{textcomp}
\usepackage{stfloats}
%\usepackage{url}
\usepackage[colorlinks=true, allcolors=blue]{hyperref} % \usepackage{url} did not provide links but prevented difficulties with underscores in URLs
\usepackage{verbatim}
\usepackage{graphicx}
%\usepackage{amssymb} %for \checkmark and \bigstar
\usepackage{wasysym} % for \Circle, \CIRCLE, \LEFTcircle
\usepackage{pifont}
\usepackage{svg}
\usepackage{enumitem} % to use itemize with options (like negative space)
\usepackage{dsfont} % for \mathds{\xi}
\newcommand{\xmark}{\ding{55}}%%for \xmark
%\usepackage[table]{xcolor}
%\href{}{}\usepackage{xcolor} % to add color into figure captions
\usepackage{rotating}%to use for full-page landscape table (sidewaystable)
\usepackage{adjustbox} % for multi-cell in tables
\usepackage{multirow}% To work
\usepackage{booktabs}
\usepackage{colortbl}
\definecolor{ltgray}{gray}{0.93}
\definecolor{mdgray}{gray}{0.84}
\usepackage{tabularx}
\usepackage{bbm} % for \mathbbm{1} % binary condition
\usepackage{makecell} % for multiline
\usepackage{array} % for p{width} column spec
\usepackage[normalem]{ulem} % for multi-line underlining
\usepackage[font=small,skip=2pt]{caption} %\usepackage{booktabs}%\usepackage{float} % skip option is for the vertical space between the caption and the figure/table
% \usepackage{subcaption}
\usepackage{cite}
\usepackage{siunitx}
\sisetup{
  table-number-alignment = center,
  table-text-alignment   = center
}
\newcolumntype{Y}{>{\raggedright\arraybackslash}X}
% \newcommand{\rothead}[1]{\rotatebox[origin=c]{90}{\textbf{#1}}}
\newcommand {\squeezeup} {\vspace {-3.8mm}}% %\squeezeup will reduce the free space before and after the figures
\newcommand {\squeezeupsection} {\vspace {-2.4mm}}

\makeatletter %defining subsubsubsection
\newcommand{\subsubsubsection}[1]{\paragraph{#1}}
\makeatother
%=============================================
% 1) COUNTER & ENVIRONMENT FOR "R" CONSTRAINTS
%=============================================
\newcounter{constraintR}
\renewcommand{\theconstraintR}{C\arabic{constraintR}-R}

\newenvironment{NamedConstraintR}{%
  \refstepcounter{constraintR}%  increments the R counter
  \par\noindent
  \textbf{\theconstraintR:}\,    % prints e.g. "C1-R:"
  \begin{math}
}{%
  \end{math}
  \par\vspace{2pt}%
}
%=============================================
% 2) COUNTER & ENVIRONMENT FOR "T" CONSTRAINTS
%=============================================
\newcounter{constraintT}
\renewcommand{\theconstraintT}{C\arabic{constraintT}-T}

\newenvironment{NamedConstraintT}{%
  \refstepcounter{constraintT}%  increments the T counter
  \par\noindent
  \textbf{\theconstraintT:}\,    % prints e.g. "C1-T:"
  \begin{math}
}{%
  \end{math}
  \par\vspace{2pt}%
}
%=============================================
% 3) COUNTER & ENVIRONMENT FOR "C" CONSTRAINTS
%=============================================
\newcounter{constraintC}
\renewcommand{\theconstraintC}{C\arabic{constraintC}-C}

\newenvironment{NamedConstraintC}{%
  \refstepcounter{constraintC}%  increments the C counter
  \par\noindent
  \textbf{\theconstraintC:}\,    % prints e.g. "C1-C:"
  \begin{math}
}{%
  \end{math}
  \par\vspace{2pt}%
}
%=============================================
% 1) COUNTER & ENVIRONMENT FOR "E" CONSTRAINTS
%=============================================
\newcounter{constraintE}
\renewcommand{\theconstraintE}{C\arabic{constraintE}-E}

\newenvironment{NamedConstraintE}{%
  \refstepcounter{constraintE}%  increments the E counter
  \par\noindent
  \textbf{\theconstraintE:}\,    % prints e.g. "C1-E:"
  \begin{math}
}{%
  \end{math}
  \par\vspace{2pt}%
}

%===============
\usepackage{environ}
% Tight flalign (keeps numbering)
\NewEnviron{tightflalign}{%
  \begingroup
  \setlength{\abovedisplayskip}{1pt}%0.17
  \setlength{\abovedisplayshortskip}{1pt}%0.17
  \setlength{\belowdisplayskip}{0pt}%0.05
  \setlength{\belowdisplayshortskip}{0pt}%0.05
  \begin{flalign}
    \BODY
  \end{flalign}
  \endgroup
}
%===============

%%%%
\hyphenation{op-tical net-works semi-conduc-tor IEEE-Xplore}

%%%%%%%%%%%%%%%%%%%%%%%
\begin{document}
\title{ENSURE: Multi-Timescale E2E Slice Resource Management for 6G Networks} 
% \title{Multi-Agent DRL-Based Multi-Timescale E2E Slice Resource Management for Open RAN-Compatible 6G Networks} 
% \title{A Novel Multi-Agent DRL-based Multi-Timescale E2E Slice Resource Management Scheme for Open RAN-Compatible 6G Networks} 

% Soufiene: Multi-Agent DRL-based Multi-Timescale E2E Slice Resource Management in Open RAN-Compatible 6G Networks
% Multi-Timescale E2E Slice Resource Management in Open RAN-Compatible 6G Networks Using Multi-Agent DRL% End-to-End Network Slice

\author{Sina Ebrahimi,~\IEEEmembership{Student Member,~IEEE},
        Hui Zhou, %,~\IEEEmembership{Member,~IEEE},
        Faouzi Bouali, %~\IEEEmembership{Member,~IEEE},
        Soufiene Djahel,~\IEEEmembership{Senior Member,~IEEE},
        Olivier C. L. Haas %,~\IEEEmembership{Senior Member,~IEEE}
\thanks{Manuscript received April 21, 2026; \\S. Ebrahimi, H. Zhou, F. Bouali, S. Djahel, and O. Haas are with the Centre for Future Transport and Cities (CFTC), Coventry University. \textit{ENSURE} is used here as an academic acronym and is not affiliated with any third-party marks. \textit{Corresponding author: Soufiene Djahel; e-mail: ae3095@coventry.ac.uk.}}} %(e-mail: \{ebrahimis, ae4818, ae3095, csx259\}@coventry.ac.uk).} %Faouzi.Bouali@wlv.ac.uk %Coventry, CV1 2LQ, UK. %Faouzi Bouali is with the School of Coding and AI, in partnership with the University of Wolverhampton.

\markboth{IEEE Transactions on Cognitive Communications and Networking,~Vol.~x, No.~x, Second Quarter~2026}{Ebrahimi \MakeLowercase{\textit{et al.}}: ENSURE: Multi-Timescale E2E Slice Resource Management for 6G Networks} %MA-DRL-Based Multi-Timescale E2E Slice Resource Management for Open RAN-Compatible 6G Networks}
%\markboth{IEEE Transactions on Network and Service Management,~Vol.~x, No.~x, Second Quarter~2026}
%\IEEEpubid{0000--0000/00\$00.00~\copyright~2023 IEEE}

\maketitle
\begin{abstract}
Efficient End-to-End (E2E) slice resource management is essential for 5G/6G, yet coordinating resources across the Radio Access Network (RAN), Transport Network (TN), and Core Network (CN) remains challenging due to cross-domain coupling and heterogeneous control timescales. This paper proposes \textit{ENSURE}, a hierarchical Multi-Agent Deep Reinforcement Learning (MA-DRL) framework for integrated E2E Network Slicing (NS). ENSURE comprises four cooperative agents at three timescales: a Soft Actor-Critic (SAC) RAN agent performs near-real-time user association and Physical Resource Block (PRB) allocation; Proximal Policy Optimization (PPO) TN and CN agents reconfigure paths and Virtual Network Function (VNF) placements at mid-level intervals, respectively; and an E2E orchestrator coordinates cross-domain Slice Instance (SI) assignments at long intervals. Satisfaction-triggered action masking suppresses unnecessary reconfigurations, while Integer Linear Program (ILP) warm-starting seeds each episode with feasible assignments. Under curriculum-based training and severe multi-domain stress, ENSURE achieves an E2E satisfaction of $0.73$, outperforming ablated variants and a baseline with adaptive RAN but non-adaptive TN/CN (representative of existing RAN slicing approaches) by at least $38\%$, while keeping SI E2E delays within their Service Level Agreement (SLA) budgets.
%Multi-Agent Deep Reinforcement Learning (MA-DRL)
%Physical Resource Block (PRB)
%Virtual Network Function (VNF)
%Slice Instance (SI)
%PPO
%SAC
%This paper proposes \textit{ENSURE}, a hierarchical multi-agent deep reinforcement learning framework for integrated E2E network slicing in Open RAN-compatible architectures
%Simulations under curriculum-based training and multi-domain stress scenarios, including concurrent TN link failures, CN server outages, and background congestion, show that ENSURE achieves 0.73 E2E satisfaction, outperforming ablated variants without masking (0.24) or warm-start (0.22), and a static RAN-only baseline (0.50), while maintaining eMBB and URLLC delays within their respective SLA budgets.

% Efficient End-to-End (E2E) slice resource management is essential for 5G/6G, yet coordinating resources across the Radio Access Network (RAN), Transport Network (TN), and Core Network (CN) remains difficult. This paper proposes \textit{ENSURE}, a hierarchical multi-agent deep reinforcement learning framework for integrated E2E slicing that captures cross-domain coupling while preserving stable operation. Aligned with Open RAN architecture, the RAN agent performs mobility-aware user association and physical resource block allocation at near-real-time intervals, while TN and CN agents reconfigure paths and virtual network function placements at slower timescales. An E2E agent coordinates the domain agents via \textit{slice instance} abstractions, selecting CN sites and aggregating cross-domain satisfaction, while satisfaction-triggered action masking suppresses unnecessary reconfigurations. Episodes are warm-started with feasible initial assignments to accelerate learning. Simulations under heterogeneous traffic demonstrate improved user satisfaction, resource efficiency, and resilience compared to single-domain baselines.

%Efficient resource management is vital for diverse 6G applications, yet End-to-End (E2E) slicing remains challenging because bottlenecks shift across the RAN, TN, and CN and the control loops operate on heterogeneous timescales. Existing solutions are often domain-specific or assume static cross-domain provisioning, which breaks under transport/core congestion or failures and can induce disruptive reconfigurations.

% Efficient resource management is vital for supporting diverse applications and services in emerging 6G networks. This paper proposes a novel multi-agent Deep Reinforcement Learning (DRL)-based multi-timescale framework for integrated resource management across the Radio Access Network (RAN), Transport Network (TN), and Core Network (CN), enabling mobility-aware and dynamic End-to-End (E2E) Network Slicing (NS) in 6G networks. The proposed E2E slice resource management (\textit{ENSURE}) framework incorporates the emerging Open RAN architecture to ensure flexibility and interoperability. The RAN agent optimizes user assignment and Physical Resource Block (PRB). Conversely, the CN and TN agents reconfigure Virtual Network Function (VNF) placement and active paths in slower timescales, respectively. An E2E agent supervises these three agents, optimizing a global reward while providing feedback to the local agents on a slower timescale. Extensive simulations demonstrate the efficacy of the proposed methodology in improving overall user satisfaction, resource utilization, and adaptability under dynamic conditions, compared with single-domain NS frameworks.

% Efficient resource management is vital for supporting diverse applications and services in emerging 6G networks. This paper presents a novel approach for joint resource management in an End-to-End (E2E) network encompassing Radio Access Network (RAN), Transport Network (TN), and Core Network (CN) segments, especially focusing on Network Slicing (NS). The proposed E2E framework also accounts for the emerging Open RAN architecture. The proposed solution leverages multi-agent Deep Reinforcement Learning (DRL) to handle the complexities and high uncertainties associated with mobility-aware environments. The RAN model, which optimizes user assignment and PRB power allocation, is designed to be compatible with Open RAN architecture, ensuring flexibility and interoperability. The Virtual Network Function (VNF) placement is handled at the core network, while bandwidth allocation is managed at the transport network. The local agent in each domain has its own MDP representation and maximizes a local reward, while a critic network in the E2E orchestrator provides a global state to these agents and maximizes a global reward. Extensive simulations demonstrate the efficacy of the proposed approach in enhancing network performance, resource utilization, and adaptability under dynamic conditions. 

% In order to address the ever-changing nature of the mobile environment and expedite convergence, we utilize meta-learning to effectively handle the fluctuations in system dynamics within our MA-DRL approach.
\end{abstract}

\begin{IEEEkeywords}
5G, 6G, Resource Allocation (RA), Network Slicing (NS), Multi-Agent Deep Reinforcement Learning (MA-DRL), End-to-End (E2E).
\end{IEEEkeywords} % could be removed : Open RAN, 

\section{Introduction} \label{introduction-section}
% \subsubsection{State-of-the-Art and Motivation} \label{subsec-introduction-motivation}
%6G context and motivation
\IEEEPARstart{T}{he} sixth generation (6G) of mobile networks is expected to face unprecedented heterogeneity in service requirements, building on the fifth generation (5G) service classes such as enhanced Mobile Broadband (eMBB) and Ultra-Reliable Low-Latency Communication (URLLC) that already capture opposing rate--latency extremes~\cite{itu23m2160}. While 6G service categories are still evolving, they are widely anticipated to be even more stringent and mutually conflicting, making a \textit{one-size-fits-all} network design impractical. \textit{Network Slicing (NS)} addresses this by creating multiple logical networks tailored to specific services over shared infrastructure, enabling the coexistence of rate-intensive (e.g., virtual reality) and latency-critical (e.g., autonomous vehicles) applications with slice-specific QoS guarantees~\cite{comst24ebrahimi}. Dynamic slice-aware Resource Allocation (RA) is therefore crucial to prevent over- and under-provisioning while meeting diverse End-to-End (E2E) demands. 
%Removed Quality of Service (QoS) due to IEEE Style Manual (assumed known acronyms); see Appnedix A: Some Common Acronyms
%\IEEEPARstart{T}{he} sixth generation (6G) of mobile networks must support unprecedented heterogeneity in service requirements, from tactile remote operations to pervasive smart-city infrastructures~\cite{itu20fgnet2030}. \textit{Network Slicing (NS)} addresses this by creating multiple logical networks tailored to specific services over shared infrastructure, enabling the coexistence of rate-intensive (e.g., enhanced Mobile Broadband (eMBB)) and latency-critical (Ultra-Reliable Low-Latency Communication (URLLC)) applications with slice-specific Quality of Service (QoS) guarantees~\cite{comst24ebrahimi}. Dynamic slice-aware Resource Allocation (RA) is therefore crucial to prevent overprovisioning and starvation while meeting diverse End-to-End (E2E) demands.
%merged these two
% \IEEEPARstart{T}{he} sixth generation (6G) of mobile networks is expected to support unprecedented connectivity and service capabilities by 2030. To enable emerging applications such as tactile remote operations and pervasive smart-city infrastructures, 6G must substantially improve capacity, latency, energy efficiency, and native intelligence over 5G \cite{itu20fgnet2030}. While 5G introduced service classes such as enhanced Mobile Broadband (eMBB), Ultra-Reliable Low-Latency Communication (URLLC), and massive Machine-Type Communication (mMTC), 6G is expected to extend these classes and impose even more stringent and heterogeneous service requirements, thereby requiring a fundamental redesign of network architecture and End-to-End (E2E) resource management to meet diverse Quality of Service (QoS) demands simultaneously \cite{comst24ebrahimi}.
% %new paragraph: Network slicing as the mechanism
% One key mechanism to handle this heterogeneity is \textit{Network Slicing (NS)}, which creates multiple logical networks (slices) tailored to specific services or tenants over a shared physical infrastructure \cite{comst24ebrahimi}. By logically isolating resources per slice, NS enables the coexistence of rate-intensive services (e.g., virtual reality) and latency-critical ones (e.g., industrial control) with slice-specific QoS guarantees. Dynamic slice-aware Resource Allocation (RA) is therefore crucial to prevent both overprovisioning and starvation and to fully exploit the flexibility and efficiency that NS offers for 6G systems \cite{comst24ebrahimi}. % By allocating (logically) isolated resources per slice, NS enables the coexistence of rate-intensive services (e.g., augmented/virtual reality) and latency-critical ones (e.g., industrial control) with slice-specific QoS guarantees.
%Open RAN and programmability 
Realizing NS requires programmable control, particularly in the Radio Access Network (RAN). The \textit{Open RAN} architecture~\cite{oran23architecture} enables this via disaggregated components and standardized interfaces, with RAN Intelligent Controllers (RICs) hosting AI/ML-driven applications for closed-loop slice management~\cite{comst23polese}.
%Realizing NS in practice, however, requires programmable and interoperable control, particularly in the Radio Access Network (RAN). The \textit{Open RAN} architecture promotes open, disaggregated RAN components with standardized interfaces, enabling multi-vendor interoperability and fine-grained resource control \cite{oran23architecture}. At its core, the RAN Intelligent Controller (RIC), with non-RT and near-RT instances, hosts AI/ML-driven control applications (rApps/xApps) that implement closed-loop, slice-aware resource management \cite{comst23polese}. This programmability makes Open RAN a natural architectural substrate for automated and adaptive NS in 6G systems \cite{oran23slicing}.

%Need for E2E, multi-domain coordination
Slice performance is constrained by the bottleneck along the E2E path spanning RAN, Transport Network (TN), and Core Network (CN) segments~\cite{comst24ebrahimi}. Meeting slice-level Service Level Agreements (SLAs) thus requires coordinated multi-domain optimization beyond isolated single-segment designs~\cite{jiot24filali,mnet23liu}. %isolated single-segment designs~\cite{jiot24filali,tnsm23johari,noms20ebrahimi}. 
%Critically, slice performance is ultimately constrained by the weakest segment along the E2E path, spanning the RAN, Transport Network (TN), and Core Network (CN) \cite{comst24ebrahimi}. Hence, meeting slice-level Service Level Agreements (SLAs) requires coordinated slicing and resource optimization across \emph{all} domains. While many works address NS within a single segment, for example dynamic RAN slicing \cite{jiot24filali}, TN slice embedding \cite{tnsm23johari}, or CN admission and resource control \cite{noms20ebrahimi}, such isolated designs ignore cross-domain dependencies and cannot systematically guarantee E2E QoS. This has motivated the development of integrated multi-domain resource management strategies for E2E slicing \cite{mnet23liu,tnsm23gharehgoli}. %This has motivated integrated multi-domain resource management strategies for E2E slicing \cite{mnet23liu,tnsm23gharehgoli}.
%Multi-timescale nature + SI abstraction + Hierarchical control and E2E agent
A further challenge is timescale heterogeneity: RAN control reacts at the millisecond scale, whereas TN routing and CN Virtual Network Functions (VNF) reconfiguration operate at slower sub-second scales. A single control loop is neither feasible nor desirable; instead, a hierarchical framework with per-domain agents and a top-level E2E coordinator is required~\cite{mnet23liu}. %instead, a hierarchical framework is needed where each domain (RAN, TN, CN) is governed by a local agent and a top-level E2E agent coordinates their actions~\cite{etsi21gszsm003,mnet23liu}. %To keep this tractable, we introduce a \textit{Slice Instance (SI)} abstraction that aggregates per-slice traffic at the TN/CN level, dramatically reducing decision-space complexity while preserving QoS guarantees. %TN routing and CN Virtual Network Functions (VNF) reconfiguration operate at slower (sub-)second scales

%Multi-timescale nature + SI abstraction
%A further challenge is the heterogeneity of control timescales across domains. RAN control must react to fast channel dynamics at the millisecond scale, whereas TN routing adaptation, as well as CN Virtual Network Functions (VNF) reconfiguration, operate at slower (sub-)second scales. A single unified control loop is therefore neither feasible nor desirable; instead, fine-grained, low-latency decisions in the RAN must be coordinated with coarser adaptations in the TN and CN. To keep this coordination tractable, we introduce a \textit{Slice Instance (SI)} abstraction that aggregates, for each slice type and Central Unit (CU), the traffic of its associated users at the TN/CN level. Operating on SIs rather than on per-user flows dramatically reduces the decision space while preserving slice-specific QoS guarantees.
%Hierarchical control and E2E agent
%These requirements naturally call for a hierarchical control framework. In line with emerging standards \cite{etsi21gszsm003,mnet23liu}, we assume that each domain (RAN, TN, CN) is governed by a local agent, whereas a top-level E2E agent maintains a global view and coordinates their actions. By steering domain agents through cross-domain constraints and high-level objectives (e.g., jointly prioritizing latency-sensitive slices in RAN scheduling, TN path selection, and CN processing), the E2E agent prevents myopic or conflicting decisions and aligns local actions with SLAs. Such hierarchical coordination is key to robust E2E slice management under dynamic conditions. % local controller/agent % aligns local actions with slice-level SLAs % robust, adaptive E2E slice management

%Why MA-DRL
Given the coupled, high-dimensional nature of E2E slicing across multiple timescales~\cite{comst24ebrahimi}, a hierarchical \textit{Multi-Agent Deep Reinforcement Learning (MA-DRL)} design provides a scalable decomposition~\cite{marl24albrecht}, with recent work demonstrating its effectiveness under dynamic NS conditions~\cite{tnsm24lu,mnet23liu}: each domain agent learns an independent policy on its own timescale, while a higher-level E2E agent acts via its own policy and guides them via compact coordination signals summarizing cross-domain slice status. This avoids monolithic joint control while retaining cross-domain awareness and aligns with Open RAN’s AI/ML-driven closed-loop operation~\cite{oran23architecture}. However, naive MA-DRL on combinatorial problems suffers from sample inefficiency and policy-update oscillation across iterations; action masking and warm-start initialization are therefore needed to enforce feasibility and accelerate convergence~\cite{marl24albrecht}.

%Given the coupled, high-dimensional nature of E2E slicing, spanning multiple decisions in different timescales~\cite{comst24ebrahimi},  % Given the dimensionality and distributed nature of this control problem, \textit{Multi-Agent Deep Reinforcement Learning (MA-DRL)} provides a natural solution paradigm\cite{marl24albrecht}. DRL has already shown strong performance in complex resource management tasks by learning policies from high-dimensional observations without explicit models. Extending DRL to a multi-agent setting enables us to train domain-specific agents (RAN, TN, CN) alongside an E2E agent that coordinates them via shared information and reward signals. Recent works report the effectiveness of (MA-)DRL for NS and wireless resource management under dynamic traffic and mobility \cite{tnsm24lu,mnet23liu,camad24ebrahimi}. Moreover, DRL-based closed-loop control is fully aligned with the AI/ML-centric vision of Open RAN.

%High-level pitch of framework
However, scalable E2E solutions that jointly address cross-domain coupling, heterogeneous timescales, and consistency remain limited. The following subsection reviews the state of the art and positions our contribution in this landscape.
%However, scalable E2E solutions for Open RAN slicing remain limited. This paper proposes \textit{\underline{EN}d-to-End \underline{S}lice Reso\underline{UR}ce Manag\underline{E}ment} (\textit{ENSURE}), a hierarchical, multi-timescale MA-DRL framework for slice resource management in Open RAN-enabled 5G/6G networks. It extends our prior RAN-only xApp study~\cite{camad24ebrahimi} to an E2E setting by jointly coordinating RAN, TN, and CN under a supervisory E2E agent. Subsec.~\ref{related-work-section} situates this research in the context of existing state-of-the-art literature.

\subsection{Related Work} \label{related-work-section}
%Domain-Specific Slicing and Open RAN Integration
Prior work on NS resource management spans single-domain studies, multi-domain extensions, and cooperative MA-DRL designs. Single-domain studies tackle each segment in isolation~\cite{comst24ebrahimi}: DRL-assisted slice embedding in the TN~\cite{tnsm23johari}, admission control in the CN~\cite{noms20ebrahimi}, and PRB/power optimization in the RAN~\cite{tvt23choi}. RAN-centric work also exploits Open RAN for hierarchical DRL-based slicing~\cite{jiot24filali}. These provide strong domain baselines but cannot deliver coordinated E2E control.
%RAN~\cite{tnsm23karbalaee,tvt23choi,icspis18dehkordi,camad24ebrahimi}

%Multi-domain E2E (single-/centralised)
Cross- and multi-domain studies are either partial or centralized. Two-timescale RAN+CN schemes~\cite{jstsp21reyhanian} and three-domain optimization with relaxed integer variables~\cite{tnse25nouruzi} capture cross-domain coupling but use a single learning entity, limiting scalability. Hierarchical NS architectures are envisioned in~\cite{mnet23liu} without a concrete multi-agent instantiation at domain timescales, while E2E admission control in~\cite{tnsm22luu} uses heuristic reconfiguration without adaptive learning. % domain-appropriate timescales
%Multi-Objective Slicing and Slice Satisfaction Metrics
% NS inherently induces conflicting objectives (rate vs.\ delay, utilization vs.\ fairness). Early works scalarize multiple metrics into weighted sums~\cite{tnsm23gharehgoli,noms20ebrahimi}. Accordingly, satisfaction metrics range from geometric-mean combinations of rate and delay~\cite{camad24ebrahimi} to threshold-based fractions~\cite{tnsm22bai}. Such measures are effective within their targeted settings, but differences in scale and semantics complicate consistent aggregation and control across coupled RAN/TN/CN objectives. Moreover, treating metrics as objectives versus constraints remains a design choice; embedding requirements as soft constraints in the reward enables adaptation but requires careful shaping~\cite{marl24albrecht}. 
%\cite{mwc24esmat} %Moreover, treating metrics as objectives versus constraints remains a design choice with no universal solution; embedding requirements as soft constraints in the reward enables adaptive recovery but requires careful shaping to avoid reward instability. %Such measures are effective within their targeted settings, but differences in scale, semantics, and entity granularity complicate consistent aggregation and control across coupled RAN/TN/CN objectives. %packet reception probabilities~\cite{twc22setayesh} . % Early works scalarize metrics into weighted sums~\cite{access20abiko,tnsm23gharehgoli}, while later studies pursue Pareto fronts via MA-DRL~\cite{tvt23zhou}

%MA-DRL for NS
Cooperative MA-DRL has been explored mostly within a single domain or across two domains. Within the RAN, the federated constrained MA-DRL framework in~\cite{dyspan26mohammadi} uses Lagrangian Proximal Policy Optimization (PPO) with federated averaging for inter-cell spectrum slicing in dense 6G, and~\cite{jsac24zangooei,jiot24filali} coordinate inter- and intra-slice PRB allocation in a two-level RAN hierarchy. At the RAN--CN boundary, \cite{access21kim} splits radio and computing decisions across cooperating agents, while~\cite{tnsm22sulaiman} uses cooperative MA-DRL for joint slice admission and RAN resource control. None of these studies extends cooperative MA-DRL to the full RAN--TN--CN chain, and, except~\cite{jsac24zangooei,jiot24filali}, each operates on a single control timescale.
%\cite{srep25mao} couples MA-DRL with evolutionary game theory for edge-enabled networks, mitigating tenant competition but modeling only a single RAN--edge boundary.

%Multi-timescale + cross-domain abstractions
RA decisions span multiple timescales~\cite{comst24ebrahimi}: two-timescale RAN schemes~\cite{jiot24filali,jsac24zangooei} and four-timescale RAN models~\cite{twc25ataeebojd} partly capture this separation but remain RAN-centric. Vehicular two-timescale MA-DRL for spectrum/computing decisions~\cite{tnsm24lu} demonstrates the viability of temporal decomposition but does not cover the full E2E path. Integrating cross-domain Slice Instance (SI)-level abstractions with cooperative MA-DRL across RAN, TN, and CN under three heterogeneous timescales remains the gap addressed by \textit{ENSURE}.
%DRL and Multi-Agent Approaches for NS
%DRL-based allocation improves QoS adherence under traffic and channel variations~\cite{tnsm23gharehgoli}, but single-agent designs face scalability limits. MA-DRL mitigates this via cooperative agents~\cite{jsac24zangooei}, including hierarchical designs that separate long- from short-term control~\cite{mnet23liu,jiot24filali}. Coordination may rely on shared rewards, message passing, or periodic parameter sharing. Recent E2E frameworks either formulate joint RAN/TN/CN optimization as a centralized combinatorial problem~\cite{tnse25nouruzi} or adopt hierarchical controllers that trade optimality for scalability under domain constraints~\cite{mnet23liu}. % two-timescale: \cite{jstsp21reyhanian} %Action masking enforces feasibility by removing invalid actions from the policy distribution prior to sampling. %MA-DRL mitigates this via cooperative agents~\cite{jsac24zangooei,tnsm22sulaiman }; tnsm22sulaiman (like us) used separate RL agents % Action masking enforces feasibility by filtering invalid actions prior to sampling~\cite{flairs22huang,tmc25nahum} 
%adopt modular/hierarchical controllers that trade optimality for scalability and deployability under domain constraints~\cite{mnet23liu

%Multi-Timescale Control and Cross-Domain Abstractions
% Finally, RA decisions span multiple timescales~\cite{comst24ebrahimi}: two-timescale RAN schemes~\cite{jiot24filali} and more elaborate four-timescale models~\cite{twc25ataeebojd} partly capture this separation but remain RAN-centric. Cross-domain abstractions map heterogeneous domain states and actions into compact slice-level summaries (e.g., per-slice PRB budgets, per-path congestion indicators), reducing signaling and control dimensionality. However, integrating such abstractions with cooperative MA-DRL across RAN, TN, and CN under heterogeneous timescales remains limited, which is the main gap addressed by this work.

%%%%%%%%%%%%%%%%%%%%%%%%%
\subsection{Contributions} \label{contribution-section}
The works above tackle either a single domain, a two-domain pair, or E2E slicing with a single learning entity. To the best of our knowledge, no prior work couples cooperative MA-DRL across all three domains (RAN, TN, CN) with domain-appropriate timescales and an SI-level abstraction that lets TN/CN/E2E control scale with SIs rather than users. \textit{ENSURE}, the proposed \textit{\underline{EN}d-to-End \underline{S}lice Reso\underline{UR}ce Manag\underline{E}ment} framework, addresses this gap and extends our RAN-only study~\cite{camad24ebrahimi} to E2E control via a supervisory E2E agent. The main contributions are:
%; RAN-specific modeling follows~\cite{camad24ebrahimi}
% To address the above limitations, we propose \textit{\underline{EN}d-to-End \underline{S}lice Reso\underline{UR}ce Manag\underline{E}ment} (\textit{ENSURE}), a hierarchical, multi-timescale MA-DRL framework for slice resource management in Open RAN-enabled 5G/6G networks. Building on our RAN-only study in~\cite{camad24ebrahimi}, \textit{ENSURE} extends the formulation to E2E control by jointly coordinating RAN, TN, and CN via a supervisory E2E agent to satisfy slice-specific QoS while preserving network consistency; RAN-specific modeling and solution details follow~\cite{camad24ebrahimi}. The main contributions include: % stability or consistency? (Olivier asked!)

\begin{itemize}[leftmargin=1.1em,labelsep=0.3em]
\item \textbf{Cooperative MA-DRL for three-timescale E2E slicing:} ENSURE coordinates four agents (RAN/TN/CN/E2E) at three timescales for full E2E slice management. Prior MA-DRL designs operate within a single~\cite{dyspan26mohammadi,jsac24zangooei} or across two domains~\cite{access21kim,tnsm22sulaiman}; existing E2E designs either centralize decisions~\cite{tnse25nouruzi} or lack explicit timescale separation~\cite{mnet23liu}.
% \item \textbf{Hierarchical multi-timescale E2E management:} \textit{ENSURE} coordinates four cooperative agents (RAN, TN, CN, and E2E) operating at three distinct timescales. Unlike prior E2E frameworks that either centralize decisions~\cite{tnse25nouruzi} or lack explicit timescale separation~\cite{mnet23liu}, our design couples near-real-time RAN control with slower TN/CN adaptation through lightweight coordination signals. %lightweight SI-level coordination signals
% \item \textbf{Hierarchical multi-timescale E2E management:} \textit{ENSURE} coordinates four cooperative agents (RAN, TN, CN, and E2E) operating at different timescales to capture cross-domain dependencies beyond single-domain baselines.

\item \textbf{Slice Instance abstraction:} We introduce a Slice Instance (SI) as the coordination unit that bundles users of the same slice type sharing a CN site and TN path, enabling SI-level TN/CN control instead of per-user decisions and thereby reducing action-space complexity. 

\item \textbf{Unified satisfaction framework:} A sigmoid-based mapping normalizes heterogeneous QoS indicators into comparable satisfaction scores. Aggregation via geometric means promotes balanced optimization across metrics, mitigating the dominance of single metrics at the expense of others.
%Aggregation via geometric means promotes balanced optimization across metrics, preventing policies that sacrifice one objective for marginal gains in others.
% \item \textbf{Unified satisfaction and reward aggregation:} A sigmoid-based mapping normalizes heterogeneous QoS indicators into comparable satisfaction scores, enabling consistent cross-domain reward aggregation and multi-objective trade-offs via geometric-mean aggregation. % explicit multi-objective trade-offs (instead of controllable trade-offs) & multi-metric (instead of multi-objective)

\item \textbf{Warm-start and satisfaction-triggered masking:} Episodes warm-start with Integer Linear Program (ILP)-based feasible assignments. Unlike standard masked PPO~\cite{flairs22huang} that lifts masks only on hard infeasibility, our mask lifts on a soft satisfaction-pivot threshold, triggering TN/CN/E2E reconfiguration before SLA violation and suppressing feasible-but-unnecessary changes for consistency.
%\item \textbf{Warm-start and event-triggered reconfiguration:} Episodes warm-start with Integer Linear Program (ILP)-based feasible assignments. \textit{Satisfaction-triggered action masking} enables TN/CN/E2E reconfigurations only when degraded performance is detected, reducing unnecessary churn while focusing adaptation on under-performing SIs.
% \item \textbf{Warm-start and stability-aware reconfiguration:} Each episode warm-starts with Integer Linear Program (ILP)-based feasible assignments, and satisfaction-triggered action masking blocks TN/CN/E2E reconfigurations when satisfaction is already adequate, reducing churn and disruption while focusing exploration on under-performing SIs. % potential service disruption

% \item \textbf{SI abstraction with unified satisfaction:} An SI abstraction aligns fast RAN control with slower TN/CN adaptations while reducing their decision complexity, and a sigmoid-based satisfaction mapping normalizes QoS metrics to enable consistent cross-domain reward aggregation.
%To keep this tractable, we introduce a \textit{Slice Instance (SI)} abstraction that aggregates per-slice traffic at the TN/CN level, dramatically reducing decision-space complexity while preserving QoS guarantees. % a sigmoid-based satisfaction mapping normalizes QoS metrics to enable coherent cross-domain reward aggregation. 

% \item \textbf{Satisfaction-triggered masking:} Episodes begin with Integer Linear Program (ILP) solutions; action masking suppresses reconfigurations unless satisfaction drops below a pivot point, supporting stability and feasibility.

% \item \textbf{Performance gains:} Simulations show substantial improvements in satisfaction, utilization, and resilience over single-domain baselines.
% \item \textbf{Multi-timescale optimization with SI abstraction:} A \textit{SI} abstraction aggregates per-slice demand at the TN/CN level, coordinating fast RAN decisions with slower TN/CN adaptations while reducing decision-space complexity.
% \item \textbf{Unified satisfaction framework:} A sigmoid-based mapping normalizes heterogeneous QoS metrics into $[0,1]$ scores, enabling consistent cross-domain reward aggregation via weighted geometric means.
%\item \textbf{Multi-agent DRL framework:} We design a hierarchical MA-DRL framework comprising domain-level DRL agents (for RAN, TN, and CN) and a higher-level E2E agent. Domain agents learn specialized policies for their respective segments, while the E2E agent oversees and guides them toward global objectives. To the best of our knowledge, this is one of the first works to apply a cooperative MA-DRL approach to an integrated E2E slicing problem with multiple timescales.
%\item \textbf{Open RAN-compatible implementation:} Our solution is explicitly developed for Open RAN deployments. In particular, the RAN agent is mapped to a near-RT RIC control application (xApp) for dynamic RAN slicing \cite{camad24ebrahimi}, while standard controllers in the TN and CN implement the corresponding domain-agent decisions. This alignment with Open RAN demonstrates the implementability of the proposed framework in realistic 6G settings.
\end{itemize}
%%%%%%%%%%%%
\subsection{Paper Organization}
The remainder of the paper is structured as follows. Sec.~\ref{section-architectural-framework} outlines ENSURE's architectural framework, and Sec.~\ref{system-model-section} details the system model. Sec.~\ref{solution-section} casts the RAN, TN, CN, and E2E problems as Partially Observable Markov Decision Processes (POMDPs) and describes the MA-DRL coordination scheme. Sec.~\ref{sec:results} presents the simulation setup and performance results, and Sec.~\ref{conclusion-section} concludes the paper.
%The remainder of the paper is structured as follows. Sec.~\ref{related-work-section} reviews related work. Sec.~\ref{section-architectural-framework} outlines ENSURE's architectural framework, and Sec.~\ref{system-model-section} details the system model and problem formulation. Sec.~\ref{solution-section} casts the RAN, TN, CN, and E2E problems as Partially Observable Markov Decision Processes (POMDPs) and describes the MA-DRL coordination scheme. Sec.~\ref{performance-evaluation} presents the simulation setup and performance results, and Sec.~\ref{conclusion-section} concludes the paper.  % The computational complexity of the proposed method is assessed in Sec.~\ref{complexity-section}, while 

%%%
\section{ENSURE Architecture} \label{section-architectural-framework} %Proposed Multi-Timescale E2E Slice Resource Management Framework Architecture
% This section presents the \textit{ENSURE} architecture, a hierarchical multi-agent framework for joint RAN, TN, and CN slice management. The architecture embeds explicit multi-timescale control and aligns with ongoing standardization~\cite{3gpp21tr28811,oran23architecture,tmf25agent}. 
Open RAN~\cite{oran23architecture} disaggregates the RAN into the Radio Unit (RU), Distributed Unit (DU), and Central Unit (CU), interconnected through standard interfaces, and introduces two RAN Intelligent Controllers (RICs): a near-Real-Time (near-RT) RIC hosting xApps for control decisions on a $10$--$1000$~ms timescale, and a non-RT RIC within the Service Management and Orchestration (SMO) plane hosting rApps for $\geq1$~s decisions~\cite{comst23polese}. ENSURE aligns naturally with this control architecture: the RAN agent operates as an xApp in the near-RT RIC, the E2E agent as an rApp in the non-RT RIC/SMO, and the TN and CN agents interact directly with the E2E agent to enable information sharing and coordination. The resulting system comprises four agents operating across three timescales (Fig.~\ref{fig:architecture}):
%etsi21gszsm003
% \begin{figure*}[t]
%     \centering
%     \subfloat[]{%
%         \includegraphics[width=0.75\textwidth]{figures/system-architecture.pdf}%
%         \label{fig:architecture}%
%     }%
%     \hfill %\vrule \hfill
%     \subfloat[]{%
%         \raisebox{1cm}{\includegraphics[width=0.41\textwidth]{figures/timescales.pdf}}%
%         \label{fig:timescales}%
%     }
%     \vspace{-1.8mm}
%     \caption{ENSURE framework. (a) Hierarchical architecture for joint RAN-TN-CN slice management; learning agents have dashed-double-dotted borders, downlink SI flows are color-coded by slice type, and dashed links indicate O-RAN/3GPP control/signaling~\cite{oran23architecture}. (b) \textcolor{red}{Multi-timescale coordination across NRTI, MTI, and LTI intervals. (TO UPDATE)}} %; per-agent functions and exchanged signals are detailed in (b).}
%     % \caption{ENSURE framework. (a)~Hierarchical architecture for joint RAN-TN-CN slice management. Learning agents are marked with dashed-double-dotted borders (vertical placement highlights higher-level control); user icons and example downlink SI paths are color-coded by slice type, where multiple same-color arrows denote distinct SIs. Colored dashed links indicate control/signaling via standard O-RAN/3GPP interfaces~\cite{oran23architecture}. (b)~Multi-timescale coordination across Near-Real-Time (NRTI), Mid-level (MTI), and Long Time (LTI) intervals.}
%     %SDN, and NFV interfaces \caption{ENSURE framework. (a)~Hierarchical architecture for joint RAN, TN, and CN slice management: decision-making agents are marked with double-dotted borders; user colors denote slice types; RAN units/interfaces follow 3GPP/Open RAN standards; CN sites host CPU/GPU servers at edge or cloud locations; TN arrows indicate downlink SI flows with thickness proportional to carried traffic. (b)~Multi-timescale control hierarchy showing temporal coordination across Near-Real-Time (NRTI), Mid-level (MTI), and Long Time (LTI) intervals.} %~\cite{3gpp21tr28811,oran23architecture}
%     \label{fig:framework}
% \end{figure*}

\begin{figure*}[t]
  \centering
  \includegraphics[width=0.91\textwidth]{figures/system-architecture.pdf}
  \caption{ENSURE framework: Hierarchical architecture for joint RAN-TN-CN slice management; learning agents have dashed-double-dotted borders, downlink SI flows are color-coded by slice type, and dashed links indicate O-RAN/3GPP control/signaling~\cite{oran23architecture}.}
  \label{fig:architecture}
\end{figure*}
%%%%%%%%%%%%

% \begin{table}
%     \centering
%     \begin{tabular}{l|l}
%         \toprule
%         \textbf{Abbreviation} & \textbf{Definition} \\ \midrule
%         AI/ML & Artificial Intelligence / Machine Learning \\
%         %AZ & Availability Zone \\        
%         CN & Core Network\\
%         CSI & Channel State Information\\
%         CPU/GPU & Central/Graphical Processing Unit \\
%         CU & Central Unit\\
%         DQN & Deep Q-Network\\
%         DDQN/D3QN & (Dueling) Double DQN \\
%         DU & Distributed Unit\\
%         (MA)-DRL &  (Multi-Agent) Deep Reinforcement Learning\\
%         E2E & End-to-End \\
%         ESRM & E2E Slice Resource Management\\
%         FH & Fronthaul\\
%         LTI & Long Time Interval\\
%         MARSRA & Mobility-Aware RAN Slicing Resource Allocation\\
%         MDP & Markov Decision Process\\
%         MNO/MVNO & Mobile (Virtual) Network Operator\\
%         MPLS & Multi-Protocol Label Switching\\
%         MTI & Mid-level Time Interval\\
%         NRTI & Near-Real-Time Interval\\
%         NIC & Network Interface Controller\\
%         NS &  Network Slicing\\
%         OFDMA & Orthogonal Frequency-Division Multiple Access\\
%         OSL & Overall Satisfaction Level\\
%         PF & Proportional Fairness\\
%         PRB & Physical Resource Block\\
%         proc & processing\\
%         prop & propagation\\
%         PPO & Proximal Policy Optimization\\
%         QoS & Quality of Service\\
%         RAM & Random-Access Memory\\
%         RAN & Radio Access Network\\
%         RA/RM & Resource Allocation/Management\\
%         RIC & RAN Intelligent Controller\\
%         RL & Reinforcement Learning\\
%         RSRQ & Reference Signal Received Quality\\
%         RT & Real-Time\\
%         RU & Radio Unit\\
%         SAC & Soft Actor-Critic\\
%         SDN & Software-Defined Networking\\
%         SI & Slice Instance\\
%         SMO & Service Management and Orchestration\\
%         SotA & State-of-the-Art \\
%         SISL & Slice Instance Satisfaction Level\\
%         TTI & Transmission Time Interval\\
%         tx & transmission\\
%         USL & User Satisfaction Level\\
%         VNF/CNF & Virtual Network Function \\
%     \end{tabular}
%     \caption{Summary of Abbreviations}
%     \label{tab:abbreviations}
% \end{table}
% \subsection{High-Level System Overview}

\begin{itemize}[leftmargin=1.1em,labelsep=0.3em]
\item \textbf{RAN Agent}: Implemented as an xApp and triggered every \textit{Near-Real-Time Interval (NRTI)} (cf. Table~\ref{tab:notations}). Jointly optimizes user association, PRB, and power allocation, subject to user demands, as specified in~\cite{camad24ebrahimi}. % and \textit{transmit power} %user-to-RU association

\item \textbf{TN Agent:} Operates at \textit{Mid-level Time Intervals (MTI)}, selects for each SI an active path from a set of precomputed routes between CN sites and CUs. Feasibility and consistency are enforced via action masking based on link status and bandwidth thresholds. % (Sec.~\ref{subsec:drl-masking}).
%Feasibility is enforced via action masking based on link status and bandwidth thresholds, with an optional training-phase soft mask to enable safe exploration.

\item \textbf{CN Agent:} Operates at the \textit{MTI} level, maps each SI's VNF chain to servers \emph{within} the CN site selected by the E2E agent, and balances per-server load under resource and placement constraints. 
%Each site comprises interconnected CPU/GPU servers, with all VNFs of an SI colocated to meet latency requirements.
%it \textit{places the VNFs} of each \textit{SI} on servers within the chosen CN site, balancing load under resource constraints. Each site comprises interconnected CPU and GPU servers, with all VNFs of an SI colocated to meet latency requirements. % The E2E agent periodically selects the hosting site for each SI's VNFs, while the CN agent performs the corresponding intra-site server assignment.

\item \textbf{E2E Agent:} A top-level rApp in the non-RT RIC/SMO acting at the \textit{Long Time Interval (LTI)}. Selects the hosting CN site for each SI and exchanges lightweight coordination summaries with domain agents (e.g., per-SI satisfaction, per-site utilization) to maintain cross-domain alignment.
%exchanges low-dimensional coordination summaries with domain agents % per-path congestion
%Each SI groups users of the same slice type within a given CU, so TN/CN/E2E decisions scale with the number of SIs rather than the number of users.
%\item \textbf{E2E Agent}: A high-level orchestrator acting at the \textit{Long Time Interval (LTI)}. It aggregates performance indicators and abstracted state information from the RAN, TN, and CN agents, \textit{selects the hosting CN site} for each \textit{SI}, and provides cross-domain guidelines so that domain-level RAs remain aligned with global slice objectives. Each SI aggregates users of the same slice type within a given CU, which keeps TN, CN, and E2E decisions scalable.
\end{itemize}
This distributed hierarchy enables scalable E2E slice management over heterogeneous timescales (Fig.~\ref{fig:timescales-integrated}) with near-RT RAN control and slower TN/CN coordination.
%This distributed, hierarchical design enables scalable E2E slice management operating over heterogeneous timescales (Fig.~\ref{fig:timescales}) with near-real-time RAN control and slower TN/CN coordination. ENSURE instantiates four cooperative agents with explicit multi-timescale coupling and satisfaction metrics aggregated into E2E rewards. %Unlike prior proposals that remain conceptual~\cite{mnet23liu} or RAN-centric~\cite{tnsm24lu,icc24mhatre,tnsm22sulaiman}, 

% \subsection{Multi-Timescale Resource Management}
% The framework couples three control intervals NRTI, MTI, and LTI (see Table~\ref{tab:notations}). We assume downlink numerology~0 with 15~kHz subcarrier spacing \cite{3gpp23ts23501}, so each TTI lasts $1$~ms. Every NRTI spans $\mathscr{N}=50$ TTIs, during which the RAN agent keeps the user-to-RU association and PRB allocation fixed. % , and transmit power
%Users arrive and depart over time and are associated with a single slice type, which determines their QoS requirements \cite{comst24ebrahimi}; the formal model is given in Sec.~\ref{system-model-section}.

% RAN actions are evaluated via the satisfaction metrics from \cite{camad24ebrahimi} and Sec.~\ref{system-model-section}, which penalize overprovisioning relative to slice rate demands and under-provisioning that leads to SLA violations. TN and CN agents act at each MTI (e.g., $500$~ms) to update TN paths and CN VNF placement in response to slow variations in traffic, failures, or resource caps. The E2E agent operates at the LTI (e.g., $1$~s), using aggregated RAN, TN, and CN statistics to refine SI-to-site placement and provide high-level guidance to domain agents.
% Resource allocations aim to ensure timely packet delivery. If allocated PRBs are sufficient to deliver all packets within $\mathscr{N}$ TTIs, the decision is acceptable. If resources exceed a defined threshold, the RAN agent receives a penalty, preventing excessive allocation to a single user at the cost of others' satisfaction. Conversely, insufficient resources, causing SLA violations, penalize user-level satisfaction scores as detailed in \cite{camad24ebrahimi}. The TN and CN orchestrators operate at MTIs (e.g., 500 ms), executing slower adjustments such as TN path updates and CN VNF placement. The E2E orchestrator acts at the LTI (e.g., 1000 ms), refining cross-domain SI-to-site placement based on domain performance and constraints.

% \subsection{Comparison to Related Architectures}
% Several works propose hierarchical or multi-timescale DRL architectures for NS. The framework in \cite{mnet23liu} advocates hierarchical DRL for E2E NS but remains largely conceptual, without concrete RAN, TN, and CN agents or explicit timescale coupling. Multi-level controllers in \cite{tnsm24lu,icc24mhatre} are instantiated within the RAN (e.g., intra-slice xApps and inter-slice rApps in an O-RAN setting), while TN and CN functions remain static or abstracted. Cloud–edge collaborative optimization in \cite{tvt25wang} employs two timescales but treats the network as a monolithic system rather than a domain-oriented stack.

% In contrast, the proposed framework instantiates four cooperative agents tied to standardized control entities: a RAN xApp, TN and CN agents, and an E2E orchestrator aligned with standard management functions \cite{etsi21gszsm003,3gpp21tr28811,oran23architecture}. All domains expose slice-aware abstractions and domain-specific satisfaction metrics aggregated into E2E rewards, allowing the E2E agent to steer long-term policies while local agents adapt on their timescales. This integration of domain abstractions, a three-level timing hierarchy, and MA-DRL coordination yields a scalable, standard-compliant E2E architecture. %a near-RT RAN xApp, TN and CN domain orchestrators, and an LTI-scale E2E agent aligned with standard management functions \cite{etsi21gszsm003,3gpp21tr28811,oran23architecture}.
%%%%%%%%%%%%
\section{System Model and Problem Formulation} \label{system-model-section}
Building on the architectural framework in Sec.~\ref{section-architectural-framework}, this section formalizes the E2E slicing system model and the associated domain-specific optimization problems. We model the downlink service path across the RAN, TN, and CN, spanning multiple timescales. Users are grouped into slice types with heterogeneous QoS requirements, and their traffic is aggregated into SIs that couple user-level fast-changing uncertainties (e.g., mobility and channel quality) in the RAN with SI-level decisions in the TN and CN. Table~\ref{tab:notations} summarizes the main symbols used throughout this section. Unless stated otherwise, $t$ is a generic time index; when a specific timescale matters, we use $t^{\mathsf{T}},t^{\mathsf{N}},t^{\mathsf{M}},t^{\mathsf{L}}$ as specified in Table~\ref{tab:notations}. %specific timescales use $t^{\mathsf{T}},t^{\mathsf{N}},t^{\mathsf{M}},t^{\mathsf{L}}$ (Table~\ref{tab:notations}).

%Slice Instances (SIs)
% As outlined in Sec.~\ref{section-architectural-framework}, the downlink resource management framework spans the RAN, TN, and CN domains \cite{comst24ebrahimi}, jointly allocating resources to meet slice-level E2E requirements. In the RAN, the objective is to maximize user QoS in accordance with slice demands \cite{camad24ebrahimi}, while the TN and CN domains operate at a higher abstraction level, targeting stability, resilience, delay minimization, and compliance with resource constraints.

\subsection{User, Slice, and SI Model} \label{user-slice-model-subsection}
We consider a set of slice types $\mathcal{S}_{\varsigma}=\{\varsigma\}_{1}^{|\mathcal{S}_{\varsigma}|}$, each characterized by a requirement tuple
\begin{align}
\label{slice-req-definition}
\mathrm{Req}_{\varsigma}
=
\{\Bar{B}_{\varsigma}, \Bar{D}_{\varsigma}, \Bar{\lambda}_{\varsigma}, \Bar{L}_{\varsigma},
\boldsymbol{\xi}_{\varsigma},
\varphi^{\mathrm{tol}}_{\varsigma},
\boldsymbol{v}_{\varsigma}\},
\end{align}
%where components are defined in Table~\ref{tab:notations} and detailed in subsequent subsections. 
where $\Bar{B}_{\varsigma}$ and $\Bar{D}_{\varsigma}$ denote the minimum rate and maximum tolerable delay, $\Bar{\lambda}_{\varsigma}$ and $\Bar{L}_{\varsigma}$ specify the packet arrival rate and mean packet length, $\boldsymbol{\xi}_{\varsigma}$ collects slice-specific shaping exponents used in the satisfaction model of Sec.~\ref{subsec:satisfaction-model}, and $\varphi^{\mathrm{tol}}_{\varsigma}\in(0,1]$ is $\varsigma$'s \emph{failure-tolerance} parameter that sets how strictly service disruptions are penalized (cf. TN path-status satisfaction in Sec.~\ref{tn-problem-formulation}).$\boldsymbol{v}_{\varsigma}$ is the tuple describing \textit{VNF chain} (ordered CN functions) and per-function resource demands. 
%\textit{VNF chain} (i.e., the ordered set of CN functions required by $\varsigma$)
%Moreover, The set $\boldsymbol{\xi}_{\varsigma}=\{\xi^{\mathrm{rate}}_{\varsigma}, \xi^{\mathrm{delay}}_{\varsigma}, \xi^{\mathrm{over.}}_{\varsigma}, \xi^{\mathrm{cong.}}_{\varsigma}, \xi^{\mathrm{status}}_{\varsigma}, \xi^{\mathrm{feas.}}_{\varsigma}, \xi^{\mathrm{stab.}}_{\varsigma}\}$ collects the slice-specific shaping exponents (cf. Sec.~\ref{subsec:satisfaction-model}) used to tune the various domain objectives jointly.

Users are denoted by $u\in\mathcal{U}$ and arrive at Transmission Time Interval (TTI) index $t_u^{\mathrm{start}}\in\mathcal{T}^{\mathsf{T}}$, remaining active until $t_u^{\mathrm{end}}>t_u^{\mathrm{start}}$. Each active user belongs to exactly one slice type $\varsigma$ and inherits the corresponding $\mathrm{Req}_{\varsigma}$. %; we write $u$ when the slice type is clear from context and $u_{\varsigma}$ only when needed.
Each agent acts at its own decision epochs (Table~\ref{tab:notations}): RAN every NRTI, TN/CN every MTI, E2E every LTI.
% Each agent acts only at its own decision epochs, i.e., the discrete instants at which it observes the system and updates its action (Table~\ref{tab:notations}). In particular, the RAN agent updates its actions every NRTI and applies them over the $\mathscr{N}$ constituent TTIs, while TN/CN actions update every MTI.
%User $u$ is active at TTIs $t^{\mathsf{T}}$ such that $t_u^{\mathrm{start}}\let^{\mathsf{T}}<t_u^{\mathrm{end}}$.
%Inspired by Open RAN architecture~\cite{oran23architecture}, fine-grained radio dynamics are handled near-real-time, while higher-layer adaptations operate more slowly. %This reflects the practical separation of timescales inspired by the Open RAN architecture \cite{oran23architecture}, where fine-grained radio dynamics are handled near-RT, while higher-layer adaptations operate more slowly.

To make TN/CN/E2E control scale with $|\mathcal{I}|$ rather than $|\mathcal{U}|$, the \textit{Slice Instance (SI)} abstraction has SI $i\in\mathcal{I}$ instantiate slice type $\varsigma(i)$ at CU $c(i)$ and serve the active user set $\mathcal{U}_i(t)$, with aggregate demand as SI traffic. Per-SI decisions are: E2E selects a hosting CN site, TN selects a path to $c(i)$, and CN places the SI's VNF chain \emph{within} the chosen site.
% To make TN/CN/E2E control scale with $|\mathcal{I}|$ slice instances rather than $|\mathcal{U}|$ individual users, we introduce the \textit{Slice Instance (SI)} abstraction: an SI $i\in\mathcal{I}$ represents slice type $\varsigma(i)$ at CU $c(i)$, serving the active user set $\mathcal{U}_i(t)$, whose aggregate demand constitutes the SI traffic. All higher-layer decisions are taken per SI: the E2E agent selects a hosting CN site, the TN agent selects a path to $c(i)$, and the CN agent places the SI's VNF chain \emph{within} that site.
%
% OH = to make it shorter = applies to other places = do you need thee math here?
%
%The \textit{Slice Instance (SI)} abstraction couples per-user RAN control with SI-level TN/CN/E2E control to serve the active user.

% To couple per-user RAN control with SI-level TN/CN/E2E control, we use the \textit{Slice Instance (SI)} abstraction: an SI $i\in\mathcal{I}$ represents slice type $\varsigma(i)$ at CU $c(i)$, serving the active user set $\mathcal{U}_i(t)$. TN/CN/E2E decisions are taken per SI: the E2E agent selects a hosting CN site, the TN agent selects a path to $c(i)$, and the CN agent places the SI’s VNF chain \emph{within} that site. Hence, SI traffic is the aggregate demand of $\mathcal{U}_i(t)$, making TN/CN/E2E control scale with $|\mathcal{I}|$ rather than $|\mathcal{U}|$.
%an \emph{SI} is a coordination unit whose carried traffic is the aggregate of its users, not “\emph{traffic}” itself.
%To couple user-level RAN decisions with SI-level TN/CN control, we define a \textit{Slice Instance (SI)} $i\in\mathcal{I}$ as the instantiation of slice type $\varsigma\in\mathcal{S}_{\varsigma}$ at CU $c\in\mathcal{C}$, yielding $\mathcal{I}=\mathcal{C}\times\mathcal{S}_{\varsigma}$. For each SI $i$, we denote its slice type and CU by $\varsigma(i)$ and $c(i)$, respectively, and its associated users at $t$ by $\mathcal{U}_{i}(t)$. Each SI represents the aggregated traffic of its users between the selected CN site and CU $c(i)$, which is served by a colocated VNF chain and routed to the RAN through the TN. %, as detailed in Secs.~\ref{cn-domain-section} and~\ref{tn-domain-section}.

%%%%%%%%
\subsection{Slice-Aware Satisfaction and Reward Aggregation}
\label{subsec:satisfaction-model}
Across all domains, we map heterogeneous QoS indicators into comparable satisfaction scores and aggregate them into scalar rewards. Let $D\in\{\mathrm{RAN},\mathrm{TN},\mathrm{CN},\mathrm{E2E}\}$ denote the domain, and let entities $e\in\mathcal{E}^\mathrm{(D)}$ be \textit{users} for $\mathrm{RAN}$ and \textit{SIs} for $\mathrm{TN}$/$\mathrm{CN}$/$\mathrm{E2E}$. For slice type $\varsigma$, we evaluate metrics $q\in\mathcal{Q}_\mathrm{(D)}$ with realized value $z^{q}_{e}(t)$ and requirement $\Bar{Z}^{q}_{\varsigma}$. A sign parameter $\rho_q\in\{+1,-1\}$ indicates whether the metric $q$ is \textit{Lower-Bounded} (\textit{LB}) ($+1$, e.g., rate) or \textit{Upper-Bounded} (\textit{UB}) ($-1$, e.g., delay). The generic performance ratio is then defined as 
%In domain $D$, we evaluate metrics $q\in\mathcal{Q}^{(D)}$; slice type $\varsigma$ conditions their requirements and shaping via $\Bar{Z}^{q}_{\varsigma}$ and $\xi^{q}_{\varsigma}$ (defined for $q\in\mathcal{Q}^{(D)}$).
%higher values are better ($+$, e.g., bitrate) or lower values are better ($-$, e.g., delay/congestion).
%For brevity, we write $e\in\mathcal{E}^{(D)}$ for a generic entity and use $t$ for the corresponding domain-specific timescale (e.g., $t^\mathsf{N}$).
\begin{equation}
\boldsymbol{\Hat{\chi}}^{\boldsymbol{\rho}}_{q,e}(t)
\triangleq
\left(\frac{z^{q}_{e}(t)}{\Bar{Z}^{q}_{\varsigma}}\right)^{\rho_q},
\label{eq:satisfaction-ratio}
\end{equation}
so that $\Hat{\chi}^{\boldsymbol{\rho}}_{q,e}(t)=1$ means exactly meeting the requirement, values $>1$ indicate over-fulfillment, and values $<1$ indicate violation (after applying $\rho_q$). We map $\boldsymbol{\Hat{\chi}}^{\boldsymbol{\rho}}_{q,e}(t)$ into a satisfaction score in $(0,1)$ using the standard logistic (sigmoid) function~\cite{camad24ebrahimi}: %\cite{twc13bouali}
\begin{equation}
\mathbb{S}^{(D)}_{q,e}(t)
=
\sigma\bigl(\xi^{q}_{\varsigma}\,\ln \boldsymbol{\Hat{\chi}}^{\boldsymbol{\rho}}_{q,e}(t)\bigr),
\label{eq:satisfaction-mapping}
\end{equation}
where $\xi^{q}_{\varsigma}\ge1$ is a slice- and metric-specific shaping exponent that controls the curve steepness. By construction, $0<\mathbb{S}^{(D)}_{q,e}(t)<1$ for $\boldsymbol{\Hat{\chi}}^{\boldsymbol{\rho}}_{q,e}(t)>0$, with $\mathbb{S}^{(D)}_{q,e}(t)=\tfrac{1}{2}$ when $z^{q}_{e}(t)=\Bar{Z}^{q}_{\varsigma}$. Thus, $\Bar{Z}^{q}_{\varsigma}$ acts as a $0.5$-satisfaction pivot, while a larger $\xi^{q}_{\varsigma}$ yields stricter, less elastic slice behavior (an example can be found on~\cite[Fig.~2]{camad24ebrahimi}). Equations \eqref{eq:satisfaction-ratio}--\eqref{eq:satisfaction-mapping} underpin the domain-specific satisfaction metrics introduced later.
%We encode whether each metric is ``\textit{larger-is-better}'' or ``\textit{smaller-is-better}'' by a sign parameter $\rho_q\in\{+,-\}$, where $\rho_q=+$ for \textit{Lower-Bounded} (\textit{LB}) metrics (e.g., bitrate) and $\rho_q=-$ for \textit{Upper-Bounded} (\textit{UB}) metrics (e.g., delay). 
% \begin{equation}
% \small
% \sigma(\xi^{q}_{\varsigma} \ln{(\boldsymbol{\Hat{\chi}}^{\boldsymbol{\rho}}_{q,e})}) \triangleq  \frac{{\boldsymbol{\Hat{\chi}}^{\boldsymbol{\rho}}_{q,e}}^{\xi^{q}_{\varsigma}}}{1+{\boldsymbol{\Hat{\chi}}^{\boldsymbol{\rho}}_{q,e}}^{\xi^{q}_{\varsigma}}} \triangleq  \frac{1}{1+{\boldsymbol{\Hat{\chi}}^{\boldsymbol{\rho}}_{q,e}}^{-\xi^{q}_{\varsigma}}} \triangleq  \frac{1}{1+{e}^{-\xi^{q}_{\varsigma} \ln{(\boldsymbol{\Hat{\chi}}^{\boldsymbol{\rho}}_{q,e})}}} ,
% \label{eq:standard-logistic}
% \end{equation}
%%%%%%%%%%%%%
% \subsubsection*{\underline{Example}} %Example (RAN bitrate)
% For RAN user $u_{i}$ with achieved bitrate $b_{u}(t^\mathsf{N})$ and requirement $\Bar{B}_{\varsigma}$, the ratio $\boldsymbol{\Hat{\chi}}^{+}_{u,\mathrm{rate}} = b_{u}/\Bar{B}_{\varsigma}$ yields $\mathbb{S}^{\mathrm{(RAN)}}_{\mathrm{rate},u} = 0.5$ when $b_u = \Bar{B}_{\varsigma}$, approaching 1 (resp.\ 0) as $b_u$ exceeds (resp.\ falls below) the requirement. The exponent $\xi^{\mathrm{rate}}_{\varsigma}$ controls how sharply each slice type penalizes deviations from its minimum rate requirement $\Bar{B}_{\varsigma}$ (Fig.~\ref{fig:sigmoid-example}). 
% \begin{figure}[t]
%     \centering
%     \includegraphics[width=0.7\linewidth]{figures/sigmoid.pdf}
%     \vspace{-2mm}
%     \caption{Example: Bitrate satisfaction $\mathbb{S}^{\mathrm{(RAN)}}_{\mathrm{rate}}$ versus normalized rate $b/\Bar{B}$ for different shaping exponents $\xi^{\mathrm{rate}}_{\varsigma}$.}
%     \label{fig:sigmoid-example}
%     \vspace{-4mm}
% \end{figure}
% SD: if you decide to keep this figure, remove the examples in it. They are obvious
%%%%%%%%%%%%%%%%%%%%%%%%%%%%%%
%CANDIDATE FOR REMOVAL
% \subsubsection*{\underline{Per-Metric Aggregation Across Entities}} \label{subsec:per-metric-agg}
% To evaluate how a \emph{single} metric $q$ is satisfied across all entities in a domain, we aggregate $\{\mathbb{S}^{(D)}_{q,e}(t)\}_{e\in\mathcal{E}^{(D)}}$ using the geometric mean. This per-metric aggregate is used in Sec.~\ref{sec:results} to present individual satisfaction components (e.g., $S_{\mathrm{delay}}$) and, unlike the arithmetic mean, is more sensitive to poorly served entities because a single low score reduces the aggregate disproportionately.

\subsubsection*{\underline{Aggregation Across Metrics}}\label{subsec:aggregation-metrics}
Given the per-metric satisfactions $\{\mathbb{S}^{(D)}_{q,e}(t)\}_{q\in\mathcal{Q}_\mathrm{(D)}}$ for $e$, we define its overall satisfaction in domain $D$ as a weighted geometric mean:
\begin{align}
\tilde{\mathbb{S}}^{(D)}_{e}(t)
&=
\prod_{q\in\mathcal{Q}_\mathrm{(D)}}
\Big(\mathbb{S}^{(D)}_{q,e}(t)\Big)^{w^{q}_{\varsigma}},
\quad
\sum_{q\in\mathcal{Q}_\mathrm{(D)}}  w^{q}_{\varsigma}=1,\;
w^{q}_{\varsigma}\ge 0.
\label{eq:entity-satisfaction}
\end{align}
Weights $\{w^{q}_{\varsigma}\}$ encode per-slice metric priorities (e.g., larger delay weight for URLLC). The geometric mean penalizes low satisfaction sharply: a single severely under-satisfied metric drags down the aggregate. This embeds QoS as \textit{soft constraints}: sub-$0.5$ scores signal SLA violation and dominate the reward until recovered. Geometric mean is also used in Sec.~\ref{sec:results} to aggregate a single metric $q$ in domain $D$ across entities (e.g., $S^{(D)}_{\mathrm{delay}}=\bigl(\prod_{e}\mathbb{S}^{(D)}_{q,e}\bigr)^{1/|\mathcal{E}^{(D)}|}$) for visualization.
% The weights $\{w^{q}_{\varsigma}\}$ (for $q\in\mathcal{Q}^{(D)}$) encode the importance of each metric for slice type $\varsigma$ (e.g., assigning a larger weight to delay for URLLC). The geometric mean strongly penalizes low satisfaction, discouraging compensation for a violation in one metric with marginal gains in others. This formulation embeds QoS requirements as \textit{soft constraints} within the reward: satisfaction below $0.5$ signals an SLA violation, and multiplicative aggregation ensures such violations sharply reduce the reward, driving the agent to prioritize under-satisfied metrics while still permitting transient violations during adaptation.

%The geometric mean strongly penalizes low satisfaction across metrics, discouraging compensating for a violation in one metric with marginal gains in others.

%This aggregation prevents trading catastrophic degradation in one dimension for small gains in others. %Throughout the paper, a \textit{tilde} on $\mathbb{S}$ denotes a \textit{geometric-mean aggregate} rather than an approximation. Other tilde-marked aggregates explicitly state their aggregation rule when not geometric means. %Because $\log \tilde{\mathbb{S}}^{(D)}_{e}(t)=\sum_{q\in\mathcal{Q}_\mathrm{(D)}}  w^{q}_{\varsigma} \log \mathbb{S}^{(D)}_{q,e}(t),$ a single severely under-satisfied metric (small $\mathbb{S}^{(D)}_{q,e}(t)$) strongly reduces $\tilde{\mathbb{S}}^{(D)}_{e}(t),$ which prevents trading catastrophic degradation in one dimension for small gains in others.

\subsubsection*{\underline{Aggregation Across Entities (Agent Rewards)}} 
 
The scalar reward for domain $D$ at time $t$ is given by:
\begin{align}
r^{(D)}(t)
&=
\left(
\prod_{e\in\mathcal{E}^{(D)}}
\bigl(\tilde{\mathbb{S}}^{(D)}_{e}(t)\bigr)^{\omega_e}
\right)^{1/(\sum_{e} \omega_e)},
\label{eq:domain-reward}
\end{align}
where higher $\omega_e>0$ prioritize entities (e.g., premium SIs). The weight ratios are bounded to ensure fairness among entities. The following subsections instantiate \eqref{eq:satisfaction-ratio}--\eqref{eq:domain-reward} for each domain’s metrics and timescales.
%REMOVED: 'The geometric aggregation discourages solutions that serve only a subset of entities well.' ;;;; as it has repetitive info shared with 'Aggregation Across Metrics'.
%where $\omega_e\ge0$ denote the entity-level weight. Equal weights $\omega_e\equiv1$ recover the uniform geometric mean, so each entity contributes equally to $r^{(D)}(t)$. Larger $\omega_e$ can prioritize selected entities, e.g., assigning higher weights to premium eMBB entities than to standard ones, while keeping weight ratios bounded to avoid starving lower-priority slices. As in the unweighted case, a single poorly served entity can sharply reduce $r^{(D)}(t)$, and maximizing $r^{(D)}(t)$ is equivalent to maximizing a weighted average log-satisfaction, which retains the fairness-promoting behavior of the geometric mean and discourages policies that trade off a subset of entities for a higher arithmetic mean. The following subsections instantiate equations \eqref{eq:satisfaction-ratio}--\eqref{eq:domain-reward} for each domain’s metrics and timescales.

%%%%%%%%

\subsection{RAN Domain} \label{ran-domain-section}
\subsubsection{\underline{Mathematical Model}} \label{ran-model-subsection}
RAN entities follow the 3GPP 7.2x split~\cite{3gpp23ts23501,oran23architecture} with fixed hierarchical RU/DU/CU placement (cf.\ Fig.~\ref{fig:architecture}). The Mobility-Aware RAN Slicing Resource Allocation (MARSRA) xApp, introduced in our previous work~\cite{camad24ebrahimi},
manages NRTI-level radio RA via user--RU association, PRB assignment, and transmit power allocation.
The resulting user rate $b_u(t^\mathsf{N})$ and RAN delay $\tau^{\mathrm{(RAN)}}_{u}(t^\mathsf{N})$ follow from~\cite[Sec.~II]{camad24ebrahimi}.
% At each NRTI $t^\mathsf{N}_i$ (spanning TTIs $t^\mathsf{T}_{i\mathscr{N}},\dots,t^\mathsf{T}_{i\mathscr{N}+\mathscr{N}-1}$), new users are attached to the RU with highest Reference Signal Received Quality (RSRQ) within their DU and PRBs are initialized via a Proportional-Fair (PF) heuristic; at subsequent NRTIs, user–RU associations and PRB allocations are jointly updated. The binary variables $\chi_{r_d}^{u}(t^\mathsf{N})$ and $\rho^{r_d}_{k,u}(t^\mathsf{N})$ denote user–RU association and PRB assignment, respectively, and the per-PRB bitrate $b^{r_d}_{k,u}$ under Orthogonal Frequency-Division Multiple Access (OFDMA) is given in~\cite[eq.~(2)]{camad24ebrahimi}. The total user bitrate is $b_u(t^\mathsf{N})=\sum_{r_d}\sum_{k} b^{r_d}_{k,u}(t^\mathsf{N})$.
% For user $u$, the RAN delay $\tau^\mathrm{(RAN)}_{u}(t^\mathsf{N})$ is dominated by air-interface propagation on the Uu link; delays on short FH/F1 segments and within the Near-RT RIC are negligible~\cite{3gpp23ts38300,oran23architecture}. We approximate
% $\tau^\mathrm{(RAN)}_{u}(t^\mathsf{N})
% \approx
% \tau^\mathrm{prop,Uu}_{u}(t^\mathsf{N})
% =
% \frac{\chi_{r_d}^{u}(t^\mathsf{N})\,\mathscr{d}_{r_d}^{u}(t^\mathsf{N})}{\mathbf{c}},$
% where $\mathscr{d}_{r_d}^{u}$ is the user–RU distance and $\mathbf{c}$ is the speed of light.
% user--RU association $\chi_{r_d}^{u}(t^\mathsf{N})\in\{0,1\}$, PRB assignment $\rho_{k,u}^{r_d}(t^\mathsf{N})\in\{0,1\}$, and transmit power $p_{k,u}^{r_d}(t^\mathsf{N})\ge0$,

\subsubsection{\underline{RAN Subproblem}} \label{ran-problem-formulation}
We instantiate the satisfaction model of Sec.~\ref{subsec:satisfaction-model} with the two SLA dimensions of $\mathrm{Req}_{\varsigma}$ observable at NRTI granularity: (i)~\textit{rate} ($\mathbb{S}^{\mathrm{(RAN)}}_{\mathrm{rate},u}$, LB) and (ii)~\textit{delay} ($\mathbb{S}^{\mathrm{(RAN)}}_{\mathrm{delay},u}$, UB); congestion, status, and feasibility manifest at the SI level and are handled by the TN/CN/E2E agents (Sec.~\ref{tn-problem-formulation}--\ref{e2e-problem-formulation}), whose aggregated signals are shared with the RAN agent so that cross-domain effects surface before they degrade user-level rate or delay. Per-user aggregation follows \eqref{eq:entity-satisfaction} and the domain reward follows \eqref{eq:domain-reward} with equal user weights, giving $R^{\mathrm{(RAN)}}(t^\mathsf{N})=\bigl(\prod_{u}\tilde{\mathbb{S}}^{\mathrm{(RAN)}}_{u}\bigr)^{1/|\mathcal{U}|}$. 
% We instantiate the satisfaction model of Sec.~\ref{subsec:satisfaction-model} with two \emph{per-user} QoS metrics: (i)~\textit{rate} ($\mathbb{S}^{\mathrm{(RAN)}}_{\mathrm{rate},u}$, LB) and (ii)~\textit{delay} ($\mathbb{S}^{\mathrm{(RAN)}}_{\mathrm{delay},u}$, UB). The per-user RAN satisfaction $\tilde{\mathbb{S}}^{\mathrm{(RAN)}}_{u}(t^\mathsf{N})$ aggregates these two scores via~\eqref{eq:entity-satisfaction}, and the RAN-domain reward $R^{\mathrm{(RAN)}}(t^\mathsf{N})$ follows~\eqref{eq:domain-reward} with equal user weights.
% \begin{itemize}[leftmargin=1.1em,labelsep=0.3em]
%     \item \textit{Bitrate ($\mathbb{S}^{\mathrm{(RAN)}}_{\mathrm{rate},u}$)}: $z^{\mathrm{rate}}_{u}=b_u$, $\Bar{Z}^{\mathrm{rate}}_{\varsigma}=\Bar{B}_{\varsigma}$, $\rho_{\mathrm{rate}}=+$, shaping $\xi^{\mathrm{rate}}_{\varsigma}$.
%     \item \textit{Delay($\mathbb{S}^{\mathrm{(RAN)}}_{\mathrm{delay},u}$)}: $z^{\mathrm{delay}}_{u}=\tau^\mathrm{(RAN)}_{u}$, $\Bar{Z}^{\mathrm{delay}}_{\varsigma}=\Bar{D}_{\varsigma}$, $\rho_{\mathrm{delay}}=-$, shaping $\xi^{\mathrm{delay}}_{\varsigma}$.
%     \item \textit{Overprovisioning ($\mathbb{S}^{\mathrm{(RAN)}}_{\mathrm{over},u}$)}: $z^{\mathrm{over}}_{u}=b_u$, $\Bar{Z}^{\mathrm{over}}_{\varsigma}=2\Bar{B}_{\varsigma}$, $\rho_{\mathrm{over}}=-$, shaping $\xi^{\mathrm{over}}_{\varsigma}$.
% \end{itemize}
% The overprovisioning metric treats excessive rate as a \textit{UB} quantity: $\mathbb{S}^{\mathrm{(RAN)}}_{\mathrm{over},u}$ stays close to $1$ for $b_u\lesssim\Bar{B}_{\varsigma}$, equals $0.5$ at $b_u=2\Bar{B}_{\varsigma}$, and decreases for larger $b_u,$ discouraging unnecessary over-allocation \cite{ojcoms24rezazadeh}, while underprovisioned users are already penalized via $\mathbb{S}^{\mathrm{(RAN)}}_{\mathrm{rate},u}$. 
%matching the MARSRA reward in~\cite[eq.~(7)]{camad24ebrahimi} under the unified notation of Sec.~\ref{subsec:satisfaction-model}. 
The RAN subproblem maximizes $R^{\mathrm{(RAN)}}$ at each NRTI subject to unique user--RU association and PRB exclusivity~\cite[eq.~(8)]{camad24ebrahimi}.
%The RAN subproblem maximizes $R^{\mathrm{(RAN)}}$ at each NRTI subject to unique user--RU association and PRB exclusivity, extensively defined in~\cite[eq.~(8)]{camad24ebrahimi}. %This formulation extends~\cite[eq.~(8)]{camad24ebrahimi} by incorporating $\mathbb{S}^{\mathrm{(RAN)}}_{\mathrm{over},u}$. % within the unified satisfaction framework of Sec.~\ref{subsec:satisfaction-model}.

% At each NRTI, the RAN subproblem is
% \begin{tightflalign} %maybe we can only mention that this problem is explained in \cite{camad24ebrahimi}
% \mathsf{P}_\text{RAN}:\quad &
% \max_{\boldsymbol{\chi},\,\boldsymbol{\rho}} 
% ~R^{\mathrm{(RAN)}}(t^\mathsf{N})
% &&
% \label{eq:obj_ran}
% \end{tightflalign}
% \textbf{s.t.}
% \begin{NamedConstraintR}\label{cons:C1-R}
%   \displaystyle
%   \sum_{r_d} \chi_{r_d}^{u}(t^\mathsf{N}) = 1,
%   \quad \forall u,\forall t^\mathsf{N},
% \end{NamedConstraintR}
% \begin{NamedConstraintR}\label{cons:C2-R}
%   \displaystyle
%   \sum_{u} \chi_{r_d}^{u}(t^\mathsf{N})\,\rho_{k,u}^{r_d}(t^\mathsf{N})
%   \le 1,
%   \quad \forall k,r_d,t^\mathsf{N},
% \end{NamedConstraintR}
% \begin{NamedConstraintR}\label{cons:C3-R}
%   \displaystyle
%   \sum_{k,u}
%   \rho_{k,u}^{r_d}(t^\mathsf{N})
%   \le K_{r_d},
%   \quad \forall r_d,t^\mathsf{N},
% \end{NamedConstraintR}
% \begin{NamedConstraintR}\label{cons:C4-R}
%   \displaystyle
%   \chi_{r_d}^{u}(t^\mathsf{N}),\,
%   \rho_{k,u}^{r_d}(t^\mathsf{N})
%   \in\{0,1\},
% \end{NamedConstraintR}
% \noindent enforcing unique RU attachment, OFDMA constraints, and binary allocations. See~\cite[eq.~(8)]{camad24ebrahimi} for the detailed formulation.

%%%%%%%%%%%%%
\begin{table}[htbp]
\centering
\footnotesize
\caption{Summary of Notations.}
\label{tab:notations}
\begin{tabular}{@{}p{0.192\linewidth}p{0.77\linewidth}@{}}
    \toprule
    \multicolumn{2}{c}{\textbf{Time Intervals}} \\
    \midrule
    $\mathcal{T}^{\mathsf{T}}, t^{\mathsf{T}}, \delta^{\mathsf{T}}$ &
    TTI set, index, and duration; RU-level scheduling interval (1 ms for $\mu=0$) \\
    $\mathcal{T}^{\mathsf{N}}, t^{\mathsf{N}}, \delta^{\mathsf{N}}$ &
    NRTI set, index, and duration; xApp acts every $\delta^{\mathsf{N}}=\mathscr{N}\delta^{\mathsf{T}}$ \\
    $\mathcal{T}^{\mathsf{M}}, t^{\mathsf{M}}, \delta^{\mathsf{M}}$ &
    MTI set, index, and duration; TN/CN agents act every $\delta^{\mathsf{M}}=\mathscr{M}\delta^{\mathsf{N}}$ \\
    $\mathcal{T}^{\mathsf{L}}, t^{\mathsf{L}}, \delta^{\mathsf{L}}$ &
    LTI set, index, and duration; E2E agent acts every $\delta^{\mathsf{L}}=\mathscr{L}\delta^{\mathsf{M}}$ \\
    $\mathscr{N}, \mathscr{M}, \mathscr{L}$ &
    Connection parameters linking TTI$\rightarrow$NRTI, NRTI$\rightarrow$MTI, MTI$\rightarrow$LTI \\ % Aggregation factors (or Connection parameters)
    \midrule    
    \multicolumn{2}{c}{\textbf{Topology and Sets}} \\
    \midrule
    $\mathcal{S}_{\varsigma}$ & Slice types; $\varsigma\in\mathcal{S}_{\varsigma}$ \\
    % $q \in \mathcal{Q}_\mathrm{(D)}$ & Index and set of QoS metrics for domain/agent $D$ \\
    $\mathcal{U},U,u$ & Users (set, cardinality, index) \\
    $\mathcal{I}$ & SIs; SI $i\in\mathcal{I}$ is the instantiation of slice type $\varsigma(i)$ at CU $c(i)$. \\
    $\mathcal{U}_{i}(t)$ & Users belonging to SI $i$ at time $t$ \\
    % $\mathcal{C},C,c$ & CUs (set, cardinality, index) \\ 
    % $\mathcal{D}_{c},D_{c},d_{c}$ & DUs attached to CU $c$ \\
    % $\mathcal{R}_d,R_d,r_d$ & RUs under DU $d$ \\    
    % $\mathcal{C},\mathcal{D}_{c},\mathcal{R}_d$ & Sets of CUs, DUs attached to CU $c$, and RUs under DU $d$ \\ 
    % $\mathcal{K}_{r_d},K,k$ & PRBs (set, cardinality, index) of RU $r_d$  \\
    $\mathcal{Z},\mathcal{N}_z$ & Set of CN sites, Set of servers in $z$ \\
    % $\mathcal{Z},z$ & CN sites (edge or cloud) \\
    % $\mathcal{N}_z$ & Servers in CN site $z$ \\
    $\mathcal{N}^{(\text{TN})},\mathcal{L}^{(\text{TN})}$ & TN nodes and links; a link between $n$ and $n'$ is denoted as $(n,n')$. \\
    $\widehat{\mathcal{G}}^{\text{(TN)}}_{m,m'}$ & Candidate TN paths between nodes $m,m'$ \\
    \midrule
    \multicolumn{2}{c}{\textbf{Slice Requirements}} \\
    \midrule
    $\mathrm{Req}_{\varsigma}, \boldsymbol{v}_{\varsigma}$ & Requirements of slice type $\varsigma$; VNF descriptor of slice $\varsigma$ \\
    $\Bar{B}_{\varsigma},\Bar{D}_{\varsigma}$ & Min. rate, max. delay \\
    $\Bar{\lambda}_{\varsigma},\Bar{L}_{\varsigma}, \varphi^{\mathrm{tol}}_{\varsigma}$ & Packet rate, mean packet size, and failure tolerance \\
    $\boldsymbol{\xi}_{\varsigma}$ & Shaping exponents $\boldsymbol{\xi}_{\varsigma}=\{\xi^{\mathrm{rate}}_{\varsigma}, \xi^{\mathrm{delay}}_{\varsigma}, \xi^{\mathrm{cong.}}_{\varsigma}, \xi^{\mathrm{status}}_{\varsigma}, \xi^{\mathrm{feas}}_{\varsigma}, \xi^{\mathrm{cons}}_{\varsigma}\}$ \\
    % $\varphi^{\mathrm{tol}}_{\varsigma}$ & Failure tolerance \\
    % $\boldsymbol{v}_{\varsigma}$ & VNF descriptor of slice $\varsigma$ \\
    % \midrule
    % \multicolumn{2}{c}{\textbf{RAN Variables}} \\
    % \midrule
    % $\chi_{r_d}^{u}(t^\mathsf{N})$ & User–RU association \\
    % $\rho^{r_d}_{k,u}(t^\mathsf{N})$ & PRB allocation indicator \\
    % $b_u(t^\mathsf{N})$ & Total user bitrate \\ %$b^{r_d}_{k,u}(t^\mathsf{N})
    % $b_u, \, \tau^{\mathrm{(RAN)}}_{u}$ & Total rate and RAN delay of user $u$ \\
    \midrule
    \multicolumn{2}{c}{\textbf{TN and CN Variables (Selected)}} \\
    \midrule
    % $W^{\text{(TN)}}_{(n,n')}, \mathbb{LS}_{(n,n')}, \mathbb{SS}_{n_z}$ & ~~~~~~~~~~~~~~Link capacity and link status (up/down) in TN; Server status in CN \\
    $W^{\text{(TN)}}_{(n,n')}, \pi_{\Hat{g}_{m,m'}}^{i}$ & ~~~~~~~Link capacity in TN; TN path selection variable for SI $i$ \\
    $\mathbb{LS}_{(n,n')}, \mathbb{SS}_{n_z}$ & ~~~~~~~Link and server status (up/down) in TN and CN, respectively \\
    % $\mathbb{SS}_{n_z}$ & Server status in CN \\
    % $\pi_{\Hat{g}_{m,m'}}^{i}$ & TN path selection for SI $i$ \\
    ${\psi}^{i}_{z}$ & SI-to-site assignment (E2E agent) \\
    ${\upsilon}^{i}_{v,n_z}$ & VNF-to-server assignment (CN agent) \\
    \midrule
    \multicolumn{2}{c}{\textbf{Generic Satisfaction Metrics and Reward Definition}} \\
    \midrule
    $z^{q}_{e}(t)$ & Achieved performance of QoS $q$ (e.g., rate) of entity $e$ (user $u_{i}$ or SI $i$) at $t$ \\
    $\,\,\Bar{Z}^{q}_{\varsigma}$ & QoS $q$'s requirement in slice type $\varsigma$ (e.g., $\Bar{D}_{\varsigma}=10$ ms) \\
    $\boldsymbol{\Hat{\chi}}^{\boldsymbol{\rho}}_{q,e}(t)$ & Performance ratio of QoS $q$ for $e$ with sign parameter $\rho_q$ \\
    % $\xi^{q}_{\varsigma}$ & Shaping exponent of $q$ for slice type $\varsigma$ \\
    $\mathbb{S}^{(D)}_{q,e}(t)$ & QoS $q$ satisfaction for $e$ in domain $D$ at $t$ \\
    $\tilde{\mathbb{S}}^{(D)}_{e}(t)$ & Aggregated satisfaction of $e$ in $D$ at $t$ \\
    $r^{(D)}(t)$ & Domain reward (aggregated) in \eqref{eq:domain-reward} \\
    $m_i(t), \Lambda_i(t)$ & Mask-lift indicator at $t$; satisfaction threshold used in masking (cf. Sec.~\ref{subsec:drl-masking})\\
    \bottomrule
\end{tabular}
\end{table}

%==================== TN: MATHEMATICAL MODEL ============
\subsection{TN Domain} \label{tn-domain-section}
\subsubsection{\underline{Mathematical Model}}\label{tn-model-subsection}
We model the TN as an undirected graph $\mathcal{G}^{(\mathrm{TN})}=(\mathcal{N}^{(\mathrm{TN})},\mathcal{L}^{(\mathrm{TN})})$ with symmetric adjacency matrix $\boldsymbol{G}^{(\mathrm{TN})}=[g^{(\mathrm{TN})}_{n,n'}]_{N \times N}$, where $g^{(\mathrm{TN})}_{n,n'}\in\{0,1\}$ indicates whether nodes $n$ and $n'$ are directly connected (i.e., $(n,n')\in\mathcal{L}^{(\mathrm{TN})}$). Each link $(n,n')$ has capacity $W^{(\mathrm{TN})}_{(n,n')}$ (Mbps) and status $\mathbb{LS}_{(n,n')}(t^\mathsf{M})\in\{0,1\}$ (up/down)~\cite{tnse25nouruzi}. TN nodes comprise CN sites $\mathcal{N}^{(\text{core})}$, CU-adjacent edge nodes $\mathcal{N}^{(\text{edge})}$ (ingress/egress points where TN connects to each CU), and intermediate nodes $\mathcal{N}^{(\text{regional})}$, with $\mathcal{N}^{(\mathrm{TN})}=\mathcal{N}^{(\text{core})}\cup\mathcal{N}^{(\text{edge})}\cup\mathcal{N}^{(\text{regional})}$. A subset $\mathcal{N}^{(\text{edgeCN})}\subseteq\mathcal{N}^{(\text{edge})}$ denotes edge nodes that also host CN compute (e.g., site \textit{LA1} in Fig.~\ref{fig:tn-topology}). %(up/down) main citation was~\cite{icc25farooq}
In our downlink model, candidate paths are defined between a CN-side source and a CU-side destination. Specifically, for $m\in\mathcal{N}^{(\text{core})}\cup\mathcal{N}^{(\text{edgeCN})}$ and $m'\in\mathcal{N}^{(\text{edge})}$, $\widehat{\mathcal{G}}^{(\mathrm{TN})}_{m,m'}$ denotes a set of precomputed candidate paths $\hat{g}_{m,m'}$, each represented as a sequence of adjacent links in $\mathcal{L}^{(\mathrm{TN})}$. For SI $i$ at MTI $t^\mathsf{M}$, $\pi^{i}_{\hat{g}_{m,m'}}(t^\mathsf{M})\in\{0,1\}$ selects one path, inducing the active link set $\mathcal{P}^{i}(t^\mathsf{M})\subseteq\mathcal{L}^{(\mathrm{TN})}$. %The local-loop case is captured by allowing $m=m'$ when $m\in\mathcal{N}^{(\text{edgeCN})}$. % sink = destination

%potential new paragraph
The aggregate SI rate is the product of the slice target rate and the average user count $\bar{b}_{i}(t^\mathsf{M})=\Bar{B}_{\varsigma}\,\Bar{U}_{i}(t^\mathsf{M})$. A utilization cap $\theta_{W^{(\mathrm{TN})}}^{\text{link}}\in(0,1]$ limits the total bandwidth demand on each link:
$$
\sum_{i:\,(n,n')\in\mathcal{P}^{i}(t^\mathsf{M})}  \bar{b}_{i}(t^\mathsf{M}) \;\le\; \theta_{W^{(\mathrm{TN})}}^{\text{link}} W^{(\mathrm{TN})}_{(n,n')},~~ \forall (n,n')\in\mathcal{L}^{(\mathrm{TN})},\forall t^\mathsf{M}, %\label{eq:bw-constraint-delta}
$$
and the link utilization is defined as 
$
\mathbb{\Hat{U}}_{(n,n')}(t^\mathsf{M})=\sum_{i:\,(n,n')\in\mathcal{P}^{i}(t^\mathsf{M})} \bar{b}_{i}(t^\mathsf{M})/W^{(\mathrm{TN})}_{(n,n')} .
$

\subsubsection{\underline{Delay in TN}}\label{tn-delay-subsection}
All users of SI $i$ share the path, hence the same TN delay $\tau^{\mathrm{(TN)}}_{i}$, modeled as cumulative propagation delay over $\mathcal{P}^{i}(t^\mathsf{M})$~\cite{tnse25nouruzi}; for local-loop routes ($m=m'\in\mathcal{N}^{(\text{edgeCN})}$), we assign a negligible delay $0<\varepsilon_{\mathrm{(TN)}}\ll1$.
% All users of SI $i$ share the same path and thus the same TN delay $\tau^{\mathrm{(TN)}}_{i}(t^\mathsf{M})$. Since transmission and processing delays are negligible, $\tau^{\mathrm{(TN)}}_{i}$ is modeled as the cumulative propagation delay over links in $\mathcal{P}^{i}(t^\mathsf{M})$, determined by inter-node distances (m) and fiber propagation speed (m/s)~\cite{tnse25nouruzi}. For the local-loop case ($m=m'\in\mathcal{N}^{(\text{edgeCN})}$), we assign a negligible delay $0<\varepsilon_{\mathrm{(TN)}}\ll1$.
%The local-loop case is captured by allowing $m=m'$ when $m\in\mathcal{N}^{(\text{edgeCN})}$.
% All users associated with SI $i$ share the same TN routing and thus the same TN delay. Transmission and processing delays are negligible compared with propagation, so the TN delay is modeled as pure propagation along the selected path 
% $\tau^{\mathrm{(TN)}}_{i}(t^\mathsf{M}) = \sum_{(n,n')\in \mathcal{P}^{i}(t^\mathsf{M})} \mu_{(n,n')}^{i}(t^\mathsf{M})\,\frac{\mathscr{d}_{n,n'}}{\mathbf{c}},$
% where $\mathscr{d}_{n,n'}$ is the physical distance between nodes $n$ and $n'$ (m), and $\mathbf{c}$ is the propagation speed in fiber (m/s). For the co-located case $m=m'\in\mathcal{N}^{(\text{edgeCN})}$ we set $\tau^{\mathrm{(TN)}}_{i}(t^\mathsf{M})=\varepsilon_{\mathrm{(TN)}}$.

%%%%%%%%
\subsubsection{\underline{TN Subproblem}}\label{tn-problem-formulation}
We instantiate four per-SI TN metrics. \textit{Delay} $\mathbb{S}^{\mathrm{(TN)}}_{\mathrm{delay},i}$ follows \eqref{eq:satisfaction-ratio}--\eqref{eq:satisfaction-mapping} with $z^{\mathrm{delay}}_{i}=\tau^{\mathrm{(TN)}}_{i}$, $\Bar{Z}^{\mathrm{delay}}_{\varsigma}=\Bar{D}_{\varsigma}$ (UB). \textit{Congestion} $\mathbb{S}^{\mathrm{(TN)}}_{\mathrm{cong},i}$ uses the bottleneck link utilization $z^{\mathrm{cong}}_{i}=\max_{(n,n')\in\mathcal{P}^{i}}\mathbb{\Hat{U}}_{(n,n')},  \Bar{Z}^{\mathrm{cong}}_{\varsigma}=\theta_{W^{(\mathrm{TN})}}^{\text{link}}$ (UB). %penalizing paths approaching or exceeding the threshold.

\textit{Path-status satisfaction} $\mathbb{S}^{\mathrm{(TN)}}_{\mathrm{status},i}(t^\mathsf{M})$ captures path health based on link status. We define the failure indicator
$$
\mathbb{F}^{i,\mathrm{(TN)}}_{\mathrm{fail}}(t^\mathsf{M}) = \mathbbm{1}\bigl[\textstyle\sum\nolimits_{(n,n')\in\mathcal{P}^{i}}(1-\mathbb{LS}_{(n,n')})>0\bigr] \cdot \Bigl(0.8 + \frac{0.2\sum_{(n,n')\in\mathcal{P}^{i}}(1-\mathbb{LS}_{(n,n')})}{|\mathcal{P}^{i}|}\Bigr), 
%\label{eq:tn-fail-soft}
$$
which equals $0$ for fully healthy paths and lies in $[0.8,1]$ otherwise, so any failure incurs a substantial penalty. We then define the status ratio $\Hat{\chi}^{\boldsymbol{\rho}}_{\mathrm{status},i}=\varphi_{\varsigma}^{\mathrm{tol}}/\mathbb{F}^{i,\mathrm{(TN)}}_{\mathrm{fail}}$, where $0<\varphi_{\varsigma}^{\mathrm{tol}}\le1$ is the slice failure-tolerance factor (larger values penalize failures less). Applying \eqref{eq:satisfaction-mapping} with $\rho_{\mathrm{status}}=+1$ yields $\mathbb{S}^{\mathrm{(TN)}}_{\mathrm{status},i}$, which remains near $1$ for healthy paths and decreases as failures occur.
%\textit{path-health ratio} $\boldsymbol{\Hat{\chi}}^{\boldsymbol{\rho}}_{\mathrm{status},i}=\varphi_{\varsigma}^{\mathrm{tol}}/(\varphi_{\varsigma}^{\mathrm{tol}}+\mathbb{F}^{i,\mathrm{(TN)}}_{\mathrm{fail}})$

\textit{Path-consistency satisfaction} $\mathbb{S}^{\mathrm{(TN)}}_{\mathrm{cons},i}$ discourages unnecessary path switching. Let path quality $\phi_{i}(t)=\min\{\mathbb{S}^{\mathrm{(TN)}}_{\mathrm{cong},i},\mathbb{S}^{\mathrm{(TN)}}_{\mathrm{status},i}\}$ capture the worst-case among congestion and status, and path-change indicator $\Delta^{\mathrm{path}}_{i}(t)\in\{0,1\}$ indicate a path is changed at MTI $t$. The consistency metric is empirically defined as 
$$
z^{\mathrm{cons}}_{\mathrm{(TN)},i}(t) =
\begin{cases}
0.5 + \phi_{i}(t), & \Delta^{\mathrm{path}}_{i}=0,\\[0.2ex]
1 + \phi_{i}(t) - \phi_{i}(t-1), & \Delta^{\mathrm{path}}_{i}=1,\;\phi_{i}(t)\ge\phi_{i}(t-1),\\[0.2ex]
\phi_{i}(t)\,/\,\phi_{i}(t-1), & \Delta^{\mathrm{path}}_{i}=1,\;\phi_{i}(t)<\phi_{i}(t-1).
\end{cases}
$$
When unchanged, the uncongested and healthy paths ($\phi>0.5$) yield $z^{\mathrm{cons}}_{\mathrm{(TN)}}>1$, while marginal ones approach $1$. When changed, improvement and degradation yield $z^{\mathrm{cons}}_{\mathrm{(TN)}}\ge1$ and $z^{\mathrm{cons}}_{\mathrm{(TN)}}<1,$ respec-tively. Applying \eqref{eq:satisfaction-ratio}--\eqref{eq:satisfaction-mapping} with $\rho_{\mathrm{cons}}=+1$ and $\Bar{Z}^{\mathrm{cons}}_{\varsigma}=1$ gives $\mathbb{S}^{\mathrm{(TN)}}_{\mathrm{cons},i}$.

The per-SI TN satisfaction $\tilde{\mathbb{S}}^{\mathrm{(TN)}}_{i}(t^\mathsf{M})$ aggregates the four metrics via \eqref{eq:entity-satisfaction}, and the TN-domain reward $R^{\mathrm{(TN)}}(t^\mathsf{M})$ follows \eqref{eq:domain-reward} as their geometric mean across all SIs. The TN subproblem maximizes $R^{\mathrm{(TN)}}$ at each MTI by selecting one path per SI from the candidate set, subject to link utilization constraints. This combinatorial problem is addressed via DRL in Sec.~\ref{solution-section}.
% \begin{tightflalign}
% \mathsf{P}_\text{TN}:\quad & \max_{\boldsymbol{\pi}}~ R^{\mathrm{(TN)}}(t^\mathsf{M}) && \label{eq:obj_tn}
% \end{tightflalign}
% \textbf{s.t.}
% \begin{NamedConstraintT}\label{cons:C1-T}
% \eqref{eq:bw-constraint-delta}
% \end{NamedConstraintT}
% \begin{NamedConstraintT}\label{cons:C2-T}
% \displaystyle \sum_{\Hat{g}_{m_s,m'_s}} \pi_{\Hat{g}_{m_s,m'_s}}^{i}(t^\mathsf{M}) = 1, \; \foralli,\forall t^\mathsf{M}, m_s\in\mathcal{N}^{(\text{core})}\cup\mathcal{N}^{(\text{edgeCN})}, m'_s\in\mathcal{N}^{(\text{edge})},
% \end{NamedConstraintT}
% \begin{NamedConstraintT}\label{cons:C3-T}
% \displaystyle \pi_{\Hat{g}_{m,m'}}^{i}(t^\mathsf{M}) \in\{0,1\}, \quad \forall i,~\forall \Hat{g}_{m,m'}\in\widehat{\mathcal{G}}^{(\mathrm{TN})}_{m,m'},~\forall t^\mathsf{M}.
% \end{NamedConstraintT}
% Constraint~\ref{cons:C1-T} enforces per-link utilization cap, \ref{cons:C2-T} selects one path per SI and source–destination pair, and \ref{cons:C3-T} enforces binary path selection. 

%%%%%%%%%%%%%%%%%%%%%%%
%==================== CN: MATHEMATICAL MODEL ====================
\subsection{CN Domain} \label{cn-domain-section}

\subsubsection{\underline{Mathematical Model}} \label{cn-model-subsection}
Each slice type $\varsigma$ has an ordered VNF chain $\boldsymbol{v}_{\varsigma}$; each VNF $v\in\boldsymbol{v}_{\varsigma}$ is characterized by
$\mathrm{Req}^{\mathrm{(CN)}}_{v_{\varsigma}}=\{p^{\varsigma}_{\mathrm{tier}},p^{v_{\varsigma}}_{\mathrm{type}},\boldsymbol{r}^{v_{\varsigma}}\}$,
where $\boldsymbol{r}^{v_{\varsigma}}=[\,r^{v_{\varsigma}}_{(j)}\,]_{j=1}^{J}$ collects the CPU, GPU, RAM, storage, and Network Interface Card (NIC) bandwidth demands. The flag $p^{\varsigma}_{\mathrm{tier}}\in\{0,1\}$ encodes the permitted tier ($0$: edge-only, $1$: edge/cloud), while $p^{v_{\varsigma}}_{\mathrm{type}}\in\{0,1\}$ selects CPU- ($0$) or GPU-based ($1$) servers. %, supporting heterogeneous CN nodes.  ($J=5$)

Let $\mathcal{Z}=\mathcal{Z}^{(\mathrm{cloud})}\cup\mathcal{Z}^{(\mathrm{edge})}$ denote CN sites, colocated with TN nodes $\mathcal{N}^{(\text{core})}$ and $\mathcal{N}^{(\text{edgeCN})}$, respectively. Site $z$ hosts a heterogeneous server set $\mathcal{N}_z=\mathcal{N}^{(\mathrm{CPU})}_z\cup\mathcal{N}^{(\mathrm{GPU})}_z$, where each $n_z\in\mathcal{N}_z$ offers capacities $(C^{(1)}_{n_z},\dots,C^{(J)}_{n_z})$ for the $J=5$ resource types. The binary \textit{SI–to-site assignment} variable ${\psi}^{i}_{z}(t^\mathsf{L})$ indicates whether the VNF instances of SI $i$ are placed on site $z$ at LTI $t^\mathsf{L}$ and is chosen by the \textit{E2E agent}. Tier consistency is enforced by setting ${\psi}^{i}_{z}(t^\mathsf{L})=0$ whenever $p^{\varsigma}_{\mathrm{tier}}=0$ and $z\notin\mathcal{Z}^{(\mathrm{edge})}$. At the MTI timescale, the \textit{CN agent} selects the \textit{VNF–server mapping} ${\upsilon}^{i}_{v,n_z}(t^\mathsf{M})\in\{0,1\}$; we impose
${\upsilon}^{i}_{v,n_z}(t^\mathsf{M})\le{\psi}^{i}_{z}(t^\mathsf{L})$
so VNFs can only be placed within the selected site, and set
${\upsilon}^{i}_{v,n_z}(t^\mathsf{M})=0$
for any server–VNF pair that violates $p^{v_{\varsigma}}_{\mathrm{type}}$ (CPU VNFs on GPU-only servers or vice versa). Server status is captured by $\mathbb{SS}_{n_z}(t^\mathsf{M})\in\{0,1\}$, analogous to $\mathbb{LS}_{(n,n')}(t^\mathsf{M})$ in the TN domain.

CPU usage on $n_z$ at $t^\mathsf{M}$ depends on traffic carried by the SIs whose VNFs are hosted on that server. Let $\Bar{L}_{\varsigma}$ be the mean packet size (Mbit) of slice type $\varsigma$ and $\Bar{U}_{i}(t^\mathsf{M})$ the average number of users of SI $i$ over MTI $t^\mathsf{M}$; the induced arrival rate is
$\lambda^{\mathrm{(CN)}}_{i}(t^\mathsf{M})=\Bar{U}_{i}(t^\mathsf{M}) \Bar{B}_{\varsigma}/\Bar{L}_{\varsigma}$.
The CPU capacity and the remaining resource constraints are then %The CPU capacity constraint then reads
$$
% \label{eq:cn-server-cpu-constraint}
\sum_{i}\sum_{v\in\boldsymbol{v}_{\varsigma}} {\upsilon}^{i}_{v,n_z}(t^\mathsf{M})\,r^{v_{\varsigma}}_{(1)}\,\Bar{L}_{\varsigma}\,\lambda^{\mathrm{(CN)}}_{i}(t^\mathsf{M}) \le C^{(1)}_{n_z},~\,\,\,\,\forall n_z\in\mathcal{N}_z,
$$
$$
%while the remaining resources satisfy
% \label{eq:cn-server-resource-constraint}
\sum_{i}\sum_{v\in\boldsymbol{v}_{\varsigma}} {\upsilon}^{i}_{v,n_z}(t^\mathsf{M})\,r^{v_{\varsigma}}_{(j)} \le C^{(j)}_{n_z},~\,\,\,\,j=2,\dots,5,~\,\forall n_z\in\mathcal{N}_z,
$$
respectively. Normalizing the left-hand sides of these equations by $C^{(j)}_{n_z}$ yields utilization ratios $\mathbb{U}^{(j)}_{z,n_z}(t^\mathsf{M})$. % reused in Sec.~\ref{cn-problem-formulation}. %which will be reused in the CN congestion satisfaction.

\subsubsection{\underline{Delay in CN}} \label{cn-delay-subsection}
Propagation and transmission delays within CN sites are negligible, so CN delay is modeled as the sum of VNF processing delays. For each VNF $v\in\boldsymbol{v}_{\varsigma}$ of SI $i$, we approximate the processing delay $\tau^{\mathrm{proc,(CN)}}_{v_{i}}(t^\mathsf{M})$ using an M/M/1 queueing model based on the VNF's service rate and traffic load~\cite{tnse22zhang}, with a large penalty (e.g., $200$\,ms) when load exceeds capacity. The total CN delay is 
$\tau^{\mathrm{(CN)}}_{i}(t^\mathsf{M})=\sum_{v\in\boldsymbol{v}_{\varsigma}} \tau^{\mathrm{proc,(CN)}}_{v_{i}}(t^\mathsf{M}).$

% Propagation and transmission delays within CN sites are negligible owing to short inter-server distances and high intra-site bandwidth, so CN delay is modeled as the \textit{sum of VNF processing delays}. For VNF $v$ of SI $i$ (denoted $v_{i}$) instantiated on server $n_z$, the per-packet service time is
% $s^{\mathrm{(CN)}}_{v_{i}}(t^\mathsf{M})=r^{v_{\varsigma}}_{(j)} \Bar{L}_{\varsigma} / C^{(j)}_{n_z}$, with $j=1$ for CPU- and $j=2$ for GPU-based VNFs. With arrival rate $\lambda^{\mathrm{(CN)}}_{i}(t^\mathsf{M})$, the load is
% $\rho^{\mathrm{(CN)}}_{v_{i}}(t^\mathsf{M})=\lambda^{\mathrm{(CN)}}_{i}(t^\mathsf{M}) s^{\mathrm{(CN)}}_{v_{i}}(t^\mathsf{M})$.
% Under an M/M/1 approximation~\cite{cambridge13harchol,tnse22zhang}, the processing delay of $v_{i}$ is
% $\tau^{\mathrm{proc,(CN)}}_{v_{i}}(t^\mathsf{M})=s^{\mathrm{(CN)}}_{v_{i}}(t^\mathsf{M}) / (1-\rho^{\mathrm{(CN)}}_{v_{i}}(t^\mathsf{M}))$
% if $\rho^{\mathrm{(CN)}}_{v_{i}}(t^\mathsf{M})<1$ and a large penalty value (e.g., $1000$~ms) otherwise. The total CN delay of SI $i$ at $t^\mathsf{M}$ is then 
% $\tau^{\mathrm{(CN)}}_{i}(t^\mathsf{M})=\sum_{v\in\boldsymbol{v}_{\varsigma}} \tau^{\mathrm{proc,(CN)}}_{v_{i}}(t^\mathsf{M})$.
%%%%%%%%%%%%%%
\subsubsection{\underline{CN Subproblem}} \label{cn-problem-formulation}
We instantiate the satisfaction model of Sec.~\ref{subsec:satisfaction-model} with five per-SI metrics: delay, congestion, feasibility, server status, and placement consistency. Metrics are computed per-VNF and aggregated to the SI level conservatively, so the bottleneck VNF dominates.

\textit{Delay satisfaction} $\mathbb{S}^{\mathrm{(CN)}}_{\mathrm{delay},i}$ mirrors the TN case, comparing $\tau^{\mathrm{(CN)}}_{i}$ against $\Bar{D}_{\varsigma}$ (UB). \textit{Congestion satisfaction} $\mathbb{S}^{\mathrm{(CN)}}_{\mathrm{cong},i}$ uses the bottleneck resource utilization across all servers hosting VNFs of SI $i$, compared to a utilization cap $\theta^{\mathrm{(CN)}}_{\mathrm{util}}$ (UB), penalizing placements approaching capacity.

\textit{Feasibility satisfaction} $\mathbb{S}^{\mathrm{(CN)}}_{\mathrm{feas},i}$ enforces hard constraints. For each VNF $v_{i}$, define indicator $z^{\mathrm{feas}}_{v_{i}}\in\{0,1\}$, equal to $0$ if the placement violates tier constraints ($p^{\varsigma}_{\mathrm{tier}}$), server-type constraints ($p^{v_{\varsigma}}_{\mathrm{type}}$), or hard capacity limits, and $1$ otherwise. Applying \eqref{eq:satisfaction-ratio}--\eqref{eq:satisfaction-mapping} with $\rho_{\mathrm{feas}}=+1$, $z^{\mathrm{feas}}_{i}=\min_{v}z^{\mathrm{feas}}_{v_{i}}$, and $\Bar{Z}^{\mathrm{feas}}_{\varsigma}=1$ yields $\mathbb{S}^{\mathrm{(CN)}}_{\mathrm{feas},i}$, which drops to $0$ if any violation occurs; while this binary formulation does not pinpoint specific violations, the CN agent observes per-VNF utilizations, enabling it to learn which resources drive infeasibility.

\textit{Server-status satisfaction} $\mathbb{S}^{\mathrm{(CN)}}_{\mathrm{status},i}$ reuses the TN path-status construction, treating the VNF chain as a path and hosting servers as links. The failure indicator $\mathbb{F}^{i,\mathrm{(CN)}}_{\mathrm{fail}}$ captures the ratio of VNFs on failed servers, transformed as in the TN case (cf.~\ref{tn-problem-formulation}) to yield $\mathbb{S}^{\mathrm{(CN)}}_{\mathrm{status},i}$.

\textit{Placement-consistency satisfaction} $\mathbb{S}^{\mathrm{(CN)}}_{\mathrm{cons},i}$ discourages unnecessary VNF migrations. For each VNF $v_{i}$, let placement quality $\phi_{v_{i}}=\min\{\mathbb{S}^{\mathrm{(CN)}}_{\mathrm{cong},v_{i}},\mathbb{S}^{\mathrm{(CN)}}_{\mathrm{status},v_{i}}\}$. The consistency metric $z^{\mathrm{cons}}_{\mathrm{(CN)},v_{i}}$ follows the same piecewise construction as $z^{\mathrm{cons}}_{\mathrm{(TN)},i}$, but tracks server changes rather than path changes and operates at the VNF level. Applying \eqref{eq:satisfaction-ratio}--\eqref{eq:satisfaction-mapping} with $\rho_{\mathrm{cons}}=+1$ (LB) and $\Bar{Z}^{\mathrm{cons}}_{\varsigma}=1$ yields $\mathbb{S}^{\mathrm{(CN)}}_{\mathrm{cons},v_{i}}$, aggregated as $\mathbb{S}^{\mathrm{(CN)}}_{\mathrm{cons},i}=\min_{v}\mathbb{S}^{\mathrm{(CN)}}_{\mathrm{cons},v_{i}}$. % it rewards stable or improved placements and penalizes unjustified migrations.

The per-SI CN satisfaction $\tilde{\mathbb{S}}^{\mathrm{(CN)}}_{i}$ aggregates the five metrics via \eqref{eq:entity-satisfaction}, and the CN-domain reward $R^{\mathrm{(CN)}}(t^\mathsf{M})$ follows \eqref{eq:domain-reward}.% as their geometric mean across all SIs. 
The CN subproblem maximizes $R^{\mathrm{(CN)}}$ at each MTI by assigning each VNF to a server within the site selected by the E2E agent, subject to resource capacity and tier/type compatibility constraints. This problem is addressed via DRL in Sec.~\ref{solution-section}.
%%%%%%%%%%%%%%%%%%%%
% The CN placement problem $\mathsf{P}_\text{CN}$ is
% \begin{tightflalign}
% \mathsf{P}_\text{CN}:\quad & \max_{\boldsymbol{\upsilon}}~R^{\mathrm{(CN)}}(t^\mathsf{M}) && \label{eq:obj_cn}
% \end{tightflalign}
% \noindent\textbf{s.t.}
% \begin{NamedConstraintC}\label{cons:C1-C}
% \eqref{eq:cn-server-cpu-constraint},~\eqref{eq:cn-server-resource-constraint}
% \end{NamedConstraintC}
% \begin{NamedConstraintC}\label{cons:C2-C}
% {\upsilon}^{i}_{v,n_z}(t^\mathsf{M})\le{\psi}^{i}_{z}(t^\mathsf{L}),\quad\forall i,v,n_z,t^\mathsf{M}
% \end{NamedConstraintC}
% \begin{NamedConstraintC}\label{cons:C3-C}
% {\psi}^{i}_{z}(t^\mathsf{L}),\,{\upsilon}^{i}_{v,n_z}(t^\mathsf{M})\in\{0,1\}
% \end{NamedConstraintC}
% %\noindent where constraint~\ref{cons:C1-C} enforces per-server resource caps, \ref{cons:C2-C} couples SI-to-site and VNF-to-server decisions (implicitly respecting tier/type compatibility), and \ref{cons:C3-C} enforces binary variables; ${\psi}^{i}_{z}(t^\mathsf{L})$ is supplied by the E2E agent and is not optimized here.

%%%%%%%%%%%
\subsection{E2E Domain} \label{e2e-domain-section}
The E2E agent operates at the LTI $t^\mathsf{L}$ and coordinates the three domains by (i) selecting a CN site for each SI via the \textit{SI-to-site mapping} ${\psi}^{i}_{z}(t^\mathsf{L})$ (cf. Sec.~\ref{cn-model-subsection}) and (ii) aggregating domain-specific satisfactions into an E2E perspective. The mapping ${\psi}^{i}_{z}(t^\mathsf{L})$ is reused by the TN agent as the CN-side ingress node of each SI’s TN path and by the CN agent as a site-level placement constraint, and is held fixed over all MTIs within that LTI.
% coompact version (but not having the maths!)
% For compacttime aggregation, $\mathsf{A}^{\mathscr{X}\rightarrow\mathscr{L}}[\cdot](t^\mathsf{L})$ denotes the per-step average over the constituent NRTI ($\mathscr{X}=\mathscr{N}$) or MTI ($\mathscr{X}=\mathscr{M}$) intervals within $t^\mathsf{L}$, and $\mathsf{G}^{\mathrm{user}}_{i}(\cdot)$ denotes the geometric mean over users of SI $i$.
To compactly handle time aggregation, we denote by
$\mathsf{A}^{\mathscr{N}\rightarrow\mathscr{L}}[X](t^\mathsf{L}) \triangleq (\mathscr{N}\mathscr{L})^{-1}\sum_{t^\mathsf{N}\in\mathcal{T}^\mathsf{N}(t^\mathsf{L})}X(t^\mathsf{N})$
the average of an NRTI-based quantity $X(t^\mathsf{N})$ over the $\mathscr{N}\mathscr{L}$ NRTIs within $t^\mathsf{L},$ and by
$\mathsf{A}^{\mathscr{M}\rightarrow\mathscr{L}}[Y](t^\mathsf{L}) \triangleq \mathscr{L}^{-1}\sum_{t^\mathsf{M}\in\mathcal{T}^\mathsf{M}(t^\mathsf{L})}Y(t^\mathsf{M})$
the average of an MTI-based quantity $Y(t^\mathsf{M})$ over the $\mathscr{L}$ MTIs. Furthermore, $\mathsf{G}^{\mathrm{user}}_{i}(\cdot)$ denotes the geometric mean over the users of SI $i$.

%To compactly handle time aggregation, we denote by $\mathsf{A}^{\mathscr{X}\rightarrow\mathscr{L}}[Y]$ the time-average of metric $Y$ over the finer timescale intervals ($\mathscr{X}$ representing either $\mathscr{N}$ (NRTI) or $\mathscr{M}$ (MTI) timescale as appropriate) within an LTI, and by $\mathsf{G}^{\mathrm{user}}_{i}(\cdot)$ the geometric mean over users of SI $i$.

\subsubsection{\underline{Aggregated E2E Satisfaction Metrics}} %\paragraph*{E2E delay and RAN-driven rate/overprovisioning}
For each SI $i$ and NRTI $t^\mathsf{N}$, we approximate the SI-level E2E delay as
$$\tau^{\mathrm{(E2E)}}_{i}(t^\mathsf{N})=\tau^{\mathrm{(RAN)}}_{i}(t^\mathsf{N})+\tau^{\mathrm{(TN)}}_{i}(t^\mathsf{M})+\tau^{\mathrm{(CN)}}_{i}(t^\mathsf{M}),$$
where $\tau^{\mathrm{(RAN)}}_{i}(t^\mathsf{N})$ is the average RAN delay over users of SI $i$ and $t^\mathsf{M}$ is the MTI covering $t^\mathsf{N}$ at which TN/CN delays apply. This sum-of-domain-delays approximation omits transient inter-domain queueing and retransmission effects; these are consistent across all evaluated methods and do not affect their relative comparisons. The LTI-level E2E delay is 
$\tau^{\mathrm{(E2E)}}_{i}(t^\mathsf{L}) = \mathsf{A}^{\mathscr{N}\rightarrow\mathscr{L}}[\tau^{\mathrm{(E2E)}}_{i}](t^\mathsf{L})$ 
mapped to an E2E delay satisfaction
$\mathbb{S}^{\mathrm{(E2E)}}_{\mathrm{delay},i}(t^\mathsf{L})$
via \eqref{eq:satisfaction-ratio}--\eqref{eq:satisfaction-mapping} with
$z^{\mathrm{delay}}_{i}=\tau^{\mathrm{(E2E)}}_{i},$
$\Bar{Z}^{\mathrm{delay}}_{\varsigma}=\Bar{D}_{\varsigma},$
and 
$\rho_{\mathrm{delay}}=-1$.
% potential new paragraph
At each NRTI, the RAN agent computes per-user rate satisfaction
$\mathbb{S}^{\mathrm{(RAN)}}_{\mathrm{rate},u_{i}}(t^\mathsf{N})$ (Sec.~\ref{ran-domain-section}). For SI $i,$ we first obtain SI-level NRTI rate metric as
$\tilde{\mathbb{S}}^{\mathrm{(RAN)}}_{\mathrm{rate},i}(t^\mathsf{N})
=
\mathsf{G}^{\mathrm{user}}_{i}\big(
\mathbb{S}^{\mathrm{(RAN)}}_{\mathrm{rate},u_{i}}(t^\mathsf{N})
\big) ,
$
and then lift it to LTIs via
$\mathbb{S}^{\mathrm{(E2E)}}_{\mathrm{rate},i}(t^\mathsf{L})
=
\mathsf{A}^{\mathscr{N}\rightarrow\mathscr{L}}\bigl[\tilde{\mathbb{S}}^{\mathrm{(RAN)}}_{\mathrm{rate},i}\bigr](t^\mathsf{L}) .
$
%respectively.
%compact (unclear version):
% For SI $i$, we aggregate per-user rate and overprovisioning satisfactions (Sec.~\ref{ran-domain-section}) to SI-level via $\mathsf{G}^{\mathrm{user}}_{i}$, then lift to LTI-level $\mathbb{S}^{\mathrm{(E2E)}}_{\mathrm{rate},i}$ and $\mathbb{S}^{\mathrm{(E2E)}}_{\mathrm{over},i}$ via $\mathsf{A}^{\mathscr{N}\rightarrow\mathscr{L}}$.

%\paragraph*{Cross-domain congestion, feasibility, status, and stability}
The remaining E2E metrics combine SI-level TN and CN satisfactions at MTI granularity and lift them to LTIs via $\mathsf{A}^{\mathscr{M}\rightarrow\mathscr{L}}[\cdot]$. For $q\in\{\mathrm{cong},\mathrm{status},\mathrm{cons}\},$ we set
$$\mathbb{S}^{\mathrm{(E2E)}}_{q,i}(t^\mathsf{L})
=
\mathsf{A}^{\mathscr{M}\rightarrow\mathscr{L}}\bigl[\min\{\mathbb{S}^{\mathrm{(TN)}}_{q,i}(t^\mathsf{M}),
\mathbb{S}^{\mathrm{(CN)}}_{q,i}(t^\mathsf{M})\}\bigr](t^\mathsf{L}),$$ 
so congestion, failures, or inconsistent placements in either domain immediately reduce the corresponding E2E satisfaction. Feasibility remains CN-driven, with
$\mathbb{S}^{\mathrm{(E2E)}}_{\mathrm{feas},i}(t^\mathsf{L})
=
\mathsf{A}^{\mathscr{M}\rightarrow\mathscr{L}}\bigl[\mathbb{S}^{\mathrm{(CN)}}_{\mathrm{feas},i}\bigr](t^\mathsf{L}),$
capturing placement feasibility over the LTI. % reflecting how consistently tier/type and capacity constraints are respected over the LTI.

\subsubsection{\underline{E2E Subproblem}}\label{e2e-problem-formulation}
Collecting the LTI-level E2E metrics
$$\bigl\{\mathbb{S}^{\mathrm{(E2E)}}_{\mathrm{delay}},
\mathbb{S}^{\mathrm{(E2E)}}_{\mathrm{rate}},
\mathbb{S}^{\mathrm{(E2E)}}_{\mathrm{cong}},
\mathbb{S}^{\mathrm{(E2E)}}_{\mathrm{cons}},
\mathbb{S}^{\mathrm{(E2E)}}_{\mathrm{status}},
\mathbb{S}^{\mathrm{(E2E)}}_{\mathrm{feas}}\bigr\}_{i,t^\mathsf{L}} \, ,$$
% we instantiate \eqref{eq:entity-satisfaction} to obtain the E2E SI satisfaction
% $\tilde{\mathbb{S}}^{\mathrm{(E2E)}}_{i}(t^\mathsf{L})$
% as a weighted geometric mean of these metrics, where weights encode priorities (e.g., more delay- and stability-oriented for URLLC, more rate-oriented for eMBB). 
we instantiate \eqref{eq:entity-satisfaction} to obtain the E2E SI satisfaction
$\tilde{\mathbb{S}}^{\mathrm{(E2E)}}_{i}(t^\mathsf{L})$
as a weighted geometric mean of these metrics. Weights are set once from slice SLA priorities (Table~\ref{tab:simulation-settings}) and held fixed across episodes; a systematic sensitivity analysis is deferred to future work.  
The E2E domain reward $R^{\mathrm{(E2E)}}(t^\mathsf{L})$ then follows from \eqref{eq:domain-reward}. % with $\mathcal{E}^{(\mathrm{E2E})}=\mathcal{I}$.
The E2E subproblem maximizes $R^{\mathrm{(E2E)}}$ by assigning a site to each SI subject to SI tier constraints (cf. Sec.~\ref{cn-model-subsection}). Together, LTI-level \textit{SI–site placement} by the E2E agent, MTI-level CN \textit{VNF–server placement} and TN \textit{path selection}, and NRTI-level \textit{user RA} by the RAN agent complete the system model. The following section formulates the domain subproblems as POMDPs and presents an MA-DRL scheme in which domain agents report SI-level metrics, while the E2E agent adapts \textit{SI–site placement} to alleviate cross-domain bottlenecks.
% At each LTI, the E2E agent selects ${\psi}^{i}_{z}(t^\mathsf{L})$ to maximize $R^{\mathrm{(E2E)}}(t^\mathsf{L})$ subject to SI–site consistency and tier constraints:
% $\sum_{z\in\mathcal{Z}}{\psi}^{i}_{z}(t^\mathsf{L})=1$,
% ${\psi}^{i}_{z}(t^\mathsf{L})\in\{0,1\}$,
% and ${\psi}^{i}_{z}(t^\mathsf{L})=0$ for edge-only slices on non-edge sites, as in Sec.~\ref{cn-model-subsection}. Together, LTI-level SI–site placement by the E2E agent, MTI-level CN VNF–server placement and TN path selection, and NRTI-level user RA by the RAN agent complete the system model. The next section formulates $\mathsf{P}_\text{RAN}$, $\mathsf{P}_\text{TN}$, $\mathsf{P}_\text{CN}$, and the E2E placement task as domain-specific POMDPs and outlines a multi-agent learning scheme in which domain agents periodically report SI-level satisfaction metrics and load statistics, while the E2E agent exploits these aggregates to adapt SI–site placement over LTIs, constraining TN routing and VNF placement, and potentially alleviating delay and congestion. % alleviating or exacerbating

%%%%%%%%%%%%%%%%%%
%\subsection{Problem Analysis} \label{problem-analysis-subsection}
%%%%%%%%%
\section{Proposed Solution} \label{solution-section}
The coupled RAN, TN, CN, and E2E subproblems jointly form a large-scale stochastic Mixed-Integer Non-linear Program (MINLP) with nonlinear objectives and tight cross-domain coupling. Classical optimization cannot guarantee timely solutions at the required timescales nor adapt to dynamic conditions without costly re-solving. We therefore adopt a hierarchical MA-DRL framework with four domain-specific agents at their natural control timescales (see Fig.~\ref{fig:timescales-integrated}). The remainder of this section proceeds as follows: Sec.~\ref{subsec:pomdp-models} casts each subproblem as a POMDP; Sec.~\ref{subsec:drl-masking} specifies the learning algorithms and the satisfaction-triggered masking that enforces feasibility at inference; Sec.~\ref{subsec:training} presents the ILP warm start, seeding each episode so agents focus on disturbance recovery rather than cold-start combinatorial search. %describes the ILP-based warm-starting that seeds each episode so that agents focus on disturbance recovery rather than cold-start combinatorial search.

% I expanded the SAC and PPO acronyms here once again as I thought it was necessary (both are previously defined.
A key design choice distinguishes the RAN agent from the others. The RAN domain exhibits high uncertainty (fast fading, user mobility) and high dimensionality (per-user decisions), favoring \textit{Soft Actor-Critic (SAC)} with entropy regularization~\cite{camad24ebrahimi}. In contrast, TN, CN, and E2E operate at the SI level, with slower dynamics, where consistency is paramount (e.g., unnecessary path toggling degrades E2E performance). We adopt \textit{Proximal Policy Optimization (PPO)}~\cite{prep17schulman} (due to its stable updates for discrete decisions) for these agents with \textit{action masking} and a \textit{toggle-based action design}: each agent decides only whether to \textit{change} from the current assignment, reducing the action space and supporting consistency satisfaction. This is enabled by \textit{warm-starting} episodes (cf. Sec.~\ref{subsec:training}). The following subsections formalize the POMDP models, the satisfaction-triggered masking, and the training procedures. % A further simplification is the \textit{toggle-based action design}: each TN/CN/E2E agent decides only whether to \textit{change} from the current assignment rather than selecting among all candidates, strongly reducing the action space and supporting domain-specific stability satisfaction (cf. Sec.~\ref{system-model-section}).%Fig.~\ref{fig:ma-drl-architecture} summarizes how the proposed MA-DRL framework (\textit{ENSURE}) functions.
% \begin{figure}[t]
%     \centering
%     % \includegraphics[width=\columnwidth]{figures/ma-drl-architecture.pdf}
%     \fbox{\parbox{0.95\columnwidth}{\centering\vspace{2cm}\textbf{[Placeholder: MA-DRL Architecture Figure]}\\\small Shows 4 agents, timescales (NRTI/MTI/LTI), and information flows between domains.\vspace{2cm}}}
%     \caption{Hierarchical MA-DRL architecture. The RAN agent (SAC) operates at NRTI scale; TN and CN agents (masked PPO) at MTI scale; E2E agent (masked PPO) at LTI scale. Arrows indicate shared observations and satisfaction feedback.}
%     \label{fig:ma-drl-architecture}
% \end{figure}

\begin{figure*}[t]
  \centering
  \includegraphics[width=1.01\textwidth]{figures/timescales-integrated.pdf}
  \caption{Multi-timescale coordination and per-agent decisions. Nested NRTI$\subset$MTI$\subset$LTI$\subset$Episode containers show decision-making frequencies.} 
  %Nested NRTI$\subset$MTI$\subset$LTI$\subset$Episode containers show the decision frequency of each agent together with its POMDP structure.
  % Left-column arrows indicate bottom-up telemetry and top-down directives (site assignment $\psi_i^z$).% Colors match Fig.~\ref{fig:architecture}(a): RAN (orange), TN (green), CN (purple), E2E (amber, SMO layer).} %(satisfactions $\tilde{\mathbb{S}}$, per-domain SI delays, and per-SI user counts $\bar{U}_i$)
  \label{fig:timescales-integrated}
\end{figure*}
%=============================================================================
\subsection{POMDP Formulations} \label{subsec:pomdp-models}
% [Content for POMDP formulations will go here]
% Subsubsections: RAN, TN, CN, E2E agents
Each agent operates as a Partially Observable Markov Decision Process (POMDP) $(\mathcal{S},\mathcal{A},\mathcal{O},\mathcal{D},\mathcal{R},\gamma)$, with latent state space $\mathcal{S}$, action space $\mathcal{A}$, observation space $\mathcal{O}$, transition dynamics $\mathcal{D}$, reward function $\mathcal{R}$, and discount factor $\gamma \in (0,1)$~\cite{marl24albrecht}. Agents observe local and aggregated cross-domain signals and receive continuous normalized rewards in $[0,1]$. TN/CN/E2E agents use discrete actions with satisfaction-triggered masking (Sec.~\ref{subsec:drl-masking}), whereas the RAN agent~\cite{camad24ebrahimi} uses continuous actions. A summary of multi-timescale orchestration and per-agent POMDP formulations is shown in Fig. \ref{fig:timescales-integrated}. Hereafter, $t$ denotes the agent's native timescale, as defined in Sec.~\ref{section-architectural-framework}. %Partial observability arises because each agent observes only local metrics and aggregated cross-domain feedback, not the full network state.
% Each agent operates as a Partially Observable Markov Decision Process (POMDP) defined by the tuple $(\mathcal{S}, \mathcal{A}, \mathcal{O}, \mathcal{TD}, \mathcal{R}, \gamma)$, where $\mathcal{S}$ is the (hidden) state space, $\mathcal{A}$ the action space, $\mathcal{O}$ the observation space, $\mathcal{TD}$ the transition dynamics, $\mathcal{R}$ the reward function, and $\gamma \in (0,1)$ the discount factor~\cite{marl24albrecht}. Partial observability arises because each agent observes only local metrics and aggregated cross-domain feedback, not the full network state. Observations are continuous and normalized to $[0,1]$; rewards likewise lie in $[0,1],$ to accelerate convergence~\cite{marl24albrecht}. The TN, CN, and E2E agents employ discrete actions with satisfaction-triggered masking (Sec.~\ref{subsec:drl-masking}), while the RAN agent~\cite{camad24ebrahimi} uses continuous actions. Hereafter, $t$ denotes the agent's native timescale index ($t^\mathsf{N}$, $t^\mathsf{M}$, or $t^\mathsf{L}$ as appropriate). %accelerate convergence~\cite{iclr21andrychowicz}.

\subsubsection{\underline{RAN Agent}} \label{subsubsec:pomdp-ran} % [Content for RAN POMDP - reference CAMAD]
The RAN POMDP extends our prior MARSRA formulation in~\cite[Sec.~III-A]{camad24ebrahimi} to incorporate the satisfaction framework of Sec.~\ref{subsec:satisfaction-model}. The observation $\mathbf{o}^{\text{RAN}}_t \in [0,1]^{4U}$ comprises, for each user $u$, satisfaction $\tilde{\mathbb{S}}^{\mathrm{(RAN)}}_{u}(t^{\mathsf{N}}-1)$, PRB fraction, normalized serving-RU index, and channel gain. The action $\mathbf{a}^{\text{RAN}}_t \in [0,1]^{3U}$ encodes, for each user, normalized RU-association preferences, alongside PRB-share and power-share fractions. A heuristic decoder~\cite[Sec.~III-B]{camad24ebrahimi} then maps $\mathbf{a}^{\text{RAN}}_t$ to RAN allocations. The reward is $R^{\mathrm{(RAN)}}(t^\mathsf{N})$ defined in Sec.~\ref{ran-problem-formulation}, balancing rate and delay satisfaction. % v0 wording for observations: allocated PRB fraction, serving RU index (normalized), and normalized channel gain from the serving RU.`

\subsubsection{\underline{TN Agent}} \label{subsubsec:pomdp-tn}
The TN agent manages path selection for each SI at the MTI timescale. For notational convenience, let $L$ denote the number of TN links and $I$ denote the number of SIs. The observation $\mathbf{o}^{\text{TN}}_t \in [0,1]^{4L + IL + 6I}$ comprises: 
(i)~per-link metrics $\bigl(\mathbb{\Hat{U}}_{\text{link}}\in[0,1]^{L}$, binary congestion flags, per-link status (1~healthy, 0~failed), and utilization margin $(\mathbb{\Hat{U}}_{\text{link}}-\theta^{(\mathrm{TN})})$ for anticipating congestion$\bigr)$, 
(ii) per-SI routing state (active route encoding (primary\,$=$\,0, secondary\,$=$\,0.5, tertiary\,$=$\,1), traffic load, link-usage matrix indicating which links each SI's active path traverses, and fraction of path links exceeding threshold), and 
(iii) per-SI previous-step aggregated TN satisfaction $\{\tilde{\mathbb{S}}^{\mathrm{(TN)}}_{i}(t^\mathsf{M}-1)\}_{i\in\mathcal{I}}$ (one scalar per SI) and the current action mask.
% (iii) per-SI $\{\tilde{\mathbb{S}}^{\mathrm{(TN)}}_{i}(t^\mathsf{M}-1)\}_{i\in\mathcal{I}}$ and action mask vectors. 
%Summarized version: per-link metrics (utilization, congestion flags, status, margins), per-SI routing state (active path encoding, traffic load, link-usage matrix, congestion exposure), previous-step satisfactions, and the current action mask."
%The TN agent manages path selection for each SI at the MTI timescale. For notational convenience, let $\mathcal{I}\triangleq\mathcal{I}$ denote the SI set with cardinality $I=|\mathcal{I}|$, and let $L$ denote the number of TN links. The observation $\mathbf{o}^{\text{TN}}_t \in [0,1]^{4L + IL + 6I}$ comprises: (i)~per-link utilization $\mathbb{\Hat{U}}_{\text{link}}\in[0,1]^{L},$ (ii)~binary congestion flags indicating links exceeding a utilization threshold, (iii)~per-link status (1~healthy, 0~failed), (iv)~utilization margin $(\mathbb{\Hat{U}}_{\text{link}}-\theta^{(\mathrm{TN})})$ for anticipating congestion, (v)~per-SI active route encoding (primary\,$=$\,0, secondary\,$=$\,0.5, tertiary\,$=$\,1), (vi)~per-SI normalized traffic load, (vii)~binary link-usage matrix indicating which links each SI's active path traverses, (viii)~per-SI congestion exposure (fraction of path links exceeding threshold), (ix)~per-SI $\{\tilde{\mathbb{S}}^{\mathrm{(TN)}}_{i}(t^\mathsf{M}-1)\}_{i\in\mathcal{I}}$ vector, and (x) the current action mask.

The action $\mathbf{a}^{\text{TN}}_t \in \{0,1\}^{I}$ is binary per SI: $a_i=1$ triggers a toggle to the next precomputed path in the cycle (primary\,$\to$\,secondary\,$\to$\,tertiary\,$\to$\,primary), while $a_i=0$ retains the current path. Candidate paths are precomputed per (site, CU) pair (cf. Sec.~\ref{subsec:training}), with the primary path active by default. 
%The three lowest-delay paths per (site, CU) pair are precomputed at initialization; once the E2E agent assigns an SI to a site, the SI's source–destination pair and candidate paths are fully determined, with the primary path active by default. 
This toggle design reduces the action space from $\mathcal{O}(3^{I})$ joint path selections to only $I$ binary decisions. Action masking (Sec.~\ref{subsec:drl-masking}) permits path changes only under degradation, preserving $\mathbb{S}^{\mathrm{(TN)}}_{\mathrm{cons}}.$ The reward $R^{\mathrm{(TN)}}$ (Sec.~\ref{tn-problem-formulation}) jointly balances delay, congestion, status, and consistency satisfactions.
%Action masking (Sec.~\ref{subsec:drl-masking}) keeps $a_i=0$ forced unless SI~$i$'s current path experiences link failure or congestion, suppressing unnecessary changes that would degrade $\mathbb{S}^{\mathrm{(TN)}}_{\mathrm{cons}}.$ 

\subsubsection{\underline{CN Agent}} \label{subsubsec:pomdp-cn}
The CN agent allocates VNFs to servers within the site determined by the E2E agent. Let $Z$ denote the number of sites, $M_z$ the number of servers at site $z,$ $M=\max_z M_z,$ and $V$ the maximum number of VNFs per SI. The observation $\mathbf{o}^{\text{CN}}_t \in [0,1]^{5MZ + IV + 2I + IZ + 2IVM}$ comprises: 
(i)~per-server utilization across five resource dimensions for all servers at all sites, 
(ii) per-SI state (previous CN action (using normalized server indices), $\{\tilde{\mathbb{S}}^{\mathrm{(CN)}}_{i}(t^\mathsf{M}-1)\}_{i\in\mathcal{I}}$, and $\{\mathbb{S}^{\mathrm{(CN)}}_{\mathrm{cons},i}(t^\mathsf{M}-1)\}_{i\in\mathcal{I}}$), and 
(iii) allocations (current VNF-to-server allocation matrix, SI-to-site assignment from the E2E agent, and the current action mask).
%The CN agent allocates VNFs to servers within the site determined by the E2E agent. Let $Z$ denote the number of sites, $M_z$ the number of servers at site $z,$ $M=\max_z M_z,$ and $V$ the maximum number of VNFs per SI. The observation $\mathbf{o}^{\text{CN}}_t \in [0,1]^{5MZ + 2IVM + IZ + IV + 2I}$ comprises: (i)~per-server resource utilization across five dimensions (CPU, GPU, RAM, storage, NIC) for all servers at all sites, (ii)~the current VNF-to-server allocation matrix encoding which server hosts each VNF of each SI, (iii)~the SI-to-site assignment from the E2E agent, (iv)~the previous CN action (normalized server indices), (v)~per-SI $\{\tilde{\mathbb{S}}^{\mathrm{(CN)}}_{i}(t^\mathsf{M}-1)\}_{i\in\mathcal{I}}$ vector, (vi)~per-SI stability satisfaction $\{\mathbb{S}^{\mathrm{(CN)}}_{\mathrm{cons},i}(t^\mathsf{M}-1)\}_{i\in\mathcal{I}}$, and (vii)~the current action mask.

The action $\mathbf{a}^{\text{CN}}_t \in \prod_{i\in\mathcal{I}} \{0,\ldots,M-1\}^{V_{i}}$ selects a target server index for each VNF, where $V_{i}$ denotes the VNF count for SI~$i$; episodes are warm-started per Sec.~\ref{subsec:training}. Action masking (Sec.~\ref{subsec:drl-masking}) restricts candidate servers to those with sufficient capacity and healthy status; for SI~$i$, migration is disabled while its current placement remains satisfactory and is enabled only upon degradation or infeasibility (e.g., server overload or failure). The reward $R^{\mathrm{(CN)}}$ (Sec.~\ref{cn-problem-formulation}) aggregates delay, utilization, status, feasibility, and consistency terms.
%Action masking restricts choices to servers with sufficient capacity and healthy status; when all satisfaction components for SI~$i$ exceed the pivot threshold (Sec.~\ref{subsec:drl-masking}), the mask locks that SI's VNFs to their current servers, permitting migration only upon performance degradation (e.g., server overload or failure), directly supporting $\mathbb{S}^{\mathrm{(CN)}}_{\mathrm{cons}}$. 

\subsubsection{\underline{E2E Agent}} \label{subsubsec:pomdp-e2e}
The E2E agent assigns each SI to a CN site at the LTI timescale, coordinating cross-domain RA. The observation $\mathbf{o}^{\text{E2E}}_t \in [0,1]^{2IZ + 5Z + L + 4I}$ comprises: (i)~the current SI-to-site allocation matrix, (ii)~per-site utilization means across five resources, (iii)~TN link utilization vector, (iv)~per-domain SI satisfactions $\bar{\mathbb{S}}^{\mathrm{(RAN)}}_i,$  $\tilde{\mathbb{S}}^{\mathrm{(TN)}}_i,$ and $\tilde{\mathbb{S}}^{\mathrm{(CN)}}_i$, (v)~per-SI E2E delay satisfaction $\bigl\{\mathbb{S}^{\mathrm{(E2E)}}_{\mathrm{delay},i}(t^\mathsf{L}-1)\bigr\}_{i\in\mathcal{I}},$ and (vi)~the current action mask. %Let $Z$ denote the number of sites and $L$ the number of TN links.

The action $\mathbf{a}^{\text{E2E}}_t \in \{0,\ldots,Z-1\}^{I}$ selects a site for each SI. This hierarchical decomposition (E2E for SI-to-site, CN for VNF-to-server) reduces the joint decision space from $\mathcal{O}(I \cdot V \cdot M \cdot Z)$ to $\mathcal{O}(I \cdot Z+I \cdot V \cdot M)$, enabling tractable learning. 
%This hierarchical decomposition (E2E for SI-to-site, CN for VNF-to-server) reduces the joint decision space from $\mathcal{O}(IVMZ)$ to $\mathcal{O}(IZ + IVM),$ enabling tractable learning. 
Episodes are initialized with high-quality assignments (cf. Sec.~\ref{subsec:training}). Action masking pins latency-critical SIs to the edge and other SIs to their current site while the E2E-delay satisfaction stays above the pivot (Table~\ref{tab:masking-criteria}); the mask lifts only upon degradation, avoiding migrations that perturb downstream TN/CN assignments.
% Action masking pins latency-critical SIs to the edge and other SIs to their current site while non-congested and delay-compliant (Sec.~\ref{subsec:drl-masking}); the mask lifts only upon degradation, avoiding migrations that perturb downstream TN/CN decisions. 
The reward $R^{\mathrm{(E2E)}}$ is defined in Sec.~\ref{e2e-problem-formulation}.
%Action masking enforces two constraints: (i)~latency-critical SIs remain pinned to the edge site, and (ii)~other SIs remain pinned to their current site while it is non-congested and meets the delay criterion in Sec.~\ref{subsec:drl-masking}. The mask lifts only when delay degrades or the site becomes congested, avoiding unnecessary migrations that would perturb downstream TN/CN decisions. The reward $R^{\mathrm{(E2E)}}$ (Sec.~\ref{e2e-problem-formulation}) integrates the seven aggregated QoS metrics. % harmonizes (instead of integrates)

%=============================================================================
\subsection{DRL Algorithms and Action Masking} \label{subsec:drl-masking}
% [Content: SAC for RAN (brief, ref CAMAD), PPO for TN/CN/E2E, 
%  formal masking mechanism, soft-to-hard schedule, justification vs constrained RL]
The four agents employ two DRL algorithms that match their respective control timescales, action spaces, and consistency requirements. The RAN agent adopts \textit{SAC}~\cite{icml18haarnoja}, whereas the TN, CN, and E2E agents are implemented as separate \textit{PPO} algorithms~\cite{prep17schulman} augmented with \textit{satisfaction-triggered action masking}~\cite{flairs22huang}. This subsection justifies these choices and formalizes the masking mechanism.

\subsubsection{\underline{SAC for the RAN Agent}}
The RAN domain presents high-dimensional, continuous decisions (per-user RU preferences, PRB fractions, power fractions) under rapidly varying channel conditions and user mobility. SAC addresses these challenges via entropy-regularized policy optimization~\cite{icml18haarnoja}, maximizing 
$J_{\text{SAC}}(\pi) = \mathbb{E}_{\pi}\left[\sum_{t}\gamma^{t}\left(R^{\mathrm{(RAN)}}_t + \alpha\,\mathcal{H}\bigl(\pi(\cdot|\mathbf{o}^{\text{RAN}}_t)\bigr)\right)\right],$
% \begin{equation}
% J_{\text{SAC}}(\pi) = \mathbb{E}_{\pi}\left[\sum_{t}\gamma^{t}\left(R^{\mathrm{(RAN)}}_t + \alpha\,\mathcal{H}\bigl(\pi(\cdot|\mathbf{o}^{\text{RAN}}_t)\bigr)\right)\right],
% \label{eq:sac-objective}
% \end{equation}% minimizing the expected KL-divergence
where $\mathcal{H}(\cdot)$ denotes entropy and $\alpha>0$ is the temperature parameter. The entropy bonus encourages sustained exploration, preventing premature convergence to suboptimal associations when channel conditions shift. The continuous action space is handled by Gaussian policy, detailed in~\cite[Sec.~III-B]{camad24ebrahimi}. %The continuous action space is handled by SAC's Gaussian policy, and the heuristic decoder (Sec.~\ref{subsubsec:pomdp-ran}) maps outputs to discrete allocations without gradient discontinuities.

\subsubsection{\underline{PPO for TN, CN, and E2E Agents}}
In contrast, TN, CN, and E2E agents operate on coarser timescales with discrete, change-sensitive decisions (path toggles, server/site assignments). \textit{PPO}~\cite{prep17schulman} is well suited to these settings: its clipped surrogate objective
$L^{\text{CLIP}}(\theta) = \mathbb{E}_t\left[\min\left(r_{t}(\theta)\Hat{A}_t,\;\text{clip}\bigl(r_{t}(\theta),1-\epsilon,1+\epsilon\bigr)\Hat{A}_t\right)\right],$
% \begin{equation}
% L^{\text{CLIP}}(\theta) = \mathbb{E}_t\left[\min\left(r_{t}(\theta)\Hat{A}_t,\;\text{clip}\bigl(r_{t}(\theta),1-\epsilon,1+\epsilon\bigr)\Hat{A}_t\right)\right],
% \label{eq:ppo-clip}
% \end{equation}
where $r_{t}(\theta)=\pi_\theta(a_t|o_t) / \pi_{\theta_{\text{old}}}(a_t|o_t)$ and $\Hat{A}_t$ is the generalized advantage estimate~\cite{prep17schulman}, constrains policy updates to a trust region without second-order overhead. This yields monotonic improvement and prevents policy shifts that would trigger path flapping or frequent VNF migrations. Unlike value-based alternatives (e.g., Deep Q-Network (DQN)), PPO handles the high-dimensional correlated observations of our TN/CN/E2E agents via direct policy parameterization and integrates seamlessly with action masking.
% where $r_{t}(\theta)=\pi_\theta(a_t|o_t) / \pi_{\theta_{\text{old}}}(a_t|o_t)$ and $\Hat{A}_t$ is the generalized advantage estimate~\cite{prep17schulman}, constrains policy updates to a trust region without the computational overhead of second-order methods. This yields monotonic improvement and prevents the significant policy shifts that could destabilize path routing or VNF placement. Unlike value-based alternatives (e.g., Deep Q-Network (DQN)), PPO handles partial observability more gracefully via direct policy parameterization and integrates seamlessly with action masking, as described next. %\cite{iclr16schulman}

\subsubsection{\underline{Satisfaction-Triggered Action Masking}}
Action masking restricts the policy to previously optimized feasible actions by zeroing the probability of undesirable choices \textit{before} sampling. Let $\mathcal{A}$ denote the full action space and $\mathcal{A}^{\text{valid}}_t\subseteq\mathcal{A}$ the subset of allowed actions at step $t$. The masked policy is
$$ % could be inline to further save space (as almost no reference is used)
\pi^{\text{mask}}_\theta(a|o_t) =
\begin{cases}
\displaystyle\frac{\pi_\theta(a|o_t)}{\sum_{a'\in\mathcal{A}^{\text{valid}}_t}\pi_\theta(a'|o_t)}, & a\in\mathcal{A}^{\text{valid}}_t,\\[1.5ex]
0, & \text{otherwise},
\end{cases}
%\label{eq:masked-policy}
$$
which is implemented via setting logits of invalid actions to $-\infty$ prior to softmax normalization, an efficient operation that integrates directly into the actor network's forward pass~\cite{flairs22huang}. %computationally efficient  % policy network's forward pass

%potential new paragraph
The key design principle is \emph{satisfaction-triggered mask lifting}: whereas standard masked PPO lifts only on hard infeasibility~\cite{flairs22huang}, our mask also lifts when the entity's satisfaction $\Lambda_e$ drops below the 0.5 pivot (cf. Sec.~\ref{subsec:satisfaction-model}), triggering reconfiguration \emph{before} a hard violation occurs. Formally, for SI $i$ in TN/CN/E2E at step $t$, the binary mask-lift indicator is $m_i(t)=\mathbf{1}\bigl[\Lambda_i(t)<0.5\;\lor\;\text{infeasibility detected}\bigr]$. 
% Formally, for SI $i$ in TN/CN/E2E at step $t$, we define the binary mask-lift indicator $m_i(t)=1$ if $\Lambda_i(t)<0.5$ or is it really infiseability?
% is detected, and $m_i(t)=0$ otherwise. 
When $m_i(t)=0$, actions that would change entity $i$'s assignment are excluded from $\mathcal{A}^{\text{valid}}_t$; when $m_i(t)=1$, the mask lifts and the agent may explore alternatives. Fig.~\ref{fig:ppo-masking} illustrates the masked PPO architecture, whilst Table~\ref{tab:masking-criteria} summarizes the per-agent instantiation of $m_i(t)$. In all cases, the mask enforces: (i)~\textit{feasibility}, by excluding actions that violate capacity, status, or tier constraints, and (ii)~\textit{reuse of the prior (warm-started or last-committed) assignment}, by suppressing changes when the current satisfaction remains above the pivot. %(ii)~\textit{reuse}, by suppressing changes when performance is acceptable. %This directly supports the stability satisfaction metrics defined in Sec.~\ref{system-model-section}. %, where \textit{infeasibility} captures hard constraint violations (e.g., link failure, server overload).

\begin{figure}[t]
    \centering
    \includegraphics[width=0.99\columnwidth]{figures/ppo-action-masking.pdf}
    % \fbox{\parbox{0.95\columnwidth}{\centering\vspace{2.2cm}\textbf{[Placeholder: PPO with Action Masking]}\\\small Actor network $\to$ logits $\to$ mask application $\to$ masked softmax $\to$ action sampling. Critic network estimates $V(o_t)$. Mask generator uses satisfaction metrics $\{\tilde{\mathbb{S}}_i\}$ and pivot $\tau=0.5$. Show which actions are masked (grayed) vs.\ allowed.\vspace{2.2cm}}}
    \caption{PPO architecture with satisfaction-triggered action masking.} %The mask generator evaluates per-entity satisfaction against $0.5$; actions violating satisfaction criteria are assigned $-\infty$ logits before softmax normalization, ensuring zero sampling probability.}
    \label{fig:ppo-masking}
\end{figure}

\begin{table}[t]
\centering
\small
\caption{Satisfaction-Triggered Masking Criteria per Agent.}
\label{tab:masking-criteria}
\begin{tabular}{@{}p{0.077\columnwidth}p{0.228\columnwidth}p{0.41\columnwidth}@{}}
\toprule
\textbf{Agent} & \textbf{Satisfaction signal $\Lambda_i(t)$} & \textbf{Infeasibility trigger} \\ \midrule
TN  & $\mathbb{S}^{\mathrm{(TN)}}_{\mathrm{cons},i}(t^\mathsf{M})$ & Link failure or $\mathbb{\Hat{U}}_{\text{link}}>\theta^{(\mathrm{TN})}$ on active path \\ \addlinespace[0.5ex]
CN  & $\mathbb{S}^{\mathrm{(CN)}}_{\mathrm{cons},i}(t^\mathsf{M})$ & Server failure or capacity violation \\ \addlinespace[0.5ex]
E2E & $\mathbb{S}^{\mathrm{(E2E)}}_{\mathrm{delay},i}(t^\mathsf{L})$ & Site overutilization or tier violation \\ \bottomrule
\end{tabular}
\end{table}
% \subsubsection{\underline{Example (TN Domain)}}
% Consider SI $i$ on its primary path. If all path links are healthy with $\mathbb{\Hat{U}}_{\text{link}}\le\theta^{(\mathrm{TN})}$, then $\Lambda_i(t)\ge0.5$ and $m_i(t)=0$, masking the toggle action. Upon link failure or congestion, $\Lambda_i(t)$ drops below $0.5$, lifting the mask and permitting a path switch. This mechanism ensures path changes occur only upon genuine performance degradation.
%To illustrate, consider SI $i$ currently routed on its primary path. At each MTI, the TN agent observes per-link utilization and status. If all links on the primary path are healthy and $\mathbb{\Hat{U}}_{\text{link}}\le\theta^{(\mathrm{TN})}$ for each, then $\Lambda_i(t)\ge0.5$ and no infeasibility is detected, yielding $m_i(t)=0$; hence the toggle action ($a_i=1$) is masked and the agent must retain the current path ($a_i=0$). If a link on the primary path fails or becomes congested, $\mathbb{S}^{\mathrm{(TN)}}_{\mathrm{status},i}$ or $\mathbb{S}^{\mathrm{(TN)}}_{\mathrm{cong},i}$ drops below the $0.5$ pivot, so $m_i(t)=1$ and the mask lifts, allowing a switch to the next best path. This mechanism ensures that path changes occur only in response to genuine performance degradation. %, directly minimizing dwell-time penalties captured by $\mathbb{S}^{\mathrm{(TN)}}_{\mathrm{cons}}$.

\subsubsection{\underline{Masking Schedule}}
To avoid over-constraining early exploration, we adopt a \textit{soft-to-hard} annealing schedule that relaxes the mask with probability $\beta(e)$:
$$
\mathcal{A}^{\text{valid}}_t =
\begin{cases}
\mathcal{A}, & \text{w.p. } \beta(e),\\
\mathcal{A}^{\text{mask}}_t, & \text{w.p. } 1-\beta(e),
\end{cases}
\quad \beta(e) = \beta_0\,(1-e/E_{\text{max}})^{\kappa},
$$
with $\kappa>0$ controlling decay curvature. At inference ($\beta=0$) masking is strictly enforced, ensuring feasibility without relying on learned constraint satisfaction.
%where $\kappa>0$ controls decay curvature. At inference, $\beta=0$ and masking is strictly enforced ($\mathcal{A}^{\text{valid}}_t=\mathcal{A}^{\text{mask}}_t$), guaranteeing stability without reliance on learned constraint satisfaction. %feasibility and stability

\subsubsection{\uline{Comparison to Alternative Constraint Handling App- roaches in RL}} % Comparison to Alternative Constraint Handling Approaches in RL
Alternatives include constrained RL via Lagrangian penalties~\cite{jsac24zangooei} and reward shaping~\cite{tnse25nouruzi}. Both introduce additional hyperparameters (penalty coefficients, saddle-point step sizes) and do not prevent violations during exploration. For our discrete-assignment setting with clearly identifiable infeasibility triggers (Table~\ref{tab:masking-criteria}), action masking enforces feasibility directly in the policy network with a single interpretable threshold (the 0.5 pivot from Sec.~\ref{subsec:satisfaction-model}). Combined with warm-start initialization (Sec.~\ref{subsec:training}), it concentrates learning on disturbance recovery. %Combined with warm-start initialization, masking keeps agents focused on disturbance recovery rather than cold-start combinatorial exploration.
% COMPRESSED THIS:
%An alternative to masking is \textit{constrained RL}, which augments the objective with Lagrangian penalties for constraint violations~\cite{jsac24zangooei}. However, such methods require solving saddle-point optimization problems and tuning penalty coefficients, adding hyperparameter sensitivity. \textit{Reward shaping} via penalties can often destabilize training (due to hefty penalties) or fail to prevent small violations~\cite{tnse25nouruzi}. In contrast, action masking \textit{directly} enforces feasibility within the policy network, requiring no additional hyperparameters beyond the satisfaction pivot, which has a natural interpretation (the 0.5 midpoint from Sec.~\ref{subsec:satisfaction-model}). Moreover, masking integrates naturally with warm-start initialization (Sec.~\ref{subsec:training}): because episodes begin with high-quality assignments, the mask remains largely closed early on, and agents learn to handle disturbances (failures, congestion) rather than solving the full combinatorial problem from scratch.

%===========================================
\subsection{Warm-Start Initialization and Training} \label{subsec:training}
Training MA-DRL agents from scratch on combinatorial E2E slicing is sample-inefficient: random initial placements frequently violate constraints. We address this via \textit{warm-start initialization}: each episode is seeded with feasible assignments obtained via lightweight optimization, so agents learn recovery rather than solving the problem from scratch. %refinement and recovery policies

\subsubsection{\underline{Warm-Start for E2E, CN, and TN}}
At the start of each episode, we solve two sequential ILPs. The \textit{E2E warm-start} determines SI-to-site assignments by minimizing aggregate TN propagation delay subject to per-site capacity limits. Given this mapping, the \textit{CN warm-start} places VNFs by minimizing weighted normalized server load, with adaptive weights that increase with utilization to encourage load balancing; constraints enforce site coupling, per-server capacity, and CPU/GPU type compatibility. Both formulations have linear objectives over binary assignment variables and linear constraints, yielding tractable ILPs that can be solved via branch-and-bound with cutting planes using Gurobi in sub-second time. For TN, the three lowest-delay paths per (CN-site, CU) pair are precomputed once at initialization using $k$-shortest simple paths on the TN graph, with the primary path active by default. The E2E and CN ILPs are re-solved \emph{once per episode} (line~2 of Algorithm~\ref{alg:training}); intra-episode adaptation is performed entirely by the masked PPO policies. This partition is deliberate: in our topology, the per-episode ILP takes ${\sim}100$~ms (well within the $500$~ms MTI budget), while per-step DRL inference is ${\sim}10$~ms; re-solving ILPs at every MTI step would scale poorly with $|\mathcal{I}|$ and risk exceeding the MTI budget at larger topologies, whereas DRL inference grows sub-linearly. Warm-starting is also a prerequisite for the satisfaction-triggered mask: without a feasible initialization, the mask would lock agents into random placements rather than enforce feasibility and consistency. Sec.~\ref{sec:results} quantifies their departure from the warm-start under stress.
%Sec.~\ref{sec:results} quantifies how far the agents depart from the warm-start under stress (Fig.~\ref{fig:tn_eval_paths}, Fig.~\ref{fig:combined_bars}).
%For TN, candidate paths are precomputed as described in Sec.~\ref{subsubsec:pomdp-tn}, with the primary path active by default. Subsequent steps inherit these placements; masked PPO policies adjust only upon performance degradation. %subject to per-site capacity limits scaled by utilization threshold $\theta^{(\mathrm{TN})}$. %yielding tractable ILPs solved via branch-and-bound with cutting planes~\cite{wolsey1998integer} using Gurobi~\cite{gurobi}

% === COMMENTED ILP FORMULATIONS FOR REFERENCE ===
% E2E Warm-Start ILP:
% \min_{\boldsymbol{\psi}} \sum_{i\in\mathcal{I}} \sum_{z\in\mathcal{Z}} {\psi}^{i}_{z}\, \widehat{\tau}^{\mathrm{(TN)}}_{z,c}
% where \widehat{\tau}^{\mathrm{(TN)}}_{z,c} is precomputed shortest-path TN delay from site z to CU c
%
% CN Warm-Start ILP:
% \min_{\boldsymbol{\upsilon}} \sum_{i} \sum_{v\in\boldsymbol{v}_{\varsigma}} \sum_{n_z\in\mathcal{N}_z} {\upsilon}^{i}_{v,n_z} \sum_{j=1}^{J} w^{(j)}_{n_z} \frac{\mathscr{D}^{(j)}_{v_{i}}}{C^{(j)}_{n_z}}
% where \mathscr{D}^{(j)}_{v_{i}} is traffic-aware demand, w^{(j)}_{n_z} \ge 1 are adaptive weights
% \subsubsection{\underline{Shortest-Path Warm-Start for TN}}
% Prior to training, we precompute the three lowest-delay candidate paths for each (site, CU) pair using NetworkX's $k$-shortest simple paths algorithm~\cite{networkx08hagberg}, with link weights corresponding to propagation delays. 
% %$O(KN^3)=O(3(Z+C)^3)$ \cite{ms71yen}.
% At the start of each episode, once the E2E agent determines the SI-to-site mapping $\boldsymbol{\psi}$, each SI $i$ is assigned the primary (minimum-delay) path from its allocated site to its associated CU $c$. The TN agent's masked PPO policy subsequently toggles to backup paths only when link failures or congestion cause satisfaction to drop below the pivot threshold, triggering mask lifting as described in Sec.~\ref{subsec:drl-masking}.
%%%%%%%%%%%
\subsubsection{\underline{Training Procedure}}
We adopt an \textit{independent learners} architecture: each agent maintains its own policy and trains separately, while a hierarchical wrapper coordinates cross-timescale interactions. This matches ENSURE's modular design. Algorithm~\ref{alg:training} summarizes the procedure.
%Rather than employing CTDE, we adopt an \textit{independent learners} architecture

\begin{algorithm}[t]
\caption{ENSURE Training Procedure}
\label{alg:training}
\normalsize
\begin{algorithmic}[1]
\Require Randomly initialized agents: RAN (SAC), TN/CN/E2E (PPO); episode count $E_{\max}$; mask relaxation schedule $\beta(e)$
\For{episode $e = 1$ to $E_{\max}$}
    \State Solve E2E and CN warm-start ILPs $\to$ initial $\{\boldsymbol{\psi}^{(0)}, \boldsymbol{\upsilon}^{(0)}\}$ % \cite{wolsey1998integer}
    \State Assign primary TN paths $\{\boldsymbol{\pi}^{(0)}\}$ via shortest-path lookup %~\cite{networkx08hagberg}
    \State Initialize environment; inject random link/server failures
    \For{each LTI $t^\mathsf{L} = 1,\ldots,T^\mathsf{L}$}
        \State E2E agent: observe $\mathbf{o}^{\text{E2E}}$, select $\mathbf{a}^{\text{E2E}}$ (masked)
        \State Update $\boldsymbol{\psi}(t^\mathsf{L})$; broadcast to TN and CN agents
        \For{each MTI $t^\mathsf{M}$ within $t^\mathsf{L}$}
            \State TN agent: observe $\mathbf{o}^{\text{TN}}$, select $\mathbf{a}^{\text{TN}}$ (masked)
            \State CN agent: observe $\mathbf{o}^{\text{CN}}$, select $\mathbf{a}^{\text{CN}}$ (masked)
            \State Update $\boldsymbol{\pi}(t^\mathsf{M})$, $\boldsymbol{\upsilon}(t^\mathsf{M})$; share metrics with E2E agent
            \For{each NRTI $t^\mathsf{N}$ within $t^\mathsf{M}$}
                \State RAN agent: observe $\mathbf{o}^{\text{RAN}}$, select $\mathbf{a}^{\text{RAN}}$
                \State Execute allocation
                \Statex \hspace{4.44em}\phantom{14:}\;Share per-SI user counts with other agents
                \Statex \hspace{4.44em}\phantom{14:}\;Share aggregate metrics with E2E agent                
                %\State Execute allocation; share aggregate user counts per-SI with other agents; share aggregate metrics with E2E agent
            \EndFor
        \EndFor
    \EndFor
    \State Store trajectories; update each agent's policy independently
    \State Decay $\beta(e) \gets \beta_0(1 - e/E_{\max})^\kappa$
\EndFor
\State \Return Trained policies $\{\pi^{\text{RAN}}, \pi^{\text{TN}}, \pi^{\text{CN}}, \pi^{\text{E2E}}\}$
\end{algorithmic}
\end{algorithm}

Each episode initializes with warm-start ILPs (line~2), shortest-path TN assignments (line~3), and randomized traffic/failures for generalization (line~4). The nested loops (lines~5--15) reflect multi-timescale control: the E2E agent acts per LTI, the TN/CN agents per MTI, and the RAN agent per NRTI. Cross-domain coordination uses lightweight exchanges (see Fig.~\ref{fig:timescales-integrated}): SI-to-site mappings flow downward (line~7), satisfaction metrics flow upward (line~11), and per-SI user counts inform TN/CN decisions (line~14). Masking follows Sec.~\ref{subsec:drl-masking} with decaying $\beta(e)$ (line~17). After each episode, SAC updates the RAN, and PPO updates the other agents.
%%%
% Each episode begins by solving the warm-start ILPs (line~2) and initializing TN paths from the precomputed shortest-path table (line~3). The environment is then reset with randomized traffic patterns and injected failures to promote generalization (line~4). The nested loop structure (lines~5--15) reflects the multi-timescale hierarchy: the E2E agent acts once per LTI, TN/CN agents act per MTI, and the RAN agent acts per NRTI. Cross-domain coordination occurs via lightweight information exchange: the E2E agent broadcasts updated SI-to-site mappings (line~7), TN and CN agents report satisfaction metrics to the E2E agent (line~11), and the RAN agent provides aggregated per-SI user counts to inform TN/CN decisions (line~14). Action masking follows Sec.~\ref{subsec:drl-masking}, with relaxation probability $\beta(e)$ decaying over episodes (line~17). After each episode, PPO updates are performed for TN, CN, and E2E agents, and SAC updates for the RAN agent, using their respective collected trajectories (line~16).
% At inference, $\beta=0$ enforces strict masking, and each agent executes its learned policy using only local observations augmented by the shared metrics described above. The warm-start ILPs are solved once at deployment to initialize assignments; subsequent adaptations are handled entirely by the trained policies, enabling real-time response without re-optimization.

% End of Section 5
%%%%%%%%%%

% %%%%%%%%%
% % \section{Computational Complexity and Convergence analysis}\label{complexity-section}

% % \subsection{Computational Complexity}

% % \subsubsection{\underline{Action Selection}

% % \subsubsection{\underline{Training}

% % % \subsubsection{\underline{Meta Learning Phase}

% % \subsection{Convergence analysis}
%%%%%%%%%%%%%%%%%%%%%

% \subsection{Signaling Overhead}
%%%%%%%%%%%%%%%%%%
% \section{Research Challenges and Future Works} \label{challenges-section}

% Future challenges are listed as follows:
% \begin{itemize}
%     \item Having redundancies for the agents (e.g., E2E or domain orchestrators, MARSRA xApp in Near-RT RIC) to combat the single point of failure issue in each domain.
%     \item Include DUs, CUs, and RICs as virtualized entities and orchestrate the RAN resources together with CN and edge resources in a comprehensive cloud orchestration solution.
%     % \item To simplify modeling, we assumed that all VNFs of each SI reside in a single CN site. However, in practice, certain VNFs (e.g., the User Plane Function (UPF)) in some slice types may be deployed in other sites (e.g., edge sites) to reduce E2E latency for user-plane traffic. In future work, we plan to modify this assumption to make our framework more realistic.
% \end{itemize}

%%%%%%%%%%%%%%%
%===============================================
\section{Performance Evaluation} \label{sec:results}

This section evaluates ENSURE through simulation: we describe the setup, present isolated TN stress testing, and conduct a multi-agent ablation study under cross-domain stress.
% This section evaluates ENSURE through simulation. We describe the setup (\S\ref{subsec:setup}), present isolated stress testing that validates the TN agent's recovery mechanisms (\S\ref{subsec:stress}), and conduct an ablation study comparing the full framework against baselines under representative stress scenarios (\S\ref{subsec:ablation}). % degraded variants (baselines)

%===============================================
\subsection{Simulation Setup}
\label{subsec:setup}

\subsubsection{Network Configuration}
% We implement ENSURE on a modified Abilene topology~\cite{networks10orlowski} with 12~nodes and 16~bidirectional fiber links (Fig.~\ref{fig:tn-topology}), using a propagation speed of $2{\times}10^8$~m/s. 
% We implement ENSURE on the Abilene topology~\cite{networks10orlowski} (12 nodes, 16~bidirectional fiber links, Fig.~\ref{fig:tn-topology}); the only modification relative to the reference topology is the addition of a second \textit{Los Angeles} node (\emph{LA2}) as a CU attachment point, so that the system exercises two CUs as required by our model. The TN graph and propagation delays are otherwise unchanged, so comparisons with prior Abilene-based studies remain well-defined. 
We implement ENSURE on the Abilene topology~\cite{networks10orlowski} with one additional CU-attachment node; see Fig.~\ref{fig:tn-topology}. We use a propagation speed of $2{\times}10^8$~m/s. The RAN comprises 2~CUs, 2~DUs per CU, and 2~RUs per DU (8~RUs total) in a $1000{\times}1000$~m$^2$ area, serving up to 40~users equally split across two slice types: eMBB and URLLC~\cite{tvt23choi}. With 2~CUs and 2~slice types, the system has $I=4$~SIs. This user scale is comparable to recent RAN-slicing studies under mobility~\cite{tvt23choi,camad24ebrahimi}; because TN/CN/E2E agents operate at the \emph{SI} rather than the user level (Sec.~\ref{user-slice-model-subsection}), user count primarily affects RAN per-step cost, whereas TN/CN/E2E complexity is set by $|\mathcal{I}|$. Scaling to larger $|\mathcal{I}|$ is flagged in Sec.~\ref{conclusion-section}. 
% serving up to 40~users equally split across two slice types: eMBB and URLLC~\cite{tvt23choi}. With 2~CUs and 2~slice types, the system has $I=4$~SIs. 

We adopt 5G numerology $\mu=0$ (15~kHz subcarrier spacing, 1~ms TTI) in the n77~band with 50~MHz shared bandwidth, yielding 273~PRBs per RU~\cite{3gpp23ts23501}. Per-RU transmit power (45~dBm) is uniformly distributed across active PRBs. Channel gains follow 3GPP~\cite{3gpp24tr38901} with $-174$~dBm/Hz noise power density, path-loss exponent~3, and inter-cell interference~\cite{tvt23choi}. User dynamics extend our prior work~\cite{camad24ebrahimi} with intra-episode velocity changes and stochastic arrivals/departures. The CN includes 3~cloud sites and 1~edge site with heterogeneous CPU/GPU servers as shown in Fig.~\ref{fig:tn-topology}; each SI's VNF chain must be placed on servers within the assigned site, subject to the resource capacities indicated in the figure. TN link and CN server utilization caps ($\theta^{\text{link}}_{W^{(\mathrm{TN})}}$ and $\theta^{(\mathrm{CN})}_{\mathrm{util}}$) are both set to $0.8$. Table~\ref{tab:simulation-settings} summarizes slice requirements, satisfaction shaping parameters, per-agent reward weights, VNF resource demands, and DRL hyperparameters. The framework is implemented using Ray~RLlib, NetworkX, and PyTorch.
%The three shortest paths per (CN~site, CU) pair are listed in Table~\ref{tab:tn-prop}.
%, and mMTC ($D_{\max}=100$\,ms, $R_{\min}=0.1$\,Mbps) %273~PRBs per RU~\cite{3gpp24ts38101}

\begin{figure}
    \centering
    \includegraphics[width=0.99\columnwidth]{figures/topology.pdf}
    \caption{TN/CN topology. \emph{LA1} is both CU and edge CN site; \emph{LA1}$\to$\emph{LA1} uses $\varepsilon_{(\mathrm{TN})}=0.01$~ms. Link labels: propagation delays. CN resource capacities per server: CPU (cycles/s), GPU (MIGs), RAM (GB), storage (GB), NIC (Mbps). Legend abbreviations: ER $=$ edge--regional (50~Gbps), RR $=$ regional--regional (100~Gbps), CR $=$ core--regional (200~Gbps), Ring $=$ edge ring (10~Gbps); $n^{\mathrm{(edgeCN)}}$ $=$ edge node co-located with a CU, RIC, and CN compute.}
    %\caption{TN/CN topology based on Abilene~\cite{networks10orlowski}. \textit{LA1} is both CU and edge CN site; \textit{LA1}$\to$\textit{LA1}\ uses $\varepsilon_{\mathrm{(TN)}}=0.01$\,ms. Link labels: propagation delays. CN capacities: CPU (cycles/s), GPU (MIGs), RAM (GB), storage (GB), NIC (Mbps).}    
    %\caption{{TN/CN Topology. `\textit{LA1}' serves as both CU and edge CN site (with co-located RIC); the co-located path `\textit{LA1}'$\to$`\textit{LA1}' uses $\varepsilon_{\mathrm{(TN)}}=0.01\,\mathrm{ms}$. Link labels show propagation delays; CN resource capacities are CPU (cycles/s), GPU (Multi-Instance GPUs (MIGs)), RAM (GB), Storage (GB), and NIC (Mbps).}}
    \label{fig:tn-topology}
\end{figure}


%%%%% MERGED TABLE %%%%%
% \begin{table}[t]
% \centering
% \caption{Simulation settings: slice requirements, satisfaction parameters, reward weights, VNF demands, and DRL hyperparameters.}
% \label{tab:simulation-settings}
% \resizebox{\columnwidth}{!}{%
% \scriptsize
% \setlength{\tabcolsep}{2pt}
% \renewcommand{\arraystretch}{1.1}
% \begin{tabular}{cl ll}
% \toprule
% \textbf{Param.} & \textbf{Description} & \textbf{eMBB} & \textbf{URLLC} \\
% \midrule
% \multicolumn{4}{c}{\textbf{Slice Requirements}} \\
% \midrule
% $|\mathcal{S}_\varsigma|$ & No.\ of slice types~\cite{icc24mhatre} & \multicolumn{2}{c}{$$2 } \\
% $U_\varsigma/U$ & User fraction per slice & 50\% & 50\% \\
% $\bar{R}_\varsigma$ & Min.\ rate (Mbps) & 16 & 3.8 \\
% $\bar{D}_\varsigma$ & Max.\ E2E delay (ms) & 20 & 2 \\
% $\bar{L}_\varsigma$ & Packet size (Bytes) & 1024 & 480 \\
% \midrule
% \multicolumn{4}{c}{\textbf{Satisfaction Shaping Parameters} (cf.\ Sec.~\ref{subsec:satisfaction-model})} \\
% \midrule
% $\xi^{\mathrm{rate}}_\varsigma,\;\xi^{\mathrm{delay}}_\varsigma,\;\xi^{\mathrm{cong}}_\varsigma$ & Rate, delay, cong.\ shaping & 5,\;2,\;3 & 2,\;5,\;6 \\
% $\xi^{\mathrm{status}}_\varsigma,\;\xi^{\mathrm{cons}}_\varsigma,\;\varphi^{\mathrm{tol}}_\varsigma$ & Status, stability shaping; failure tol. & 3,\;3,\;0.4 & 5,\;5,\;0.15 \\
% \midrule
% \multicolumn{4}{c}{\textbf{Reward Weights} $w^q_\varsigma$ in \eqref{eq:entity-satisfaction} per agent (sum to 1)} \\
% \midrule
% RAN & $w_{\mathrm{rate}},\;w_{\mathrm{delay}}$~\cite{camad24ebrahimi} & .8,\;.2 & .2,\;.8 \\ % \multicolumn{2}{c}{0.89,\;0.11} \\
% TN & $w_{\mathrm{delay}},\;w_{\mathrm{cong}},\;w_{\mathrm{status}}$ & .4,\;.4,\;.2 & .5,\;.1,\;.4 \\
% CN & $w_{\mathrm{delay}},\;w_{\mathrm{cong}},\;w_{\mathrm{status}},\;w_{\mathrm{feas}},\;w_{\mathrm{cons}}$ & .2,\;.2,\;.1,\;.3,\;.2 & .3,\;.2,\;.2,\;.1,\;.2 \\
% E2E & $w_{\mathrm{delay}},\;w_{\mathrm{rate}},\;w_{\mathrm{cong}},\;w_{\mathrm{status}},\;w_{\mathrm{feas}},\;w_{\mathrm{cons}}$ & .1,\;.3,\;.2,\;.1,\;.1,\;.2 & .3,\;.1,\;.2,\;.2,\;.1,\;.1 \\
% \midrule
% % E2E & $\omega_{\mathrm{eMBB}},\;\omega_{\mathrm{URLLC}}$ (entity weights in \eqref{eq:domain-reward}) & \multicolumn{2}{c}{3.0,\;5.0} \\
% $\omega_{\mathrm{eMBB}},\;\omega_{\mathrm{URLLC}}$ & eMBB/URLLC SIs' weight in \eqref{eq:domain-reward} (only applied for E2E reward) & 3 & 5 \\
% \midrule
% \multicolumn{4}{c}{\textbf{VNF Characteristics~\cite{tnse25nouruzi}} $\mathrm{Req}^{\mathrm{CN}}_{\boldsymbol{v}_\varsigma}$\;=\;$(p^\varsigma_{\mathrm{tier}},\,p^{v_\varsigma}_{\mathrm{type}},\,\boldsymbol{r}^{v_\varsigma})^{\dagger}$} \\
% \midrule
% $v_1$ & AMF & 1,\;0,\;8.5,\;0,\;16,\;60,\;1 & 0,\;0,\;6,\;0,\;12,\;50,\;1.5 \\
% $v_2$ & UPF & 1,\;0,\;11,\;0,\;48,\;180,\;2.5 & 0,\;0,\;12,\;0,\;16,\;100,\;3 \\
% $v_3$ & SMF (eMBB) / NWDAF (URLLC) & 1,\;0,\;7.2,\;0,\;12,\;90,\;0.5 & 0,\;1,\;0,\;.00025,\;24,\;50,\;2 \\
% \midrule
% \multicolumn{4}{c}{\textbf{DRL Hyperparameters}} \\
% \midrule
% & Masked PPO (TN/CN/E2E) & \multicolumn{2}{l}{$\alpha=10^{-4}$,\;$\gamma=0.99$,\;$\lambda=0.95$,\;$\epsilon=0.2$} \\
% & \qquad entropy coeff.,\;batch size,\;NN arch. & \multicolumn{2}{l}{0.005,\;512,\;$256{\times}256$} \\
% & SAC (RAN)~\cite{camad24ebrahimi} & \multicolumn{2}{l}{As in~\cite{camad24ebrahimi}} \\
% \bottomrule
% \multicolumn{4}{@{}p{1.55\columnwidth}@{}}{%
% \scriptsize $^{\dagger}$$p^\varsigma_{\mathrm{tier}}$: $$0\,=\,edge-only, $$1\,=\,any tier; $p^{v_\varsigma}_{\mathrm{type}}:$ $$0\,=\,CPU, $$1\,=\,GPU. $\boldsymbol{r}^{v_\varsigma}=$(CPU ($10^6 $ cyc/Mbit), $$GPU $$(s/Mbit), $$RAM $$(GB), $$Storage $$(GB), $$NIC $$(Gbps)).%, matching the resource units in Fig.~\ref{fig:tn-topology}.%
% }
% \end{tabular}%
% }
% \end{table}

%%%%%%%%%%%
% Ensure you have \usepackage{tabularx} and \usepackage{booktabs} in your preamble

\begin{table*}[t]
\centering
\caption{Simulation settings: slice requirements, satisfaction parameters, reward weights, VNF demands, and DRL hyperparameters.}
\label{tab:simulation-settings}
\resizebox{0.99\columnwidth}{!}{%
\footnotesize
\setlength{\tabcolsep}{1.2pt}
\renewcommand{\arraystretch}{1.2} % Adds breathing room between rows
\begin{tabularx}{\textwidth}{@{} c >{\raggedright\arraybackslash}p{6.95cm} >{\raggedright\arraybackslash}X >{\raggedright\arraybackslash}X @{}}
\toprule
\textbf{Param.} & \textbf{Description} & \textbf{eMBB} & \textbf{URLLC} \\
\midrule
\multicolumn{4}{c}{\textbf{Slice Requirements}} \\
\midrule
%$|\mathcal{S}_\varsigma|$, $U_\varsigma/U$ & No.\ of slice types~\cite{icc24mhatre}, User fraction & \multicolumn{2}{c}{2 types total; \quad 50\% eMBB / 50\% URLLC} \\
$\bar{R}_\varsigma, \bar{D}_\varsigma, \bar{L}_\varsigma$ & Min.\ rate (Mbps), Max.\ E2E delay (ms), Packet size (Bytes) & 16, \quad 20, \quad 1024 & 3.8, \quad 2, \quad 480 \\
\midrule
\multicolumn{4}{c}{\textbf{Satisfaction Shaping Parameters} (cf.\ Sec.~\ref{subsec:satisfaction-model})} \\
\midrule
$\xi^{\mathrm{rate}}_\varsigma, \xi^{\mathrm{delay}}_\varsigma, \xi^{\mathrm{cong}}_\varsigma$ & Rate, delay, cong.\ shaping & 5, 2, 3 & 2, 5, 6 \\
$\xi^{\mathrm{status}}_\varsigma, \xi^{\mathrm{cons}}_\varsigma, \varphi^{\mathrm{tol}}_\varsigma$ & Status, consistency shaping; failure tol. & 3, 3, 0.4 & 5, 5, 0.15 \\
\midrule
\multicolumn{4}{c}{\textbf{Reward Weights} $w^q_\varsigma$ in \eqref{eq:entity-satisfaction} per agent (sum to 1)} \\
\midrule
RAN & $w_{\mathrm{rate}}, w_{\mathrm{delay}}$~\cite{camad24ebrahimi} & .8, .2 & .2, .8 \\ 
TN & $w_{\mathrm{delay}}, w_{\mathrm{cong}}, w_{\mathrm{status}}$ & .4, .4, .2 & .5, .1, .4 \\
CN & $w_{\mathrm{delay}}, w_{\mathrm{cong}}, w_{\mathrm{status}}, w_{\mathrm{feas}}, w_{\mathrm{cons}}$ & .2, .2, .1, .3, .2 & .3, .2, .2, .1, .2 \\
E2E & $w_{\mathrm{delay}}, w_{\mathrm{rate}}, w_{\mathrm{cong}}, w_{\mathrm{status}}, w_{\mathrm{feas}}, w_{\mathrm{cons}}$ & .1, .3, .2, .1, .1, .2 & .3, .1, .2, .2, .1, .1 \\
$\omega_{\mathrm{eMBB}}, \omega_{\mathrm{URLLC}}$ & eMBB/URLLC SIs' weight in \eqref{eq:domain-reward} & 3 & 5 \\
\midrule
\multicolumn{4}{c}{\textbf{VNF Characteristics~\cite{tnse25nouruzi}} $\mathrm{Req}^{\mathrm{CN}}_{\boldsymbol{v}_\varsigma}$\;=\;$(p^\varsigma_{\mathrm{tier}},\,p^{v_\varsigma}_{\mathrm{type}},\,\boldsymbol{r}^{v_\varsigma})^{\dagger}$} \\
\midrule
$v_1$ & AMF & 1, 0, 8.5, 0, 16, 60, 1 & 0, 0, 6, 0, 12, 50, 1.5 \\
$v_2$ & UPF & 1, 0, 11, 0, 48, 180, 2.5 & 0, 0, 12, 0, 16, 100, 3 \\
$v_3$ & SMF (eMBB) / NWDAF (URLLC) & 1, 0, 7.2, 0, 12, 90, 0.5 & 0, 1, 0, .00025, 24, 50, 2 \\
\midrule
\multicolumn{4}{c}{\textbf{DRL Hyperparameters}} \\
\midrule
Masked PPO & The algorithm used for TN/CN/E2E agents & \multicolumn{2}{l}{$\alpha=10^{-4}$, $\gamma=0.99$, $\lambda=0.95$, $\epsilon=0.2$, batch size$=512$,} \\
& & \multicolumn{2}{l}{~entropy coeff.$=0.005$, NN arch.$=256{\times}256$} \\
SAC & The algorithm used for RAN agent~\cite{camad24ebrahimi} & \multicolumn{2}{l}{As in~\cite{camad24ebrahimi}} \\
\bottomrule
\multicolumn{4}{@{}p{\textwidth}@{}}{%
\footnotesize $^{\dagger} p^\varsigma_{\mathrm{tier}}$: 0\,=\,edge-only, 1\,=\,any tier; $p^{v_\varsigma}_{\mathrm{type}}:$ 0\,=\,CPU, 1\,=\,GPU. $\boldsymbol{r}^{v_\varsigma}=$(CPU ($10^6 $ cyc/Mbit), GPU (s/Mbit), RAM (GB), Storage (GB), NIC (Gbps)).%
}
\end{tabularx}
}
\end{table*}
%%%%%%
\subsubsection{Training Configuration}
The RAN agent uses SAC~\cite{camad24ebrahimi}; all other agents use Masked PPO~\cite{prep17schulman} with hyperparameters in Table~\ref{tab:simulation-settings}. Training runs offline for 1000~episodes of 100~NRTI steps each; the RAN, TN/CN, and E2E agents act every 1, 5, and 10~NRTIs, yielding 100, 20, and 10 steps/episode, respectively. Per-episode cumulative rewards are divided by these step counts for cross-agent comparison. The 1000-episode budget was selected empirically as the smallest value at which all four agents reach plateau under the three-stage curriculum; halving it left them un-converged, while doubling gave negligible gains. Failures and load bursts are injected during training; at deployment, the trained policies run online with the ILP warm-start re-solved once per operational epoch, and periodic retraining is triggered when traffic or topology statistics drift from the training distribution.
%The RAN agent uses SAC~\cite{camad24ebrahimi}; all other agents use Masked PPO~\cite{prep17schulman} with hyperparameters in Table~\ref{tab:simulation-settings}. Training runs for 1000~episodes of 100~NRTI steps each, with the E2E agent acting every 10~NRTIs (10~LTI steps/episode), TN/CN agents every 5~NRTIs (20~MTI steps/episode), and the RAN agent every NRTI (100~steps/episode). Per-episode cumulative rewards are divided by the respective step counts for cross-agent comparison. %normalized cross-agent comparison.
 
\subsubsection{Evaluation Metrics}
Performance is quantified via the satisfaction metrics of Sec.~\ref{subsec:satisfaction-model}. %For each agent, we report the geometric-mean reward normalized by its per-episode step count. 
At the E2E level, $\tilde{\mathbb{S}}^{\mathrm{(E2E)}}$ (Sec.~\ref{e2e-domain-section}) captures rate, delay, congestion, status, feasibility, and consistency satisfaction across all SIs. Per-SI total E2E delay is examined to reveal slice-type-specific impacts. All evaluation results are averaged over 20~Monte Carlo episodes with fixed failure patterns for reproducibility.

%==============================
\subsection{Stress Testing} \label{subsec:stress}
We evaluate ENSURE's resilience under representative failure and congestion scenarios, focusing on the TN agent. Isolated RAN/CN/E2E stress tests yielded consistent findings and are subsumed by the \emph{multi-agent} evaluation in Sec.~\ref{subsec:ablation}.
% We evaluate ENSURE's resilience under representative failure and congestion scenarios, focusing on the TN agent, which must jointly handle link failures and congestion. Isolated RAN/CN/E2E stress tests yielded consistent findings and are subsumed by the \emph{multi-agent} evaluation in Sec.~\ref{subsec:ablation}, which exercises all four agents jointly under multi-domain stress.
% We evaluate ENSURE's resilience under failure and congestion scenarios representative of real network conditions. As a representative case, we present detailed stress-test results for the TN agent, which must jointly handle link failures and traffic-driven congestion; analogous tests for RAN/CN/E2E agents yielded consistent findings but are omitted for brevity.

% \subsubsection{Isolated TN-Agent Stress Test}
% \label{subsubsec:tn_isolated}
The \emph{TN agent} is instructive because it must balance delay, congestion, and link-status objectives against the topology's limited alternatives, making it well-suited to demonstrate how satisfaction-triggered masking enables targeted adaptation. %Because the mask suppresses updates when the current assignment is acceptable, fewer policy-network forward passes are spent on futile changes, lowering wall-clock cost per episode (see Table~\ref{tab:e2e_curriculum}).
% The \emph{TN agent} is particularly instructive because it must balance multiple competing objectives (delay minimization, congestion avoidance, and link-status awareness) while the underlying topology constrains available alternatives. This makes it well-suited for demonstrating how satisfaction-triggered masking enables targeted adaptation without excessive reconfiguration. Because the mask suppresses action updates when the current assignment remains acceptable, agents spend fewer policy-network forward passes on futile changes, translating into a lower wall-clock cost per episode (see Table~\ref{tab:e2e_curriculum}). %, right-most columns).
%This makes it well-suited for demonstrating how satisfaction-triggered masking enables targeted adaptation without excessive reconfiguration.

\subsubsection{Training Methodology}\label{subsubsec:tn_curriculum}
We employ curriculum learning~\cite{nfvsdn23shahbazi} to progressively expose the agent to harder scenarios, as abrupt exposure diverges early-stage learning. Table~\ref{tab:tn_curriculum} details the three-stage curriculum: a \emph{ramp} (Initial$\to$Intermediate$\to$Final) increases load and failure probability, while periodic \emph{interleaving} of Initial/Final episodes mitigates catastrophic forgetting~\cite{marl24albrecht}. Each episode includes at most one failed link, consistent with the low probability of simultaneous failures in operational networks~\cite{tnse25nouruzi}; multi-link failures are a limitation flagged for future work. Training runs in a high-fidelity simulator; online fine-tuning to close the sim-to-real gap is an orthogonal direction.
% We employ curriculum learning~\cite{nfvsdn23shahbazi} to progressively expose the agent to increasingly challenging scenarios, as abrupt exposure to difficult episodes can destabilize early-stage policy learning. Table~\ref{tab:tn_curriculum} details the three-stage curriculum. The \emph{ramp} progression (Initial$\to$Intermediate$\to$Final) gradually increases traffic load and failure probability, while periodic \emph{interleaving} (i.e., injecting Initial and Final episodes at regular intervals) mitigates catastrophic forgetting by revisiting previously mastered scenarios~\cite{marl24albrecht}. This maintains robustness across the full difficulty spectrum, avoiding overfitting to the latest stage. Each episode includes at most one failed link, consistent with the low probability of simultaneous failures in operational networks~\cite{tnse25nouruzi}.  %This combination maintains %at most one failed link main citation~\cite{imc18meza}

\subsubsection{Evaluation scenario}\label{subsubsec:tn_scenario}
We construct a controlled stress test combining deterministic link failure with monotonically increasing congestion to assess worst-case resilience. Link~$l_2$ fails throughout steps~8--15 while eMBB users ramp linearly; URLLC users remain fixed to isolate congestion effects. Results are averaged over 20 runs with identical failure patterns.
% We construct a controlled stress test that combines deterministic link failure with monotonically increasing congestion to assess worst-case resilience. Link~$l_2$ fails throughout steps~8--15, while eMBB users ramp linearly, inducing progressive congestion. URLLC users remain fixed to isolate congestion effects. Results are averaged over 20 Monte Carlo runs with identical failure patterns for fair comparison.

\subsubsection{Baselines}\label{subsubsec:tn_baselines}
We compare four TN-only variants (the methodology itself is generic across TN/CN/E2E): (i)~\textbf{Proposed} (PPO, with both warm-start and masking), (ii)~\textbf{NoMask PPO} (warm-start retained, action masking removed), (iii)~\textbf{NoWarm PPO} (action masking and warm-start removed), and (iv)~\textbf{Static} (shortest path, no learning).
%We compare: (i)~\textbf{Masked PPO} (proposed), (ii)~\textbf{NoMask PPO} (same architecture without action masking), (iii)~\textbf{NoWarm PPO} (masked PPO without warm-start initialization), and (iv)~\textbf{Static} (always selects the shortest path regardless of network conditions).

\subsubsection{Results and analysis}\label{subsubsec:tn_results}
Fig.~\ref{fig:tn_isolated_eval} presents per-step evaluation metrics. Under nominal conditions (Fig.~\ref{fig:tn_eval_reward}), all schemes perform comparably. Upon link failure at step~8, Static drops sharply and remains degraded because it cannot reroute. In Proposed, a below-pivot $S_{\mathrm{status}}$ on affected SIs lifts the mask ($m_i(t)=1$, Table~\ref{tab:masking-criteria}), triggering a path switch on the next MTI step and restoring satisfaction within two steps; these lifting events coincide with the jumps in Fig.~\ref{fig:tn_eval_paths}, evidencing satisfaction-triggered (rather than blind) adaptation.
%Upon link failure (step~8), Static drops sharply and remains degraded because it cannot reroute, whereas Masked PPO detects the failure via $S_{\mathrm{status}}$ and switches affected SIs to backup paths within two steps. Post-recovery, Masked PPO outperforms all baselines owing to its constraint-driven adaptation. 
% Fig.~\ref{fig:tn_isolated_eval} presents per-step evaluation metrics. In Fig.~\ref{fig:tn_eval_reward}, all schemes initially achieve ${\sim}0.9$ reward under nominal conditions. When the failure occurs (step~8), the Static baseline drops sharply to ${\sim}0.35$ and remains degraded until link recovery, whereas Masked PPO detects the failure via the $S_{\mathrm{status}}$ component and switches affected SIs to backup paths within two steps, maintaining reward above $0.7$. Post-recovery, Masked PPO restores reward to ${\sim}0.8$, outperforming NoWarm PPO (${\sim}0.7$) and Static (${\sim}0.65$).
Fig.~\ref{fig:tn_eval_delay} shows that Static achieves the highest $S_{\mathrm{delay}}$ because shortest paths minimize propagation delay; however, Fig.~\ref{fig:tn_eval_cong} reveals its critical weakness: without load redistribution, $S_{\mathrm{cong}}$ degrades monotonically as eMBB load increases. Proposed partially mitigates this by rerouting one eMBB SI (visible in Fig.~\ref{fig:tn_eval_paths}).
% Fig.~\ref{fig:tn_eval_delay} reveals that Static maintains high $S_{\mathrm{delay}}$ (${\sim}0.92$) because shortest paths inherently minimize propagation delay; however, this metric alone is insufficient for assessing overall performance. Fig.~\ref{fig:tn_eval_cong} exposes Static's critical weakness: as eMBB load increases, $S_{\mathrm{cong}}$ degrades monotonically to ${\sim}0.45$ because Static cannot redistribute traffic. Masked PPO adapts by rerouting one eMBB SI around step~18 (visible in Fig.~\ref{fig:tn_eval_paths}), thereby slowing the decline of $S_{\mathrm{cong}}$ relative to Static.

Fig.~\ref{fig:tn_eval_paths} confirms that Proposed achieves the fewest path changes, permitting switches only upon satisfaction violations (failure at steps~8--10; congestion at steps~16--18). NoMask oscillates throughout, while NoWarm switches prematurely due to suboptimal initialization. Their wide shaded bands reflect the well-known high run-to-run variance of unconstrained policy-gradient DRL on combinatorial tasks~\cite{marl24albrecht}; masking (bounding the feasible set) and warm-starting (removing the cold-start seed dependence) jointly reduce this variance by over an order of magnitude (Table~\ref{tab:tn_curriculum}, $\pm$ columns), alongside Proposed's highest evaluation reward. Static's strong Initial-stage performance does not transfer to failure scenarios.
% NoMask oscillates throughout, and NoWarm switches prematurely due to suboptimal  initialization, both operationally undesirable as unnecessary switching increases control-plane overhead~\cite{tnse25nouruzi}. 
% Fig.~\ref{fig:tn_eval_paths} demonstrates that Masked PPO achieves superior adaptation with minimal reconfiguration overhead. The masking mechanism permits path changes only when satisfaction thresholds are violated: during failure detection (steps~9--10) and congestion onset (steps~16--18). This constraint-driven approach guarantees that necessary switches occur precisely when needed. In contrast, NoMask and NoWarm PPO rely entirely on learned policy quality; despite curriculum training, both exhibit undesirable behavior. NoWarm PPO switches paths in early steps (before any failure) due to suboptimal initialization, while NoMask PPO oscillates throughout the failure window. This instability is operationally significant, as unnecessary path switching increases control-plane overhead and risks transient disruptions~\cite{tnse25nouruzi}.
%Table~\ref{tab:tn_curriculum} summarizes stage-wise and evaluation rewards. Masked PPO achieves the highest mean evaluation reward ($0.79{\pm}0.01$) with the lowest variance, demonstrating consistent performance across the combined failure-congestion scenario. The Static baseline's poor evaluation reward ($0.63$) despite strong Initial-stage performance ($0.94$) underscores the necessity of adaptive policies for deployments where failures and load variations are inevitable.

%%%
\begin{table*}[t]
\centering
\caption{Isolated TN-agent curriculum (training) and evaluation.} % Training uses ramp progression with global interleaving~\cite{nfvsdn23shahbazi}.}
%\caption{Isolated TN-agent curriculum (training) and evaluation. Training uses ramp progression with global interleaving: every 10th episode reverts to Initial; every $(10k{+}8)$th advances to Final \cite{nfvsdn23shahbazi}. At most one link fails per episode.}%interleaving main citation \cite{wsc24mobius}
\label{tab:tn_curriculum}
\scriptsize
\setlength{\tabcolsep}{2.2pt}
\begin{tabular}{@{}l c c c c c c c c | c c c c@{}}
\toprule
& \multicolumn{8}{c|}{\textbf{Scenario Configuration}} & \multicolumn{4}{c}{\textbf{Mean Reward $\pm$ Std}} \\
\cmidrule(lr){2-9} \cmidrule(lr){10-13}
\textbf{Stage} & \textbf{Episodes} & \textbf{eMBB} & \textbf{URLLC} & \textbf{Load} & $\boldsymbol{p}_{\textbf{fail}}$ & \textbf{Failed} & \textbf{Start} & \textbf{Duration} & \multirow{2}{*}{\textbf{Proposed}} & \textbf{NoMask} & \textbf{NoWarm} & \multirow{2}{*}{\textbf{Static}} \\
& (range; count) & (users/SI) & (users/SI) & Pattern & & Links & Range$^*$ & (steps)$^\dagger$ &  & \textbf{PPO} & \textbf{PPO} & \\
\midrule
Initial  & 0--15\%; ${\sim}$1.5k & 10--35   & 5--20  & Flat           & 0.00 & n/a              & n/a  & n/a        & $\mathbf{0.94} \pm$0.01 & 0.90$\pm$0.04 & 0.87$\pm$0.05 & $\mathbf{0.94} \pm$0.00 \\
Intermediate   & 15--60\%; ${\sim}$4.5k & 20--170  & 10--40 & Triangle  & 0.25 & \{2,3,5,8,14\}   & 25--80\%  & $\sim$6   & $\mathbf{0.88} \pm$0.03 & 0.86$\pm$0.04 & 0.84$\pm$0.05 & 0.66$\pm$0.00 \\
Final  & 60--100\%; ${\sim}$4k & 50--260  & 20--70 & Random+burst & 0.80 & All$^\ddagger$    & 5--98\%  & $\sim$20     & $\mathbf{0.84} \pm$0.06 & 0.82$\pm$0.07 & 0.78$\pm$0.08 & 0.56$\pm$0.00 \\
\midrule
\textbf{Evaluation}        & 20 runs & 50$\to$250 & 50 (fixed)   & Linear ramp    & 1.00 & \{$l_2$\}        & 20--40\%$^\S$  & 8     & $\mathbf{0.79} \pm$0.01 & 0.76$\pm$0.02 & 0.69$\pm$0.05 & 0.63$\pm$ 0.00 \\
\bottomrule
\multicolumn{13}{@{}l@{}}{\scriptsize $^*$Allowed range for stochastic start step. $^\dagger$Mean of Poisson-sampled duration. $^\ddagger$All links except $l_0$ (LA1--LA2). $^\S$Deterministic (steps 8--15).}
\end{tabular}
\end{table*}
%%%%%%%%%%%
% \begin{figure*}[t]
% \centering
% \subfloat[Reward\label{fig:tn_eval_reward}]{\includegraphics[width=0.245\textwidth]{figures/tn_eval_reward.pdf}}
% \hfill
% \subfloat[$S_{\mathrm{delay}}$ (aggregated)\label{fig:tn_eval_delay}]{\includegraphics[width=0.245\textwidth]{figures/tn_delay_satisfaction_agg.pdf}}
% \hfill
% \subfloat[$S_{\mathrm{cong}}$ (aggregated)\label{fig:tn_eval_cong}]{\includegraphics[width=0.245\textwidth]{figures/tn_congestion_satisfaction_agg.pdf}}
% \hfill
% \subfloat[Cumulative path changes\label{fig:tn_eval_paths}]{\includegraphics[width=0.245\textwidth]{figures/tn_total_path_changes.pdf}}
% \caption{Isolated TN-agent evaluation under joint link failure (steps 8--16) and increasing eMBB load (50$\to$250 users/SI). The color coded shaded regions indicate variability across the Monte Carlo runs for each algorithm.}
% \label{fig:tn_isolated_eval}
% \end{figure*}
%%%%%%%%%%%%%%%%%%%%
\begin{figure*}[t]
\centering
\subfloat[Reward\label{fig:tn_eval_reward}]{\includegraphics[width=0.485\textwidth]{figures/tn_eval_reward.pdf}}
\hfill
\subfloat[$S_{\mathrm{delay}}$ (aggregated)\label{fig:tn_eval_delay}]{\includegraphics[width=0.485\textwidth]{figures/tn_delay_satisfaction_agg.pdf}}
\hfill
\subfloat[$S_{\mathrm{cong}}$ (aggregated)\label{fig:tn_eval_cong}]{\includegraphics[width=0.485\textwidth]{figures/tn_congestion_satisfaction_agg.pdf}}
\hfill
\subfloat[Cumulative path changes\label{fig:tn_eval_paths}]{\includegraphics[width=0.485\textwidth]{figures/tn_total_path_changes.pdf}}
\caption{Isolated TN-agent evaluation under link failure and eMBB load ramp. Shaded regions: Monte Carlo variability.}
\label{fig:tn_isolated_eval}
\end{figure*}
%===============================================
\subsection{Ablation Study: E2E Multi-Agent Evaluation}
\label{subsec:ablation}

Having validated the TN agent's masking and warm-start mechanisms in isolation, we now evaluate the complete ENSURE framework in multi-agent E2E settings. This ablation study quantifies the individual contribution of each component, i.e., action masking, ILP warm-starting, and adaptive multi-domain DRL, by comparing ENSURE against three degraded variants under progressively challenging scenarios.

\subsubsection{Baselines}
All four schemes share the same SAC-based RAN agent~\cite{camad24ebrahimi}, differing only in TN/CN/E2E configurations:
\begin{itemize}[leftmargin=*]
\item \textbf{ENSURE}: Proposed framework with satisfaction-triggered masked PPO for TN/CN/E2E agents, ILP warm-start for E2E and CN at episode onset, and shortest-path initialization for TN (Sec.~\ref{subsec:training}).
\item \textbf{NoMask}: Standard PPO (without action masking) for TN/CN/E2E; all other components identical to ENSURE, including warm-start.
\item \textbf{NoWarm}: Standard (unmasked) PPO with random initial assignments; masking is disabled, as without a feasible warm-start it would pin agents to an infeasible starting point. % (cf. Sec.~\ref{subsec:training}).
% \item \textbf{NoWarm}: Masked PPO for TN/CN/E2E agents, but with random initial assignments instead of ILP warm-start; agents learn from scratch.
\item \textbf{Static}~\cite{camad24ebrahimi}: ILP-only initialization for CN/E2E and shortest-path TN; no intra-episode adaptation.
% \item \textbf{Static}~\cite{camad24ebrahimi}: ILP-only initialization for CN/E2E at episode start (no subsequent adaptation) and shortest-path heuristic for TN. This scheme represents a RAN-only DRL baseline where TN and CN are statically provisioned and do not respond to intra-episode dynamics; functionally equivalent to existing RAN-centric slicing approaches that assume fixed TN/CN infrastructure~\cite{camad24ebrahimi}.
% TODO: Add a DDPG-based Static variant (Static-DDPG) following \cite{tvt23choi} as an additional RAN-only baseline.
% The Static-DDPG variant replaces SAC with DDPG for the RAN agent to benchmark against \cite{tvt23choi}.
% Expected to perform at or below the current SAC-based Static since SAC generally outperforms DDPG in continuous control.
\end{itemize}
% We deliberately benchmark against (i)~Static, which captures the dominant family of RAN slicing approaches that assume fixed TN/CN provisioning~\cite{camad24ebrahimi,jiot24filali,jsac24zangooei}, and (ii)~two ablations of ENSURE itself (NoMask, NoWarm) that isolate the contribution of each novel mechanism. A direct head-to-head comparison with a published cooperative MA-DRL E2E framework was not possible: existing MA-DRL slicing works~\cite{access21kim,tnsm22sulaiman,dyspan26mohammadi,srep25mao} differ in (a) their domain coverage (RAN-only or RAN+CN), (b) their reward and observation/action spaces (e.g., scalarized utility~\cite{access21kim}, Lagrangian-constrained PPO~\cite{dyspan26mohammadi}, evolutionary game theory~\cite{srep26mao}), and (c) their underlying network model (single-cell, edge-only, or non-Open RAN). Re-implementing any of them on our three-domain Open RAN topology with our SI abstraction and three-timescale training loop would require substantial reformulation that risks misrepresenting the original method. The Static baseline therefore plays the role of a representative RAN-centric proxy~\cite{camad24ebrahimi,jiot24filali}, while ablations isolate the marginal contribution of masking, warm-starting, and timescale separation. We acknowledge this as a limitation and identify a fair, common-testbed evaluation against published MA-DRL E2E baselines as a priority for follow-up work.
The \textbf{Static} baseline represents \textit{RAN-only slicing} approaches~\cite{tvt23choi,camad24ebrahimi,jiot24filali}: an adaptive SAC RAN agent over TN/CN/E2E frozen at their warm-start, mirroring the fixed-TN/CN provisioning assumption of those works. An accurate comparison with the discussed works on \emph{cooperative} MA-DRL E2E slicing (cf. Sec.~\ref{related-work-section}) was not feasible because these works differ in domain coverage, rewards, and network model. Therefore, adapting them to our three-domain Open RAN topology with the SI abstraction and three-timescale loop would demand substantive reformulation and risk misrepresentation.
%Common-testbed evaluation is prioritized for follow-up work.

\subsubsection{Training Methodology}\label{subsubsec:e2e_curriculum}
Following the curriculum strategy validated in Sec.~\ref{subsec:stress}, all four methods undergo similar three-stage training (Table~\ref{tab:e2e_curriculum}: stage configurations, per-method per-agent normalized rewards, and per-episode wall-clock time). Training uses sequential progression over 1000~episodes across three stages of increasing difficulty (\emph{Initial} / \emph{Intermediate} / \emph{Final}: 200/350/450 episodes), with global interleaving of Initial- and Final-stage episodes throughout to mitigate catastrophic forgetting (cf. Table~\ref{tab:tn_curriculum} for the TN-isolated analogue). 
% Training uses sequential stage progression over 1000~episodes with a 20\%/35\%/45\% Initial/Intermediate/Final split and global interleaving (cf.\ Table~\ref{tab:tn_curriculum}). 
Each stage scales up active user fractions, user mobility, TN-link/CN-server failure probabilities and durations, and injected background TN and CN load. User arrivals and departures follow Poisson processes with stage-dependent rates.
% Each stage progressively increases active user fractions, user mobility, and failure probabilities for both TN~links and CN~servers, with growing failure durations; background TN congestion and CN server load are injected with increasing severity across stages. User arrivals and departures follow Poisson processes with stage-dependent rates. %($\lambda_{\mathrm{arr}}/\lambda_{\mathrm{dep}}$: 0.02/0.01 $\to$ 0.03/0.02 $\to$ 0.05/0.03).
%with failure durations growing from 2--5 to 7--19 steps (TN) and 3--8 to 10--19 steps (CN).

\subsubsection{Training Convergence}\label{subsubsec:convergence}
Fig.~\ref{fig:e2e_convergence} reports per-agent normalized training rewards, each agent's cumulative per-episode reward divided by its step count (RAN: 100, TN/CN: 20, E2E: 10) for comparability. The panels are read together with each scheme's designed capabilities (cf. Baselines above).
%We interpret them in light of what each scheme's TN/CN/E2E block can actually do — adaptive with both masking and warm-start in ENSURE, masking removed in NoMask and both removed in NoWarm, and non-adaptive in Static.

In the RAN panel, ENSURE and Static perform comparably (same SAC agent); ENSURE's marginal advantage in later stages reflects better cross-domain conditions from its adaptive TN/CN agents. NoMask and NoWarm underperform as their erratic TN/CN behavior inflates E2E delay, indirectly degrading RAN satisfaction. Reward drops at episodes~${\sim}$200, ${\sim}$550 confirm the curriculum's increasing difficulty.
% In the RAN panel, ENSURE and Static perform comparably, as both employ the same SAC agent; ENSURE's marginal advantage in later stages reflects improved cross-domain conditions created by its adaptive TN/CN agents. NoMask and NoWarm underperform as their erratic TN/CN behavior increases E2E delay, indirectly degrading RAN-level satisfaction. The reward drop at stage transitions (episodes~${\sim}$200 and~${\sim}$550) confirms that the curriculum progressively increases difficulty.

% To aid interpretation, recall that the four schemes differ only in their TN/CN/E2E blocks: \textbf{Static} re-solves the ILP \emph{once per episode} (at onset) and then holds the assignment fixed — it is optimized at the episode level, not truly stateless; \textbf{NoMask} and \textbf{NoWarm} use adaptive PPO but remove masking and warm-start, respectively; \textbf{ENSURE} retains both. 
In the TN panel, ENSURE achieves the highest and most converged reward, owing to the joint effect of masking and warm-start. Static matches ENSURE under benign Initial conditions but oscillates under failures/congestion. NoMask and NoWarm fail to converge due to \emph{sample inefficiency}~\cite{marl24albrecht}: the training budget needed to reach a given policy quality grows sharply with the effective action-space size (unbounded under NoMask) and worsens under random initialization (NoWarm). Prolonging training 5--10$\times$ is both prohibitive — the current 1000-episode run already costs 233 and 271 minutes of aggregate wall-clock time for NoMask and NoWarm, respectively, on an RTX 3080 GPU (see Table~\ref{tab:e2e_curriculum}) — and unnecessary, as ENSURE converges within the same budget. CN and E2E panels follow the same pattern, with the ENSURE gap widening as failures demand coordinated adaptation; the E2E panel separates most sharply because $R^{(\mathrm{E2E})}$ aggregates cross-domain satisfactions via geometric mean, amplifying any domain's deficiencies. Per-episode wall-clock (Table~\ref{tab:e2e_curriculum}) ranks Static $<$ ENSURE $<$ NoMask $<$ NoWarm: Static is cheapest because TN/CN/E2E do not update intra-episode; ENSURE pays a modest overhead for masked inference that stays below NoMask and NoWarm, whose unmasked policies perform more forward passes on futile changes, which is the same mechanism that drives their oscillatory behavior in Fig.~\ref{fig:combined_time}.
% NoMask and NoWarm fail to converge due to \emph{sample inefficiency}: the unconstrained action space (NoMask) and random initialization (NoWarm) drastically increase required training episodes. 
% In the RAN panel, ENSURE and Static perform comparably within each stage, with marginal superiority for ENSURE in the Intermediate and Final stages. This is expected: both employ the same SAC RAN agent, but the adaptive TN/CN agents in ENSURE create more favorable cross-domain conditions (e.g., lower TN delays via timely path switching) that indirectly improve RAN-level $\mathbb{S}^{\mathrm{(RAN)}}_{\mathrm{delay}}$. NoMask and NoWarm consistently underperform ENSURE due to their unreliable TN/CN behavior, whose erratic path and VNF placement decisions increase the E2E delay observed by the RAN agent, thereby degrading its aggregated satisfaction. The reward drop at stage transitions (episodes~${\sim}$200 and~${\sim}$550) is visible for all methods, confirming that the curriculum effectively increases scenario difficulty.
% The TN panel reveals that ENSURE achieves the highest and most stable reward across Intermediate and Final stages, owing to the combination of masking (which prevents unnecessary path switches) and warm-start (which provides high-quality initial assignments). Although Static performs similarly to ENSURE in the Initial stage, where its heuristic coincides with the optimal policy under benign conditions, it exhibits oscillations during Intermediate and Final stages as failures and congestion emerge, rendering its behavior unreliable. NoMask and NoWarm fail to converge to acceptable reward levels at any stage. This is attributable to \emph{sample inefficiency}: without masking, the effective action space remains unconstrained, and without warm-start, the agent must explore from random initializations. Both factors drastically increase the number of episodes needed for policy improvement. In principle, substantially longer training (e.g., $5$--$10\times$ more episodes per stage) might partially close this gap; however, such an increase is computationally prohibitive as the current 1000-episode training for NoMask and NoWarm methods already requires 233 and 271 minutes of wall-clock time, respectively, on an RTX~3080 GPU, and unnecessary given that ENSURE demonstrates that masking and warm-start suffice to achieve stable convergence within practical training budgets.
% Similar patterns are observed for the CN and E2E agents. ENSURE consistently achieves the highest normalized rewards, and the gap relative to other baselines widens in the Intermediate and Final stages, where frequent failures and elevated load demand coordinated cross-domain adaptation. The CN agent's reward for NoMask/NoWarm suffers from the same sample inefficiency issue explained above. Static's CN reward is stable across Initial and Intermediate, since the ILP warm-start provides a reasonable initial VNF placement, but it cannot adapt when background load shifts or servers fail mid-episode, resulting in a growing gap relative to ENSURE as stage difficulty increases. %The CN agent's reward for NoMask/NoWarm suffers from the same sample inefficiency, compounded by the larger per-SI action space ($\prod_{i} M^{V_i}$ server selections) that amplifies the exploration burden. 
% The E2E agent panel exhibits the most pronounced separation across methods. Because $R^{(\mathrm{E2E})}$ aggregates cross-domain satisfactions via the geometric mean (Sec.~\ref{e2e-problem-formulation}), deficiencies in \emph{any} domain are amplified. Consequently, the TN/CN instabilities of NoMask and NoWarm propagate to very low E2E rewards, while ENSURE maintains satisfaction above the $0.5$ threshold even under the most challenging conditions, confirming that joint masking and warm-start across all agents yield compounding benefits at the E2E level.

\begin{figure}[t]
\centering
\includegraphics[width=1.02\columnwidth]{figures/02b_comp_per_agent_reward_normalized.pdf}
\caption{Per-agent normalized reward during E2E multi-agent training (rolling average, window$=$10).} 
%The logits shown are raw outputs from the actor network before normalization; the softmax function converts these into a valid probability distribution, assigning zero probability to masked actions via \infty logits.
% Each agent's cumulative episode reward is divided by its per-episode step count (RAN: 100, TN/CN: 20, E2E: 10) for comparability.}% Vertical dashed lines indicate curriculum stage transitions.}
\label{fig:e2e_convergence}
\end{figure}


%%%% E2E CURRICULUM SUMMARY TABLE — aligned and presentation-safe %%%%
\begin{table*}[t]
\centering
\caption{E2E multi-agent curriculum (training) and evaluation scenario.} % Training uses sequential stage progression (Initial$\to$Intermediate$\to$Final) over 1000 episodes with interleaving (cf.\ Table~\ref{tab:tn_curriculum}).}
\label{tab:e2e_curriculum}
\resizebox{\textwidth}{!}{%
\footnotesize
\setlength{\tabcolsep}{1.28pt}
\renewcommand{\arraystretch}{1.45}
\begin{tabular}{
  @{}l c c c c c c
  r >{\columncolor{ltgray}}r r >{\columncolor{mdgray}}c
  r >{\columncolor{ltgray}}r r >{\columncolor{mdgray}}c
  r >{\columncolor{ltgray}}r r >{\columncolor{mdgray}}c
  r >{\columncolor{ltgray}}r r >{\columncolor{mdgray}}c
  @{\hspace{3pt}}*{4}{>{\centering\arraybackslash}m{6.2mm}}
  @{}
}
\toprule
& \multicolumn{6}{c}{\textbf{Scenario Configuration}}
& \multicolumn{16}{c}{\textbf{Normalized Mean Reward per Agent}}
& \multicolumn{4}{c}{\textbf{Time$^{\P}$ (s/ep)}} \\
\cmidrule(lr){2-7} \cmidrule(lr){8-23} \cmidrule(lr){24-27}

& & & & & &
& \multicolumn{4}{c}{\textbf{ENSURE}}
& \multicolumn{4}{c}{\textbf{NoMask}}
& \multicolumn{4}{c}{\textbf{NoWarm}}
& \multicolumn{4}{c}{\textbf{Static}}
& \multicolumn{1}{c}{}
& \multicolumn{1}{c}{}
& \multicolumn{1}{c}{}
& \multicolumn{1}{c}{} \\
\cmidrule(lr){8-11} \cmidrule(lr){12-15} \cmidrule(lr){16-19} \cmidrule(lr){20-23}

\textbf{Stage}
& \textbf{Eps.}
& \textbf{\shortstack{Active\\Users}}
& \textbf{\shortstack{Vel.\\(m/s)}}
& $\boldsymbol{p}^{\textbf{TN}}_{\textbf{f}}$/$\boldsymbol{p}^{\textbf{CN}}_{\textbf{f}}$
& \textbf{\shortstack{Fail dur.$^{*}$\\TN/CN}}
& \textbf{\shortstack{BG$^{\diamond}$\\TN/CN}}
& \textbf{RAN} & \cellcolor{ltgray}\textbf{TN} & \textbf{CN} & \cellcolor{mdgray}\textbf{E2E}
& \textbf{RAN} & \cellcolor{ltgray}\textbf{TN} & \textbf{CN} & \cellcolor{mdgray}\textbf{E2E}
& \textbf{RAN} & \cellcolor{ltgray}\textbf{TN} & \textbf{CN} & \cellcolor{mdgray}\textbf{E2E}
& \textbf{RAN} & \cellcolor{ltgray}\textbf{TN} & \textbf{CN} & \cellcolor{mdgray}\textbf{E2E}
& \multicolumn{1}{c}{\rothead{ENSURE}}
& \multicolumn{1}{c}{\rothead{NoMask}}
& \multicolumn{1}{c}{\rothead{No-Warm}}
& \multicolumn{1}{c}{\rothead{Static}} \\
\midrule

Initial
  & 0-199 & 10--30\% & 0--5 & .45/.45 & 2--5/3--8 & Lo/Lo
  & .90 & \cellcolor{ltgray}.90 & .57 & \cellcolor{mdgray}\textbf{.69}
  & .73 & \cellcolor{ltgray}.34 & .10 & \cellcolor{mdgray}\textbf{.13}
  & .71 & \cellcolor{ltgray}.28 & .14 & \cellcolor{mdgray}\textbf{.03}
  & .90 & \cellcolor{ltgray}.90 & .63 & \cellcolor{mdgray}\textbf{.70}
  & 6.6 & 6.9 & 8.0 & 3.7 \\
Intermediate
  & 200-549 & 30--60\% & 4--10 & .70/.70 & 8--16/10--18 & M/M
  & .72 & \cellcolor{ltgray}.86 & .65 & \cellcolor{mdgray}\textbf{.66}
  & .54 & \cellcolor{ltgray}.28 & .50 & \cellcolor{mdgray}\textbf{.22}
  & .65 & \cellcolor{ltgray}.36 & .45 & \cellcolor{mdgray}\textbf{.23}
  & .70 & \cellcolor{ltgray}.79 & .54 & \cellcolor{mdgray}\textbf{.59}
  & 11.3 & 11.9 & 13.8 & 3.9 \\
Final
  & 550-999 & 60--100\% & 9--20 & 1/1 & 7--19/10--19 & Hi/Hi
  & .53 & \cellcolor{ltgray}.67 & .58 & \cellcolor{mdgray}\textbf{.52}
  & .46 & \cellcolor{ltgray}.23 & .25 & \cellcolor{mdgray}\textbf{.19}
  & .48 & \cellcolor{ltgray}.38 & .45 & \cellcolor{mdgray}\textbf{.16}
  & .52 & \cellcolor{ltgray}.58 & .37 & \cellcolor{mdgray}\textbf{.40}
  & 17.6 & 18.7 & 21.9 & 9.7 \\
\midrule

% \texttt{CN\_S}$^{\dagger}$
%   & 20 & 85--100\% & 10--40 & ---/\checkmark & ---/20--60\% & ---/Hi
%   & .64 & \cellcolor{ltgray}.90 & .91 & \cellcolor{mdgray}\textbf{.70}
%   & .50 & \cellcolor{ltgray}.34 & .10 & \cellcolor{mdgray}\textbf{.13}
%   & .48 & \cellcolor{ltgray}.21 & .03 & \cellcolor{mdgray}\textbf{.10}
%   & .64 & \cellcolor{ltgray}.90 & .59 & \cellcolor{mdgray}\textbf{.57}
%   & 6.6 & 6.5 & 6.6 & 6.1 \\
\texttt{All\_S}$^{\dagger}$
  & 20 & 90--100\% & 15--50 & \checkmark/\checkmark & 15--55/20--65\% & Hi/Hi
  & .61 & \cellcolor{ltgray}.81 & .90 & \cellcolor{mdgray}\textbf{.66}
  & .47 & \cellcolor{ltgray}.27 & .10 & \cellcolor{mdgray}\textbf{.10}
  & .45 & \cellcolor{ltgray}.17 & .03 & \cellcolor{mdgray}\textbf{.07}
  & .57 & \cellcolor{ltgray}.64 & .48 & \cellcolor{mdgray}\textbf{.41}
  & 6.8 & 6.7 & 6.7 & 6.4 \\
\bottomrule

\multicolumn{27}{@{}p{1.32\textwidth}@{}}{%
\small
$^{\natural}$Normalized reward$=$cumulative episode reward$/ $agent steps; differs from per-metric satisfaction $S_{\mathrm{E2E}}$ in Figs.~\ref{fig:combined_bars}--\ref{fig:combined_delay}, which averages satisfaction components across steps and SIs. 
$^{*}$Training: stochastic failure durations in steps; eval: deterministic windows as \% of episode length.
$^{\diamond}$Background (BG) stress level: Lo\,=\,low, M\,=\,medium, Hi\,=\,high injection of additional TN link utilization\,/\,CN server load.
$^{\P}$Mean wall-clock time per episode (seconds).
$^{\dagger}$\texttt{All\_Stress} evaluation scenario: TN links $\{l_1,l_2\}$ fail (15--55\%); CN failure at STTL (20--65\%); BG TN util.\ on $\{l_1,l_2\}$ (0.30--0.45 extra); BG CN load on DNVR \& STTL (0.28--0.45 extra).%
}
\end{tabular}%
}
\end{table*}
%%%%%%%%%%%%%%%%%%%%%%

\subsubsection{Evaluation Under Stress Scenarios}\label{subsubsec:e2e_eval}
We evaluate the trained agents on \texttt{All\_Stress}, a multi-domain scenario combining TN link failures (bottleneck links $\{l_1, l_2\}$), a CN server failure (STTL site), TN/CN background load, and high user mobility (Table~\ref{tab:e2e_curriculum}). By stressing all three domains concurrently, it is the most demanding test of ENSURE's cross-domain coordination and subsumes single-domain failure modes (e.g., CN-only disruptions) omitted here for brevity.
% We evaluate the trained agents on a comprehensive multi-domain stress scenario, \texttt{All\_Stress}, that simultaneously exercises TN, CN, and RAN failure/congestion modes. This scenario combines TN link failures (bottleneck links $\{l_1,l_2\}$), a CN server failure (STTL site), TN/CN background load, and high user mobility (Table~\ref{tab:e2e_curriculum}). By stressing all three domains concurrently, \texttt{All\_Stress} is the most demanding and representative test of ENSURE's cross-domain coordination capabilities; it subsumes single-domain failure modes (e.g., CN-only disruptions) that were also evaluated during development but are omitted here for brevity. $S_{\mathrm{rate}}$ of the E2E agent is excluded from the time-series analysis, as rate satisfaction is dominated by the RAN agent and varies minimally across methods.

%To comply with the 13-page initial submission limit, we prioritized the most representative evaluation scenario in the main manuscript. Additional scenario results have been omitted from the main paper for space reasons and can be incorporated in a revision if the reviewers or editor consider them necessary.
% \subsubsubsection{\texttt{CN\_Stress} scenario}

% Fig.~\ref{fig:cn_stress_bars} presents the per-metric satisfaction metrics of the E2E agent averaged across all LTI steps and SIs over 20~evaluation runs. ENSURE achieves the highest overall $S_{\mathrm{E2E}}=0.77$, followed by Static ($0.64$), NoMask ($0.28$), and NoWarm ($0.27$). The $S_{\mathrm{rate}}$ and $S_{\mathrm{delay}}$ results immediately reveal that the performance separation is primarily driven by TN/CN behavior rather than RAN: ENSURE and Static, both of which employ stable TN path selection (via masking and heuristic, respectively), achieve high delay satisfaction ($S_{\mathrm{delay}}>0.9$), whereas NoMask and NoWarm suffer from erratic path changes that inflate E2E delay, yielding $S_{\mathrm{delay}}<0.65$. The most significant differentiation between ENSURE and Static appears in $S_{\mathrm{cong}}$ and $S_{\mathrm{cons}}$: Static's inability to adaptively migrate VNFs under increasing background load leads to congestion buildup ($S_{\mathrm{cong}}\approx0.45$ vs.\ ENSURE's ${\sim}0.95$), which in turn degrades stability.

% Fig.~\ref{fig:cn_stress_time} shows the temporal evolution of E2E-level satisfaction metrics, aggregated across SIs via the geometric mean. The $S_{\mathrm{delay}}$ subplot confirms that ENSURE and Static maintain high delay satisfaction throughout the episode, as no TN link failures are present in this scenario; the TN domain (the principal contributor to E2E delay) remains unperturbed. By contrast, NoMask and NoWarm exhibit persistently low $S_{\mathrm{delay}}$ (${\sim}$0.6) due to their unnecessary and counterproductive TN path switches, consistent with the instability observed in the isolated TN stress test (Sec.~\ref{subsec:stress}). The $S_{\mathrm{cong}}$ subplot reveals the key distinction between ENSURE and Static. As CN background load intensifies, Static's congestion satisfaction drops to ${\sim}$0.4 because the one-shot ILP placement cannot redistribute VNF load across servers mid-episode. ENSURE maintains $S_{\mathrm{cong}}$ above 0.8 through adaptive VNF migration. Interestingly, NoMask and NoWarm partially recover $S_{\mathrm{cong}}$ toward the end (${\sim}$0.85), suggesting that their CN policies can eventually respond to congestion; however, this recovery is insufficient to offset the sustained $S_{\mathrm{delay}}$ penalty caused by their unstable TN behavior, resulting in low overall $S_{\mathrm{E2E}}$.

% The $S_{\mathrm{status}}$ metric is similar across all methods: the scheduled CN server failure at STTL (site~2, server~1) does not affect the specific servers selected by the ILP warm-start, rendering the failure event largely ineffective in this scenario. Consequently, $S_{\mathrm{feas}}$ is equally high for both ENSURE and Static. However, NoMask and NoWarm exhibit $S_{\mathrm{feas}}<0.4$, as their frequent VNF migrations often produce placements that violate server-type (CPU/GPU) or tier (edge/cloud) constraints, a direct consequence of unconstrained exploration in NoMask and poor initialization in NoWarm. The $S_{\mathrm{cons}}$ behavior mirrors $S_{\mathrm{cong}}$, since $S_{\mathrm{status}}$ is constant over the LTI steps. The $S_{\mathrm{E2E}}$ subplot, confirms ENSURE's consistent superiority across the entire episode. Note that the E2E agent aggregates metrics from domain agents at the \emph{beginning} of each LTI, which introduces a one-step initialization transient visible in the first LTI of all $S_{\mathrm{E2E}}$ curves.

% Fig.~\ref{fig:cn_stress_delay} presents the per-SI total E2E delay (RAN$+$TN$+$CN), grouped by slice type. For URLLC SIs (particularly SI3), NoMask and NoWarm produce mean delays significantly exceeding the 2~ms URLLC budget, with large variance across runs, confirming their unsuitability for delay-sensitive services. ENSURE maintains delay well within the URLLC budget across all SIs. For eMBB SIs, ENSURE and Static both consistently select shortest TN paths (since no TN failures occur), yielding comparable and low delays. However, NoMask and NoWarm exhibit delay fluctuations that occasionally violate the 20~ms eMBB budget, further illustrating how unconstrained TN path changes can degrade performance even without external failures.

% \begin{figure}[t]
% \centering
% \includegraphics[width=0.97\columnwidth]{figures/01_eval_satisfaction_bars-CN_Stress.pdf}
% \caption{Per-metric satisfaction comparison under \texttt{CN\_Stress}.}% (averaged over 20 runs). Error bars indicate $\pm$1 standard deviation.}
% \label{fig:cn_stress_bars}
% \end{figure}

% \begin{figure}[t]
% \centering
% \includegraphics[width=1.02\columnwidth]{figures/02_eval_satisfaction_time-CN_Stress.pdf}
% \caption{E2E satisfaction metrics over LTI steps under \texttt{CN\_Stress}: $S_{\mathrm{delay}}$, $S_{\mathrm{cong}}$, $S_{\mathrm{status}}$, $S_{\mathrm{feas}}$, $S_{\mathrm{cons}}$, and $S_{\mathrm{E2E}}$ (geometric mean across SIs).}
% \label{fig:cn_stress_time}
% \end{figure}

% \begin{figure}[t]
% \centering
% \includegraphics[width=0.9\columnwidth]{figures/03c_eval_delay_total_per_si_grouped_split-CN_Stress.pdf}
% \caption{Per-SI E2E delay under \texttt{CN\_Stress}, grouped by slice type.}
% \label{fig:cn_stress_delay}
% \end{figure}

% \subsubsubsection{\texttt{All\_Stress} scenario}
Fig.~\ref{fig:combined_bars} presents per-metric satisfaction averaged over all LTI steps, SIs, and 20 runs. ENSURE reaches $S_{\mathrm{E2E}}=0.73$, against Static ($0.53$), NoMask ($0.24$), and NoWarm ($0.22$). The Static-vs-ablation gap is itself informative: a feasible-but-frozen ILP solution outperforms adaptive DRL agents lacking the initialization or action-space pruning to converge within the training budget. This is consistent with prior observations that policy-gradient DRL on combinatorial network management is dominated by well-initialized deterministic baselines under limited training~\cite{marl24albrecht}; ENSURE's contribution is to make adaptive DRL \emph{viable} here via warm-start plus masking. Two metrics drive the separation: $S_{\mathrm{cong}}$, where ENSURE's coordinated path switching and VNF migration sustain near-unity congestion satisfaction whereas Static cannot reroute around failed/congested links; and $S_{\mathrm{cons}}$, where Static's persistent congestion cascades into oscillatory behavior. ENSURE's $S_{\mathrm{delay}}$ also stays high via masking-guided path selection, while NoMask/NoWarm inflate delay through unwarranted reconfigurations. This is not due to geometric-mean aggregation (shared across schemes), but because unmasked policies treat path changes as low-cost exploration and continue toggling at evaluation. Their below-threshold $S_{\mathrm{feas}}$ confirms systematic placement-constraint violations from unconstrained exploration and poor initialization.
% Fig.~\ref{fig:combined_bars} presents the per-metric satisfaction averaged across all LTI steps and SIs (20 runs). ENSURE achieves $S_{\mathrm{E2E}}=0.73$, outperforming Static, NoMask, and NoWarm. The gap between Static ($0.53$) and NoMask/NoWarm ($0.24$/$0.22$) is itself informative: a feasible-but-frozen ILP solution outperforms an adaptive RL agent that has neither the initialization nor the action-space pruning to converge within the training budget. This is consistent with prior observations that policy-gradient RL on combinatorial network-management tasks is dominated by well-initialized deterministic baselines when adequate training is unavailable~\cite{marl24albrecht}; ENSURE's contribution is precisely to make adaptive RL \emph{viable} in this regime by combining warm-start and masking. The separation is primarily driven by two metrics. $S_{\mathrm{cong}}$: ENSURE maintains near-unity congestion satisfaction through coordinated path switching and VNF migration, whereas Static cannot reroute traffic away from failed/congested links. $S_{\mathrm{cons}}$: Static's persistent congestion cascades into transient oscillatory behavior. ENSURE also achieves high $S_{\mathrm{delay}}$ through masking-guided path selection, while NoMask/NoWarm's unwarranted reconfigurations inflate delay. The below-threshold $S_{\mathrm{feas}}$ of NoMask/NoWarm confirms that unconstrained exploration and poor initialization systematically violate placement constraints.
% Fig.~\ref{fig:combined_bars} presents the per-metric satisfaction averaged across all LTI steps and SIs over 20 evaluation runs. ENSURE achieves $S_{\mathrm{E2E}}=0.73$, substantially outperforming Static ($0.53$), NoMask ($0.24$), and NoWarm ($0.22$). Two metrics primarily drive this separation. First, $S_{\mathrm{cong}}$: ENSURE maintains ${\sim}0.97$ through coordinated path switching and VNF migration, whereas Static collapses to ${\sim}0.43$ as the fixed shortest paths traverse the failed links, accumulating congestion that cannot be alleviated without rerouting. Second, $S_{\mathrm{cons}}$: ENSURE (${\sim}0.78$) substantially surpasses Static (${\sim}0.38$) because the latter's degraded $S_{\mathrm{cong}}$ and inability to adapt to link/server failures continuously destabilizes the satisfaction components. Moreover, ENSURE achieves high delay satisfaction ($S_{\mathrm{delay}}>0.8$) through stable, masking-guided TN path selection, whereas NoMask and NoWarm suffer from erratic path changes that inflate E2E delay, yielding $S_{\mathrm{delay}}{\sim}0.5$. The $S_{\mathrm{feas}}$ gap between NoMask/NoWarm (${\sim}0.35$) and ENSURE (${\sim}1.0$) shows that unconstrained exploration (NoMask) and poor initialization (NoWarm) systematically produce infeasible VNF placements that violate server-type or tier constraints.

Fig.~\ref{fig:combined_time} reveals the temporal dynamics. Static initially matches ENSURE in $S_{\mathrm{delay}}$ but collapses when TN failures activate (step~2) and remains degraded until recovery. ENSURE dips mid-episode from the one-LTI coordination lag but converges to ${\sim}0.8$, demonstrating cross-timescale recovery; NoMask/NoWarm stay persistently low due to unnecessary path switches (cf. Sec.~\ref{subsec:stress}). The $S_{\mathrm{cong}}$ subplot is most revealing: Static, by design, cannot reroute within an episode, whereas ENSURE sustains $S_{\mathrm{cong}}>0.97$ throughout via adaptive rerouting and VNF migration. $S_{\mathrm{status}}$ shows ENSURE near-unity while Static drops during the failure window, and the composite $S_{\mathrm{E2E}}$ confirms ENSURE's dominance after the initial LTI-boundary aggregation transient.
% Static drops immediately due to background load on fixed paths, while ENSURE sustains $S_{\mathrm{cong}}>0.97$ throughout via adaptive rerouting and VNF migration. 
% Fig.~\ref{fig:combined_time} shows the temporal evolution of E2E-level satisfaction metrics, aggregated across SIs via the geometric mean. In the $S_{\mathrm{delay}}$ subplot, Static initially matches ENSURE (${\sim}0.93$) but drops sharply to ${\sim}0.5$ at step~2 when the TN link failures activate and remains degraded until recovery around step~7. ENSURE's delay satisfaction dips to ${\sim}0.68$ mid-episode (as the E2E agent absorbs the MTI-level failure response with a one-LTI lag) but stabilizes at ${\sim}0.8$, confirming effective cross-timescale coordination. By contrast, NoMask and NoWarm exhibit persistently low $S_{\mathrm{delay}}$ (${\sim}$0.6) due to unnecessary and counterproductive TN path switches, consistent with the instability observed in the isolated TN stress test (Sec.~\ref{subsec:stress}). The $S_{\mathrm{cong}}$ subplot exposes the most severe Static degradation: it drops to ${\sim}0.37$ at step~1 and remains flat thereafter, as the static TN paths cannot divert traffic from congested links (due to extra background load). ENSURE sustains $S_{\mathrm{cong}}>0.97$ throughout, owing to its ability to adaptively migrate VNFs and reroute TN paths. The $S_{\mathrm{status}}$ subplot shows that ENSURE maintains ${\sim}0.98$, whereas Static drops to ${\sim}0.55$ during the concurrent failure window (steps~2--6), then recovers once failures are lifted. NoMask and NoWarm follow a similar mid-episode dip in $S_{\mathrm{status}}$ (${\sim}0.55$--$0.65$) as their unstable path reconfigurations did not coincide with the failure event. The $S_{\mathrm{E2E}}$ trajectory confirms ENSURE's dominance: it stabilizes at ${\sim}0.8$ after the initial LTI transient, while Static oscillates between $0.45$ and $0.65$, and NoMask/NoWarm remain below $0.35$. Note that the E2E agent aggregates metrics from domain agents at the \emph{beginning} of each LTI, which introduces a one-step initialization transient visible in the first LTI of all curves.

Fig.~\ref{fig:combined_delay} exposes the most operationally critical distinction. Static's eMBB delay far exceeds the 20~ms budget because failed TN links impose a 200~ms penalty on the fixed shortest paths; ENSURE reroutes around failures and stays within budget, while NoMask/NoWarm achieve intermediate but highly variable delays as their policies only occasionally discover functional paths. For URLLC, ENSURE and Static both meet the 2~ms budget, but NoMask/NoWarm exhibit elevated variance for SI3 due to path changes the unmasked policies have not learned to suppress within the 1000-episode budget. These changes are feasible but sub-optimal: they yield no satisfaction gain and incur extra control-plane actions.
%but NoMask/NoWarm exhibit elevated variance for SI3 due to erratic path changes.
%Fig.~\ref{fig:combined_delay} reveals the most operationally critical finding. For eMBB SIs, the Static method incurs mean delays of ${\sim}85$~ms (far exceeding the 20~ms budget) because the failed TN links impose a 200~ms penalty delay on the static shortest paths in affected steps. ENSURE successfully reroutes eMBB SIs to alternative paths that bypass the failed links, maintaining mean delay below ${\sim}20$~ms, within the eMBB budget. NoMask and NoWarm produce intermediate delays (${\sim}48$--$65$~ms) with high variance, as their policies occasionally discover functional paths but lack the consistency provided by masking. For URLLC SIs, the inset zoom shows that typical delays for ENSURE and Static methods remain well below the 2~ms budget ($<0.15$~ms); however, NoMask and NoWarm schemes exhibit elevated mean delay with large standard deviations for SI3, driven by erratic and unnecessary path changes. ENSURE exhibits the lowest mean delay among the four methods for all SIs, with the tightest variance, further illustrating that satisfaction-triggered masking and warm-start jointly provide the most reliable delay performance under severe multi-domain stress.

\begin{figure}[t]
\centering
\includegraphics[width=0.86\columnwidth]{figures/01_eval_satisfaction_bars-All_Stress.pdf}
\caption{Per-metric satisfaction comparison under \texttt{All\_Stress}.}% (averaged over 20 runs). Error bars indicate $\pm$1 standard deviation.}
\label{fig:combined_bars}
\end{figure}

\begin{figure}[t]
\centering
\includegraphics[width=1.01\columnwidth]{figures/02_eval_satisfaction_time-All_Stress.pdf}
\caption{E2E satisfaction metrics over LTI steps under \texttt{All\_Stress}: $S_{\mathrm{delay}}$, $S_{\mathrm{cong}}$, $S_{\mathrm{status}}$, $S_{\mathrm{feas}}$, $S_{\mathrm{cons}}$, and $S_{\mathrm{E2E}}$ (geom. mean across SIs).}
\label{fig:combined_time}
\end{figure}

% Fig.~\ref{fig:combined_time} reveals distinctive temporal dynamics not observed in \texttt{CN\_Stress}. In the $S_{\mathrm{delay}}$ subplot, Static initially matches ENSURE (${\sim}0.93$) but drops sharply to ${\sim}0.5$ at step~2 when the TN link failures activate, and remains degraded until recovery around step~7; this contrasts with \texttt{CN\_Stress}, where Static maintained high $S_{\mathrm{delay}}$ throughout. ENSURE's delay satisfaction dips to ${\sim}0.68$ mid-episode (as the E2E agent absorbs the MTI-level failure response with a one-LTI lag) but stabilizes at ${\sim}0.8$, confirming effective cross-timescale coordination. The $S_{\mathrm{cong}}$ subplot exposes the most severe Static degradation: it drops to ${\sim}0.37$ at step~2 and remains flat thereafter, as the static TN paths cannot divert traffic from congested failed links. ENSURE sustains $S_{\mathrm{cong}}>0.97$ throughout. The $S_{\mathrm{status}}$ subplot shows that ENSURE maintains ${\sim}0.98$, whereas Static drops to ${\sim}0.55$ during the concurrent failure window (steps~2--6), then recovers once failures are lifted, a pattern absent in \texttt{CN\_Stress} where no TN failures occurred. NoMask and NoWarm follow a similar mid-episode dip in $S_{\mathrm{status}}$ (${\sim}0.55$--$0.65$) as their unstable path reconfigurations did not coincide with the failure event. The $S_{\mathrm{E2E}}$ trajectory confirms ENSURE's dominance: it stabilizes at ${\sim}0.8$ after the initial LTI transient, while Static oscillates between $0.45$ and $0.65$, and NoMask/NoWarm remain below $0.35$.

% Fig.~\ref{fig:combined_delay} reveals the most operationally critical finding. For eMBB SIs, the Static method incurs mean delays of ${\sim}85$~ms (far exceeding the 20~ms budget) because the failed TN links impose a 200~ms penalty delay on the static shortest paths in affected steps. ENSURE successfully reroutes eMBB SIs to alternative paths that bypass the failed links, maintaining mean delay below ${\sim}20$~ms, within the eMBB budget. NoMask and NoWarm produce intermediate delays (${\sim}48$--$65$~ms) with high variance, as their policies occasionally discover functional paths but lack the consistency provided by masking. For URLLC SIs, the inset zoom shows that typical delays for all methods remain well below the 2~ms budget ($<0.15$~ms); however, SI3 exhibits elevated mean delay with large standard deviations across all methods, driven by outlier episodes where concurrent TN and CN failures momentarily disrupt the E2E path. ENSURE exhibits the lowest mean delay among the four methods for all SIs, with the tightest variance, further illustrating that satisfaction-triggered masking and warm-start jointly provide the most reliable delay performance under severe multi-domain stress.
\begin{figure}[t]
\centering
\includegraphics[width=0.86\columnwidth]{figures/03c_eval_delay_total_per_si_grouped_split-All_Stress.pdf}
\caption{Per-SI E2E delay under \texttt{All\_Stress}, grouped by slice type.} % The Static method's eMBB delay spike results from the 200~ms penalty on failed TN links.}
\label{fig:combined_delay}
\end{figure}

%%% COMMENTED: Scalability Analysis (to be added once results are available) %%%
% \subsubsection{Scalability Analysis}\label{subsubsec:scalability}
% To assess ENSURE's computational overhead as the problem size grows, we vary the number of CUs
% while keeping the episode length ($T=100$), user count (40), and number of training episodes
% constant. With 2~slice types, $I=2 \cdot |\mathcal{C}|$ SIs, so doubling CUs doubles the SI count
% and accordingly increases TN/CN/E2E action and observation dimensions.
%
% Table~\ref{tab:scalability} reports the average wall-clock time per episode during training and
% evaluation (inference), as well as the normalized E2E reward at each curriculum stage and the two
% evaluation scenarios. The first row (2~CUs, 4~SIs) corresponds to the configuration used throughout
% this paper; subsequent rows will be populated once simulations complete.
%
% \begin{table*}[t]
% \centering
% \caption{Scalability analysis: wall-clock time and E2E reward vs.\ number of CUs (SIs). All runs use $T=100$, 40~users, and 1000~training episodes. Empty cells are pending simulation results.}
% \label{tab:scalability}
% \scriptsize
% \setlength{\tabcolsep}{2.5pt}
% \begin{tabular}{@{}c c | c c | c c c | c c @{}}
% \toprule
% & & \multicolumn{2}{c|}{\textbf{Avg. Wall-clock / Episode (s)}} & \multicolumn{3}{c|}{\textbf{Norm. E2E Reward (Training)}} & \multicolumn{2}{c}{\textbf{Norm. E2E Reward (Eval)}} \\
% \cmidrule(lr){3-4} \cmidrule(lr){5-7} \cmidrule(lr){8-9}
% \textbf{CUs} & \textbf{SIs} & \textbf{Training} & \textbf{Inference} & \textbf{Initial} & \textbf{Intermediate} & \textbf{Final} & \texttt{CN\_Stress} & \texttt{All\_Stress} \\
% \midrule
% 2 & 4  & \textcolor{red}{XX} & \textcolor{red}{XX} & \textcolor{red}{XX} & \textcolor{red}{XX} & \textcolor{red}{XX} & \textcolor{red}{XX} & \textcolor{red}{XX} \\
% 4 & 8  & --- & --- & --- & --- & --- & --- & --- \\
% 6 & 12 & --- & --- & --- & --- & --- & --- & --- \\
% 8 & 16 & --- & --- & --- & --- & --- & --- & --- \\
% \bottomrule
% \end{tabular}
% \end{table*}
%
% The key scalability concern is the E2E agent's action space, which grows as $\mathcal{O}(Z^I)$ (SI-to-site
% assignments). However, the toggle-based action design (Sec.~\ref{solution-section}) reduces this to
% $\mathcal{O}(Z \cdot I)$ binary decisions, and satisfaction-triggered masking further constrains the
% effective action space at each step. We expect sub-linear wall-clock growth due to these mechanisms.
%%% END COMMENTED SCALABILITY %%%
%=================================

\subsection{Summary of Findings}
\label{subsec:findings}
The evaluation yields three findings. \textbf{(1) Components compound:} removing masking or warm-start each collapses $S_{\mathrm{E2E}}$ by ${\sim}$70\%. \textbf{(2) Multi-domain adaptation is necessary:} against the RAN-centric Static baseline~\cite{jiot24filali,camad24ebrahimi}, ENSURE holds eMBB delay within SLA whereas Static exceeds it by $>4\times$. \textbf{(3) Slice-aware resilience:} ENSURE keeps both eMBB and URLLC within SLA under multi-domain stress; all baselines violate at least one.
% The evaluation yields three principal findings:
% \textbf{(1)~Each ENSURE component contributes measurably.} Removing either action masking or ILP warm-start causes $S_{\mathrm{E2E}}$ to collapse by ${\sim}$70\%, confirming that both are essential and their benefits compound.
% \textbf{(2)~Multi-domain adaptation is necessary.} Against the RAN-centric Static baseline (a proxy for fixed-infrastructure approaches~\cite{camad24ebrahimi,jiot24filali}), ENSURE preserves eMBB delay within SLA, whereas Static exceeds it by over $4\times$ when TN links fail.
% \textbf{(3)~Slice-aware resilience.} ENSURE maintains both eMBB and URLLC delays within SLA under multi-domain stress, whereas all evaluated baselines violate at least one. The TN agent's rapid failure response propagates to E2E recovery through hierarchical coordination.
% The evaluation yields three principal findings: 
% \textbf{(1)~Each ENSURE component contributes measurably.} Removing either action masking or ILP warm-start causes $S_{\mathrm{E2E}}$ to collapse by ${\sim}$70\%, confirming that both are essential and their benefits compound. 
% \textbf{(2)~Multi-domain adaptation is necessary.} Static performs well under benign conditions but degrades severely under stress, with eMBB delay exceeding its SLA by over $4{\times}$ due to its inability to reroute around failed links. 
% \textbf{(3)~ENSURE provides slice-aware resilience.} ENSURE maintains both eMBB and URLLC delays within SLA budgets under multi-domain stress, whereas all baselines violate at least one. The TN agent's rapid failure response propagates to E2E-level recovery via hierarchical coordination. 
%Consistent conclusions were obtained under additional single-domain scenarios (omitted for space).

% The evaluation yields three principal findings: 
% \textbf{(1)~Each ENSURE component contributes measurably.} Action masking improves $S_{\mathrm{E2E}}$ from 0.24 (NoMask) to 0.73 under \texttt{All\_Stress}, while ILP warm-start raises it from 0.22 (NoWarm) to 0.73, confirming that both mechanisms are essential and their benefits compound at the E2E level.
% \textbf{(2)~Multi-domain adaptation is necessary.} Static achieves competitive rewards under benign conditions (Table~\ref{tab:e2e_curriculum}, Initial stage) but degrades to $S_{\mathrm{E2E}}=0.5$ under \texttt{All\_Stress}, a 46\% deficit relative to ENSURE, with eMBB delay reaching ${\sim}85$~ms (4.25$\times$ the SLA) due to failed TN links that the static policy cannot circumvent.
% \textbf{(3)~ENSURE provides slice-aware resilience.} Under the multi-domain stress scenario, ENSURE maintains delay within eMBB/URLLC SLA budgets, whereas NoMask/NoWarm produce unreliable URLLC delay and Static incurs penalty delays. The TN agent's two-step failure response (Sec.~\ref{subsec:stress}) propagates to E2E-level recovery via multi-timescale coordination. Consistent conclusions were obtained under additional single-domain stress scenarios (not reported here for space), further validating ENSURE's robustness.

%%%%%%%%%%
\section{Conclusions and Future Directions} \label{conclusion-section}
This paper introduced ENSURE, a hierarchical cooperative MA-DRL framework for E2E slice resource management in Open RAN-compatible 5G/6G networks. Four agents (SAC RAN; PPO TN, CN, E2E) operate at three timescales (NRTI/MTI/LTI) and coordinate through an SI abstraction that decouples per-user radio control from per-slice TN/CN decisions. Satisfaction-triggered action masking restricts reconfiguration to entities below the requirement pivot, and ILP warm-starting seeds each episode with feasible assignments. In simulation, under a three-stage curriculum and severe multi-domain stress, ENSURE reaches $S_{\mathrm{E2E}}=0.73$, outperforming a RAN-centric Static baseline and two internal ablations by at least 38\%, while keeping both eMBB and URLLC delays within SLA. The advantage widens with scenario difficulty, indicating that timescale separation, satisfaction-triggered masking, and ILP warm-starting are jointly responsible for the gap; a common-testbed comparison against published MA-DRL E2E frameworks, using criteria shared with prior work, is expected to narrow this margin while preserving ENSURE's multi-domain resilience advantage.
% This paper demonstrated, using simulation, the benefits of ENSURE, a hierarchical MA-DRL framework for E2E slice resource management in Open RAN-compatible 6G networks. ENSURE combines satisfaction-triggered action masking with ILP-based warm-starting across four cooperative agents at three timescales. The proposed method achieves the highest E2E satisfaction score under multi-domain stress, outperforming all baselines by at least 46\% while maintaining slice-specific delay budgets. %Future directions include scalability to larger topologies, agent fault tolerance, and virtualized DU/CU/RIC orchestration.
% This paper introduced \textit{\underline{EN}d-to-End \underline{S}lice Reso\underline{UR}ce Manag\underline{E}ment} (\textit{ENSURE}), a hierarchical MA-DRL framework for E2E slice resource management in Open RAN-compatible 6G networks. Four cooperative agents (RAN (SAC), TN, CN, and E2E (Masked PPO)) operate at three distinct timescales, coordinated via SI-level information exchange. Satisfaction-triggered action masking suppresses unnecessary reconfigurations, while warm-starting seeds each episode with feasible assignments to accelerate convergence. Simulation results under curriculum-based training and a multi-domain stress scenario show that ENSURE achieves the highest E2E satisfaction, outperforming all baselines by at least 46\%. ENSURE maintains eMBB and URLLC delays within their respective budgets.%,whereas the Static baseline (equivalent to RAN-only approaches) incurs ${\sim}85$~ms eMBB delay when TN links fail. The isolated TN stress test further confirms that masked PPO detects and responds to link failures within two MTI steps, and this rapid domain-level recovery propagates to E2E-level satisfaction through the hierarchical coordination structure.

Future directions include (i)~the common-testbed comparison noted above (cf.\ Sec.~\ref{subsec:ablation}); (ii)~scalability to larger topologies with $\ge30$ SIs; (iii)~xApp/rApp redundancy and recovery for agent fault tolerance; (iv)~virtualized DU/CU/RIC orchestration over disaggregated placement targets. %; (v)~extension beyond eMBB/URLLC to deterministic and joint communication--sensing slices.
% Future directions include scalability analysis to larger topologies with more CUs and SIs, integration of redundancy mechanisms for agent fault tolerance, and extension to virtualized DU/CU/RIC orchestration for comprehensive cloud--edge resource orchestration. 
%Another promising direction consists in extending
%%%%%%%%%%%%%%%%%%%%
%%% OUTDATED CONCLUSIONS WITH SOUFIENE COMMENTS%%%%%%%%%%%%%
% This paper introduced a novel MA-DRL-based multi-timescale framework for efficient E2E resource management in Open RAN-compatible 6G networks. Specifically, we proposed an independent MA-PPO framework to simultaneously optimize resources across the RAN, TN, and CN domains, guided by a higher-level E2E orchestrator agent. By operating on multiple timescales and leveraging PPO’s stability and efficiency, our approach substantially reduces the complexity of the joint optimization problem, ensuring convergence towards optimal policies tailored for each domain.

% The simulation results demonstrated the effectiveness of the proposed methodology, particularly within the RAN domain. ENSURE significantly outperformed a heuristic-based baseline, achieving a remarkable \textcolor{red}{X\%} improvement in Overall Satisfaction Level (OSL). Detailed analysis also highlighted substantial improvements for eMBB and URLLC users, with an increase of rate and delay satisfaction levels by \textcolor{red}{X\%} and \textcolor{red}{X\%}, respectively, clearly demonstrating the benefits of intelligent DRL-based resource allocation.

% Looking ahead, several research challenges must be addressed to fully unlock the potential of this multi-agent approach. A key future research direction is the integration of redundancy strategies within the agent ecosystem to mitigate single point of failure, particularly critical in high-availability scenarios inherent to 6G networks. Another promising direction consists in extending the virtualization paradigm to encompass DUs, CUs, and RICs as fully virtualized entities to enable comprehensive orchestration across cloud, edge, and network resources, further advancing the practicality and scalability of 6G slicing frameworks.

\section*{Acknowledgments}
This work was supported by the Centre for Future Transport and Cities (CFTC), Coventry University, Coventry, UK. %The authors express their gratitude to the anonymous reviewers who contributed to enhancing the quality of the manuscript.

\bibliographystyle{IEEEtran_modified} 
\bibliography{main}{}

% \begin{IEEEbiography}[{\includegraphics[width=1in,height=1.1in,clip,keepaspectratio]{figures/sina.jpg}}]{Sina Ebrahimi}
% received the B.Sc. and M.Sc. degrees in software engineering from the University of Zanjan, Iran, and Tarbiat Modares University, Tehran, Iran, in 2016 and 2020, respectively. He is currently pursuing the Ph.D. degree at Coventry University, UK. His research interests include network slicing, DRL, AI-enabled 5G/6G networking technologies, and resource allocation/optimization. He has served as a reviewer for IEEE TMC, OJ-COMS, and LNET, and as a TPC member/reviewer for IEEE ICC, VTC, WCNC, and INFOCOM.
% \end{IEEEbiography}

% \begin{IEEEbiography}
% [{\includegraphics[width=1in,height=1.1in,clip,keepaspectratio]{figures/hui.png}}]{Hui Zhou} received . 
% \end{IEEEbiography}

% \begin{IEEEbiography}[{\includegraphics[width=1in,height=1.1in,clip,keepaspectratio]{figures/faouzi.png}}]{Faouzi Bouali} received a Ph.D. in Signal Theory and Communications from the University Politecnica de Catalunya (UPC), Spain in 2013. He is an Assistant Professor (Senior Lecturer) with the Centre for Future Transport and Cities, Coventry University, UK. His research interests span the field of communications and networking, with a focus on radio resource management, architectural innovations for next-generation networks (e.g., network slicing and Open RAN), 5G/6G experimental research and testbeds, and connected and automated industries.
% \end{IEEEbiography}

% \begin{IEEEbiography}[{\includegraphics[width=1in,height=1.1in,clip,keepaspectratio]{figures/Djahel.jpg}}]{Soufiene Djahel} (SM'16) holds Ph.D (2010) from USTL, France. He is a Professor in the Centre for Future Transport and Cities at Coventry University. His research interests include the design and evaluation of communication, planning, optimization and security algorithms and solutions for wireless networked systems with a focus on CAVs and UAVs. His research was supported by the Newton Mosharrafa Fund, JSPS, IntelliCentrics UK LTD, EPSRC DTP and the Transport Systems Catapult.
% \end{IEEEbiography}

% \begin{IEEEbiography}[{\includegraphics[width=1in,height=1.1in,clip,keepaspectratio]{figures/olivier.jpg}}]{Olivier C. L. Haas} (Senior Member, IEEE, 2022) received a DUT (Diplôme Universitaire de Technologie) from the Université Joseph Fourier, Grenoble, France, B.Eng., M.Sc., and Ph.D. from Coventry University, Coventry, UK. He is an Associate Professor with 25 years’ experience, 28 PhD completions. He heads the Intelligent Transportation Systems and 5G, Centre for Future Transport and Cities, Coventry University. \href{https://pureportal.coventry.ac.uk/en/persons/olivier-haas}{pureportal.coventry.ac.uk/en/persons/olivier-haas}
% \end{IEEEbiography}

% \clearpage
% \begin{table}[!htbp] % we can remove this if we really need to cut stuff
% \caption{{TN propagation paths from CN sites to CUs.}}
% \label{tab:tn-prop}
% \small
% \centering
% \setlength{\tabcolsep}{3pt} % tighter columns so it fits in two columns
% \renewcommand{\arraystretch}{1.15}
% \newcolumntype{P}[1]{>{\raggedright\arraybackslash}p{#1}}

% \begin{tabular}{P{1.9cm} P{1.42cm} P{1.95cm} P{9.49cm} P{1.5cm}}
% \textbf{\shortstack{CN site\\(core)}} & \textbf{\shortstack{CU\\(edge)}} & \textbf{ \shortstack{Selected\\Path}} & \textbf{Path (excluding src, dst)} & \textbf{\shortstack{$\tau^{\mathrm{prop}}$\\(ms)}} \\
% \hline

% \multirow{3}{*}{DNVR} & \multirow{3}{*}{LA1} & primary   & SNVA $\to$ LA2 & 10.15 \\
%                         &                       & secondary & STTL $\to$ SNVA $\to$ LA2 & 16.11 \\
%                         &                       & tertiary  & KSCY $\to$ HSTN & 19.77 \\ \hline

% \multirow{3}{*}{WASH} & \multirow{3}{*}{LA1} & primary   & ATLA $\to$ HSTN & 20.81 \\
%                         &                       & secondary & ATLA $\to$ IPLS $\to$ KSCY $\to$ DNVR $\to$ SNVA $\to$ LA2 & 25.82 \\
%                         &                       & tertiary  & NYCM $\to$ CHIN $\to$ IPLS $\to$ KSCY $\to$ DNVR $\to$ SNVA $\to$ LA2 & 27.07 \\ \hline

% \multirow{3}{*}{STTL} & \multirow{3}{*}{LA1} & primary   & SNVA $\to$ LA2 & 8.26 \\
%                         &                       & secondary & DNVR $\to$ SNVA $\to$ LA2 & 18.00 \\
%                         &                       & tertiary  & DNVR $\to$ KSCY $\to$ HSTN & 27.63 \\ \hline

% LA1    & LA1 & primary   &  & 0.01 \\  \hline

% \multirow{3}{*}{DNVR} & \multirow{3}{*}{LA2} & primary   & SNVA & 10.06 \\
%                         &                       & secondary & STTL $\to$ SNVA & 16.02 \\
%                         &                       & tertiary  & KSCY $\to$ HSTN $\to$ LA1 & 19.87 \\ \hline

% \multirow{3}{*}{WASH} & \multirow{3}{*}{LA2} & primary   & ATLA $\to$ HSTN $\to$ LA1 & 20.90 \\
%                         &                       & secondary & ATLA $\to$ IPLS $\to$ KSCY $\to$ DNVR $\to$ SNVA & 25.73 \\
%                         &                       & tertiary  & NYCM $\to$ CHIN $\to$ IPLS $\to$ KSCY $\to$ DNVR $\to$ SNVA & 26.98 \\ \hline

% \multirow{3}{*}{STTL} & \multirow{3}{*}{LA2} & primary   & SNVA & 8.17 \\
%                         &                       & secondary & DNVR $\to$ SNVA & 17.91 \\
%                         &                       & tertiary  & DNVR $\to$ KSCY $\to$ HSTN $\to$ LA1 & 27.72 \\ \hline

% \multirow{3}{*}{LA1}    & \multirow{3}{*}{LA2} & primary   &  & 0.09 \\
%                         &                       & secondary & HSTN $\to$ KSCY $\to$ DNVR $\to$ SNVA & 29.83 \\
%                         &                       & tertiary  & HSTN $\to$ KSCY $\to$ DNVR $\to$ STTL $\to$ SNVA & 35.80 \\

% \hline
% \end{tabular}
% \end{table}

\end{document}


